\chapter{Homological algebra and group cohomology}\label{ch:cohomology}

The paper's Tamagawa factor $\Tam_q(\ad T)$ is the Fitting ideal of the Frobenius-invariants of the torsion of $H^1(I_q,\ad T)$, and its proof of Lemma
lem:tamfactor computes that group in one paragraph: inflation-restriction, the identity $H^1(\langle t\rangle,M)=M/(t-1)M$, the conjugation action of Frobenius,
an explicit $3\times3$ matrix, the invariants. Around Lemma lem:dfg21 it uses long exact sequences, Pontryagin duals, the vanishing of the cohomology of the real
place, and the formula $\Hur=M^{I}/(\Frob-1)M^I$. This chapter supplies the group cohomology behind all of it, for abstract and for profinite groups, with
continuous cochains and with the three kinds of coefficients the paper uses (finite or torsion discrete modules, finitely generated $\Zl$-modules, $\Ql$-vector
spaces). Its centre is \Cref{sec:cohomology:model}, where the group-cohomological part of Lemma lem:tamfactor is proved line by line. Galois cohomology proper
(Kummer theory, local duality, Selmer groups) is \Cref{ch:galcoh}; the Bloch-Kato conditions and the Tamagawa factors are \Cref{ch:blochkato}.

\section*{Core results at a glance}

Throughout, $\ell$ is a prime; groups act on the left; all cohomology of profinite groups is continuous cochain cohomology with coefficients of type (D) (discrete,
finite or $\ell$-primary torsion), (C) (finitely generated $\Zl$-modules) or (V) (finite-dimensional $\Ql$-spaces); $\Frob$ is an arithmetic Frobenius.

\begin{enumerate}[label=\textbf{\thechapter.\arabic*}]
\item \textbf{Conventions.} Types (D), (C), (V); $\Frob$ arithmetic, $\Frob\,t\,\Frob^{-1}=t^q$, $\Frob^{-1}t\Frob=t^{q^{-1}}$; conjugation action
$(\sigma\cdot c)(h_1,\dots,h_n)=\sigma c(\sigma^{-1}h_1\sigma,\dots)$; $\ell$ odd is used only in \Cref{sec:cohomology:model} and for groups of order $2$. \emph{(convention)}
\Cref{conv:cohomology:conventions}.
\item \textbf{Complexes.} Cochain complexes, cocycles, coboundaries, $H^n=Z^n/B^n$, morphisms and homotopies. \emph{(definition)} \Cref{def:cohomology:complex}.
\item \textbf{Functoriality.} Morphisms induce maps on $H^n$ functorially; homotopic morphisms induce the same map. \emph{(proved here)} \Cref{lem:cohomology:complex-basics}.
\item \textbf{Snake lemma.} For a map of exact rows $A'\to A\to A''\to0$, $0\to B'\to B\to B''$: $\ker f'\to\ker f\to\ker f''\xrightarrow{\delta}\coker f'\to\coker f\to\coker f''$ exact, $\delta$ by
lift, apply $f$, pull back. \emph{(stated; proved in \Cref{ch:modules})} \Cref{lem:cohomology:snake}.
\item \textbf{Long exact sequence.} A short exact sequence of complexes gives a natural long exact sequence with $\delta[c]=[a]$, $i(a)=db$, $p(b)=c$. \emph{(proved here)}
\Cref{thm:cohomology:les}.
\item \textbf{Resolutions.} Projective and injective modules and resolutions. \emph{(definition)} \Cref{def:cohomology:resolutions}.
\item \textbf{Existence and comparison.} Free and injective resolutions exist; comparison theorem; over a PID resolutions of length one, divisible modules
injective. \emph{(stated; reference)} \Cref{prop:cohomology:resolutions}.
\item \textbf{$\Ext$ and $\Tor$.} $\Ext^n_R(M,N)=H^n\Hom_R(P_\bullet,N)=H^n\Hom_R(M,I^\bullet)$, $\Tor_n^R(M,N)=H_n(P_\bullet\otimes N)$. \emph{(definition)} \Cref{def:cohomology:ext-tor}.
\item \textbf{Properties.} Independence of resolutions, long exact sequences, vanishing; over a PID nothing in degree $\ge2$; $\Ext^1_\Z(\Z/m,N)=N/mN$,
$\Tor_1^\Z(\Z/m,N)=N[m]$, $\Ext^1_{\Zl}(\Z/\ell^a,\Z/\ell^b)=\Z/\ell^{\min(a,b)}$. \emph{(stated; reference; examples computed here)} \Cref{prop:cohomology:ext-tor}.
\item \textbf{$\Ext^1$ and extensions.} $\Ext^1_R(M,N)$ classifies extensions of $M$ by $N$; $0$ is the split one. \emph{(stated; reference)} \Cref{thm:cohomology:ext1}.
\item \textbf{$G$-modules.} $M^G$, $M_G=M/I_GM$, $N_G$; $\Hom_G(M,N)=\Hom(M,N)^G$; $\Zl[G]$-modules give $\Zl$-linear cohomology. \emph{(definition)} \Cref{def:cohomology:gmodule}.
\item \textbf{Inhomogeneous cochains.} $C^n(G,M)=\operatorname{Map}(G^n,M)$ with the standard $d$; $(d^0m)(g)=gm-m$, $(d^1\varphi)(g,h)=g\varphi(h)-\varphi(gh)+\varphi(g)$;
$H^n(G,M)=H^n(C^\bullet(G,M))$. \emph{(definition)} \Cref{def:cohomology:cochains}.
\item \textbf{$d\circ d=0$; functoriality in $M$.} \emph{(proved here in degrees $0$, $1$; degree $2$ an exercise; general degree by reference)} \Cref{lem:cohomology:dd}.
\item \textbf{Cohomology as $\Ext$.} $H^n(G,M)\simeq\Ext^n_{\Z[G]}(\Z,M)$ via the bar resolution; $H^n(G,-)$ are the derived functors of $M\mapsto M^G$. \emph{(stated; reference)}
\Cref{thm:cohomology:as-ext}.
\item \textbf{$H^0$ and $H^1$.} $H^0=M^G$; $H^1$ = crossed homomorphisms modulo principal ones, $\varphi(1)=0$, $\varphi(g^{-1})=-g^{-1}\varphi(g)$; trivial action: $H^1=\Hom(G,M)$;
$0\to M'^G\to M^G\to M''^G\xrightarrow{\delta}H^1(M')\to H^1(M)\to H^1(M'')$ exact with $\delta(m'')=[g\mapsto gm-m]$. \emph{(proved here)} \Cref{prop:cohomology:h0-h1}.
\item \textbf{The cocycle in the corner.} $\rho=\sm{\eta_1}{\kappa}{0}{\eta_2}$ is a homomorphism iff $\kappa\eta_2^{-1}\in Z^1(G,R(\eta_1\eta_2^{-1}))$; base change $e_2\mapsto e_2+re_1$ changes it
by $d^0r$; split iff the class is $0$; every class occurs. For $\sm{\chi_\ell\varepsilon}{\varepsilon\kappa}{0}{\varepsilon}$ the cocycle is $\kappa\in Z^1(G,\Zl(1))$, a homomorphism on $I$.
\emph{(proved here)} \Cref{prop:cohomology:corner}.
\item \textbf{$H^2$ and extensions.} $H^2(G,M)$ classifies extensions $1\to M\to\mathcal E\to G\to1$ inducing the action; cocycle $s(g)s(h)s(gh)^{-1}$. \emph{(proof sketched here;
Brown)} \Cref{thm:cohomology:h2-ext}.
\item \textbf{Long exact sequence in $M$.} $0\to M'\to M\to M''\to0$ gives $0\to H^0(G,M')\to\cdots\to H^n(G,M'')\xrightarrow{\delta}H^{n+1}(G,M')\to\cdots$, natural, $\Zl$-linear.
\emph{(proved here)} \Cref{prop:cohomology:les-groups}.
\item \textbf{$H^0$ and torsion.} $H^0(G,W[m])=H^0(G,W)[m]$; for a $\Zl[G]$-module $T$ without $\ell$-torsion, $W[\ell^k]\simeq T/\ell^kT$ for $W=T\otimes\Qlmod$, so $H^0(G,W)[\ell^k]=H^0(G,T/\ell^kT)$; a finite $\ell$-group with
$\#A[\ell]\le\ell$ is cyclic of length the largest $k$ with an element of order $\ell^k$; if $\#H^0(G,T/\ell T)\le\ell$, $H^0(G,W)$ is finite cyclic or $\simeq\Qlmod$, and when finite its
length is the largest $k$ with an invariant element of $T/\ell^kT$ outside $\ell T/\ell^kT$. \emph{(proved here, given \Cref{ch:modules})} \Cref{lem:cohomology:h0-torsion}.
\item \textbf{Compatible pairs.} $(f\colon G'\to G,\ \varphi\colon M\to M')$ with $\varphi(f(g')m)=g'\varphi(m)$ induces $(f,\varphi)^*\colon H^n(G,M)\to H^n(G',M')$ on cochains. \emph{(proved here)}
\Cref{prop:cohomology:compatible}.
\item \textbf{$\Res$, $\Inf$, $\Cor$.} Restriction and inflation as compatible pairs; corestriction for finite index, $\Cor(m)=\sum g_im$, $(\Cor\varphi)(g)=\sum_ig_{\tau(i)}\varphi(h_i(g))$;
$\Cor\circ\Res=[G:H]$; for $H$ normal $\Res\circ\Cor=\sum_{\sigma\in G/H}\sigma_*$. \emph{(proved here in degrees $0$, $1$; higher degrees stated)} \Cref{prop:cohomology:res-inf-cor}.
\item \textbf{The order kills.} For $G$ finite of order $N$, $N\cdot H^n(G,M)=0$ ($n\ge1$), explicitly $N\varphi=(-1)^nd\psi$, $\psi=\sum_g\varphi(\dots,g)$; so $H^n=0$ if $N$ is invertible on $M$,
$H^n$ finite for $M$ finitely generated, and for $\ell\nmid N$ and $\ell$-primary or $\Zl$-modules $H^n=0$ and $M^G$ is a direct summand. \emph{(proved here)}
\Cref{thm:cohomology:order-kills}.
\item \textbf{Conjugation action.} For $H\trianglelefteq G$, $(\sigma\cdot\varphi)(h_1,\dots)=\sigma\varphi(\sigma^{-1}h_1\sigma,\dots)$ is a $\Zl$-linear left action on $H^n(H,M)$, trivial on $H$
($(h_0\cdot\varphi)-\varphi=d^0\varphi(h_0)$), so an action of $G/H$; $\Res$ lands in the invariants ($\sigma\cdot\Res c-\Res c=d^0c(\sigma)$). Same for profinite groups.
\emph{(proved here for $n\le1$; triviality on $H$ for $n\ge2$ stated)} \Cref{prop:cohomology:conjugation}.
\item \textbf{Inflation-restriction.} $0\to H^1(G/H,M^H)\to H^1(G,M)\to H^1(H,M)^{G/H}$ exact, extended by the transgression to $H^2(G/H,M^H)\to H^2(G,M)$; $\Inf$ is an isomorphism on
$H^1$ if $H^1(H,M)=0$; valid with continuous cochains, the five-term sequence for types (D) and, when $G/H\simeq\widehat\Z$, (C) and (V)
(\Cref{prop:cohomology:inf-res-zhat}). \emph{(three-term sequence and corollary proved here; five-term sequence and higher version stated, transgression
sketched)} \Cref{thm:cohomology:inf-res}.
\item \textbf{Hochschild-Serre.} $E_2^{p,q}=H^p(G/H,H^q(H,M))\Rightarrow H^{p+q}(G,M)$; if $H^q(H,M)=0$ for $q\ge1$, $\Inf$ is an isomorphism in all degrees. \emph{(stated;
reference)} \Cref{thm:cohomology:hochschild-serre}.
\item \textbf{Finite cyclic groups.} $H^1(\langle t\rangle,M)=\ker N/(t-1)M$ via $\varphi\mapsto\varphi(t)$; $H^2=M^G/NM$; period $2$. \emph{($H^1$ proved here; the rest stated)}
\Cref{thm:cohomology:cyclic}.
\item \textbf{Tate cohomology.} $\hat H^0=M^G/N_GM$, $\hat H^{-1}=\ker N_G/I_GM$, $\hat H^n=H^n$ ($n\ge1$). \emph{(definition)} \Cref{def:cohomology:tate}.
\item \textbf{Properties of $\hat H$.} Long exact sequences in all degrees; killed by $\#G$; induced modules acyclic; period $2$ for cyclic $G$. \emph{(stated; reference)}
\Cref{prop:cohomology:tate}.
\item \textbf{Groups of order $2$.} $H^n(\Z/2,M)$ ($n\ge1$) is killed by $2$; for $\ell$ odd and $\ell$-adic $M$ it is $0$ and $M=M^+\oplus M^-$; so $H^1(\R,M)$ has exponent $2$ and
vanishes for $\ell$-adic $M$, $\ell$ odd. \emph{(proved here)} \Cref{cor:cohomology:z2}.
\item \textbf{Discrete modules.} Continuous action on a discrete module. \emph{(definition)} \Cref{def:cohomology:discrete}.
\item \textbf{Discrete equivalences.} Continuous iff stabilisers open iff $M=\bigcup_UM^U$. \emph{(proved here)} \Cref{lem:cohomology:discrete}.
\item \textbf{Continuous cohomology.} $C^n_{\mathrm{cont}}(G,M)$; types (D), (C), (V); the convention that all profinite cohomology is continuous. \emph{(definition)}
\Cref{def:cohomology:continuous}.
\item \textbf{Basic properties.} $H^1$ = continuous crossed homomorphisms mod principal; discrete $M$: $H^n(G,M)=\varinjlim_UH^n(G/U,M^U)$; type (C) continuity criterion; type (V)
has a stable lattice. \emph{(proved here)} \Cref{prop:cohomology:continuous}.
\item \textbf{Long exact sequences, continuous.} Exist when $p$ has a continuous section, in particular for $M''$ discrete, surjections of finitely generated $\Zl$-modules,
$T\to V\to W$, $T\xrightarrow{\ell^k}T\to T/\ell^k$; and $\ker(H^1(G,T)\to H^1(G,V))=H^1(G,T)_{\mathrm{tor}}=\delta(H^0(G,W))$. \emph{(proved here)} \Cref{prop:cohomology:les-cont}.
\item \textbf{Tate's limit theorem.} $0\to\varprojlim^1H^{n-1}(G,T_k)\to H^n(G,T)\to\varprojlim H^n(G,T_k)\to0$; $H^1(G,T)=\varprojlim H^1(G,T/\ell^kT)$ always; $H^n(G,V)=H^n(G,T)\otimes\Ql$.
\emph{(stated; reference; the case $n=1$ and the statement for $V$ proved here)} \Cref{thm:cohomology:tate-limits}.
\item \textbf{Finiteness of $H^1$.} $H^1(G,T)=\varprojlim H^1(G,T/\ell^k)$; if $H^1(G,A)$ is finite for finite $\ell$-primary $A$ (for instance $G$ finite, procyclic onto $\Zl$, or $I_q$ with $q\ne\ell$), $H^1(G,T)$ is finitely generated and its
torsion finite. \emph{(proved here, given topological Nakayama, \Cref{rem:cohomology:top-nakayama})} \Cref{prop:cohomology:h1-fg}.
\item \textbf{Supernatural orders.} $\#G$, pro-$\ell$ groups, order prime to $\ell$, $G^{(\ell)}=G/K_\ell$. \emph{(definition)} \Cref{def:cohomology:supernatural}.
\item \textbf{Pro-$\ell$ quotients.} $G^{(\ell)}$ is the maximal pro-$\ell$ quotient; $\ell\nmid\#G\Rightarrow G^{(\ell)}=1$, not conversely ($A_4$, $\ell=2$); $K_\ell^{(\ell)}=1$; closed subgroups of
pro-$\ell$ groups and finitely generated $\Zl$-modules are pro-$\ell$; Sylow subgroups. \emph{(proved here; Sylow stated)} \Cref{lem:cohomology:pro-l}.
\item \textbf{Prime-to-$\ell$ vanishing.} If $\ell\nmid\#G$: $H^n(G,M)=0$ ($n\ge1$) for $M$ of type (D) $\ell$-primary, (C), (V), and $M^G$ is a direct summand; if $G^{(\ell)}=1$ and the action is
trivial, $\Hom_{\mathrm{cont}}(G,M)=0$. \emph{(proved here for $H^1$ and for type (D) in all degrees; type (C), (V) in degrees $\ge2$ given \Cref{thm:cohomology:tate-limits})}
\Cref{thm:cohomology:prime-to-l}.
\item \textbf{Through the pro-$\ell$ quotient.} $K$ normal acting trivially: $\Inf\colon H^1(G/K,M)\simeq H^1(G,M)$ if $K^{(\ell)}=1$; in all degrees if $\ell\nmid\#K$; pro-$\ell$ image of $G$ gives
$H^1(G,M)=H^1(G^{(\ell)},M)$. \emph{(proved here, except the all-degrees statement, given Hochschild-Serre)} \Cref{cor:cohomology:pro-l-quotient}.
\item \textbf{Inertia at $q\ne\ell$.} With $K=\ker t_\ell$ ($\ell\nmid\#K$) and $M$ of type (C) or (V) with pro-$\ell$ inertia image: $H^1(I_q,M)\simeq H^1(I_q/K,M)\simeq M/(t-1)M$ via $c\mapsto c(t)$;
$H^n(I_q,M)=0$ for $n\ge2$; for every $M$ of type (C), $H^1(I_q,M)\simeq M^K/(t-1)M^K$ is finitely generated with finite torsion. \emph{(proved here, given
\Cref{ch:local}; the vanishing for $n\ge2$ given stated results)} \Cref{prop:cohomology:tame-inertia}.
\item \textbf{Procyclic groups.} Quotients of $\widehat\Z$; with a surjection $\pi$ onto $\Zl$, $\pi(F)$ is a unit; continuous actions of $\widehat\Z$, $\Zl$ on type (C) modules; unipotent
elements generate $\Zl$-actions. \emph{(proved here)} \Cref{lem:cohomology:procyclic}.
\item \textbf{Cocycles on powers.} $c(F^k)=(1+F+\dots+F^{k-1})c(F)$, and $c(\gamma)\equiv\pi(\gamma)c(F)$ modulo $(F-1)M$; in particular $c(t^a)\equiv a\,c(t)$, $a\in\Zl$. \emph{(proved here)}
\Cref{lem:cohomology:powers}.
\item \textbf{Procyclic, discrete coefficients (Serre).} $\Gamma=\widehat\Z$, $M$ discrete torsion, or $\Gamma=\Zl$, $M$ discrete $\ell$-primary: $H^1(\Gamma,M)\simeq M/(F-1)M$ via $c\mapsto c(F)$,
$\#H^1=\#H^0$ for finite $M$, $H^n=0$ for $n\ge2$. \emph{($H^0$, $H^1$ proved here; $n\ge2$ stated)} \Cref{thm:cohomology:procyclic-discrete}.
\item \textbf{Geometric sums.} For $u\in\Aut(M)$, $M$ finitely generated over $\Zl$: $(1+u+\dots+u^{j-1})M\subseteq\ell^NM$ whenever $a_0\ell^N\mid j$, $a_0$ the order of $u\bmod\ell$. \emph{(proved here)}
\Cref{lem:cohomology:geometric-sums}.
\item \textbf{Procyclic, compact coefficients.} $\Gamma$ procyclic with a surjection onto $\Zl$ (e.g. $\Zl$, $\Zl(1)$, $\widehat\Z$, $I_q/P_q$), $M$ of type (C) or (V): $H^0=\ker(F-1)$,
$H^1(\Gamma,M)\simeq M/(F-1)M$ via $c\mapsto c(F)$, $H^n=0$ for $n\ge2$; $H^1(\Gamma,V)=0$ iff $\det(1-F\mid V)\ne0$. \emph{((i), (ii) proved here; (iii) stated)}
\Cref{thm:cohomology:procyclic-compact}.
\item \textbf{Inflation-restriction over $\widehat\Z$.} $G$ profinite, $H\trianglelefteq G$ closed, $G/H\simeq\widehat\Z$, $M$ of type (D) (torsion), (C) or (V): $0\to H^1(G/H,M^H)\to H^1(G,M)\to
H^1(H,M)^{G/H}\to0$ is exact; with $H^2(G/H,M^H)=0$ this is the five-term sequence of \Cref{thm:cohomology:inf-res}(b) for these coefficients. \emph{(proved here; the vanishing of
$H^2$ stated)} \Cref{prop:cohomology:inf-res-zhat}.
\item \textbf{The model setting.} $G\trianglerighteq I\supseteq K$, $\ell\nmid\#K$, $I/K\simeq\Zl$ generated by $t$, $\Frob$ generating $G/I\simeq\widehat\Z$, $\Frob\,t\,\Frob^{-1}\equiv t^q$, $M$ of type (C),
$K$ acting trivially; dictionary with $G_{\Q_q}$, $I_q$, $\ad T_\ell E$. \emph{(construction; imports from \Cref{ch:local,ch:galrep})} \Cref{constr:cohomology:model}.
\item \textbf{Frobenius on $H^1$ of inertia.} $G$ acting on $M$ by $\rho_M$: $H^1(I,M)\simeq M/(t-1)M$; $\overline\Frob$ acts by $q^{-1}\rho_M(\Frob)$ (for $M=\ad T$:
$q^{-1}\Ad\rho(\Frob)$), well defined since $\rho_M(\Frob)(t-1)M\subseteq(t-1)M$, independent of the lift;
$(H^1(I,M)_{\mathrm{tor}})^{G/I}=\{\bar m\in(M/(t-1)M)_{\mathrm{tor}}:\rho_M(\Frob)\bar m=q\bar m\}$. \emph{(proved here)} \Cref{prop:cohomology:frob-h1}.
\item \textbf{The Tate shape.} For $\rho(t)=\sm1{n'}01$, $\rho(\Frob)=\sm{\varepsilon q}{c'}0\varepsilon$ on $T$, and $\Ad\rho$ on $\ad T$ (basis $e,h,f$): $(t-1)e=0$, $(t-1)h=-2n'e$, $(t-1)f=n'h-n'^2e$; $\Frob\cdot e=qe$,
$\Frob\cdot h=h-2\varepsilon c'e$, $\Frob\cdot f=q^{-1}(f+\varepsilon c'h-c'^2e)$. \emph{(proved here)} \Cref{lem:cohomology:tate-shape}.
\item \textbf{The model computation.} $\ell$ odd, $w=v_\ell(n')$ (the paper's $v$ is $v_\ell(q-1)$): $(t-1)\ad T=n'(\Zl e\oplus\Zl h)$; $H^1(I,\ad T)=(\Zl/n')^2\oplus\Zl$, torsion $(\Z/\ell^w)^2$; $\overline\Frob$: $\bar e\mapsto\bar e$,
$\bar h\mapsto q^{-1}(\bar h-2\varepsilon c'\bar e)$; invariants $\simeq\Z/\ell^w\oplus\Z/\ell^m$, $m=\min(w,v_\ell(q-1),v_\ell(c'))$, length $w+m$; $m$ depends on $c'$ only modulo $(n',q-1)$, the
indeterminacy of $c'$; for $T$: invariants $\Z/\ell^w$ (split) or $0$ (nonsplit). \emph{(proved here)} \Cref{prop:cohomology:model}.
\item \textbf{Fitting ideal of the invariants.} $\Fitt_{\Zl}((H^1(I,\ad T)_{\mathrm{tor}})^{G/I})=\ell^{w+m}\Zl$, the Tamagawa factor; $v_\ell(\Tam_q(\ad T))=v_\ell(n')+\min(v_\ell(n'),v_\ell(q-1),v_\ell(c'))$.
\emph{(proved here, given \Cref{ch:modules}; identifications of $n'$, $c'$ in \Cref{ch:galrep,ch:galcoh})} \Cref{cor:cohomology:fitting}.
\item \textbf{Cup products.} $(a\cup b)(g_1,\dots,g_{p+q})=a(g_1,\dots,g_p)\otimes g_1\cdots g_pb(\dots)$. \emph{(definition)} \Cref{def:cohomology:cup}.
\item \textbf{Properties of $\cup$.} Leibniz rule; well defined on classes; associative, graded commutative, compatible with $\Res$, $\Inf$, $\delta$, $\Cor$. \emph{(degree $(1,1)$ proved
here; the rest stated)} \Cref{prop:cohomology:cup}.
\item \textbf{Pontryagin duality.} $X^*$ for finite, compact and discrete $\ell$-primary modules: exact, reflexive, length-preserving; $(X^*)^G=(X_G)^*$; $H^n(G,A)^*=H_n(G,A^*)$.
\emph{(recalled from \Cref{ch:modules}; duality of cohomology stated)} \Cref{prop:cohomology:pontryagin}.
\item \textbf{Lengths.} $\len A^*=\len A$; in $0\to A\to B\to C\to D\to0$, $\len B=\len A+\len C-\len D$; hence the length identity of Lemma lem:dfg21. \emph{(proved here, given
\Cref{ch:modules})} \Cref{lem:cohomology:lengths}.
\item \sloppy\textbf{Galois cohomology.} $H^n(K,M)=H^n(G_K,M)$; $H^n(\R,M)$; $\Hur(\Q_p,M)=H^1(G_{\Q_p}/I_p,M^{I_p})$, which is the kernel of $H^1(\Q_p,M)\to H^1(I_p,M)$. \emph{(convention; the equality proved here)}
\Cref{conv:cohomology:galois}.
\item \textbf{Unramified cohomology.} $\Hur(\Q_p,M)=M^{I_p}/(\Frob_p-1)M^{I_p}$ for $M$ of type (C) or (V); for $V$ it vanishes iff $\det(1-\Frob_p\mid V^{I_p})\ne0$. \emph{(proved here)}
\Cref{prop:cohomology:unramified}.
\end{enumerate}

\section*{Where the paper uses this chapter}

\begin{itemize}
\item \textbf{Lemma lem:tamfactor} (the Tamagawa factor at a multiplicative prime), every group-cohomological step: the Tate shape as a cocycle
(\Cref{prop:cohomology:corner}); ``by inflation-restriction $H^1(I,M')=H^1(\langle t\rangle,M')=M'/(t-1)M'$'' (\Cref{thm:cohomology:inf-res}, \Cref{cor:cohomology:pro-l-quotient},
\Cref{prop:cohomology:tame-inertia}, \Cref{thm:cohomology:procyclic-compact}); the conjugation action of $G/I$ (\Cref{prop:cohomology:conjugation}); ``$c(t^k)\equiv k\,c(t)$''
(\Cref{lem:cohomology:powers}); ``$\Frob$ acts on $M/(t-1)M$ as $q^{-1}\rho(\Frob)$'', where $q^{-1}\Ad\rho(\Frob)$ is meant (\Cref{prop:cohomology:frob-h1}, \Cref{pl:cohomology:tamfactor}); the explicit $(t-1)M$, the invariants and their
Fitting ideal (\Cref{lem:cohomology:tate-shape}, \Cref{prop:cohomology:model}, \Cref{cor:cohomology:fitting}); $H^1(G/I,\Zl(1)^I)=\Zl(1)/(q-1)$
(\Cref{ex:cohomology:procyclic-examples}). The paperlink \Cref{pl:cohomology:tamfactor} matches each sentence of the proof.
\item \textbf{Section 1.2} (the definition $\Tam_p(M)=\Fitt_{\Zl}((H^1(I_p,M)_{\mathrm{tor}})^{G_{\Q_p}})$) and \textbf{Lemma lem:fpr}: the object $(H^1(I_p,M)_{\mathrm{tor}})^{G_{\Q_p}}$ is
defined by \Cref{prop:cohomology:conjugation} and computed by \Cref{prop:cohomology:frob-h1}; the definition itself is \Cref{ch:blochkato}'s.
\item \textbf{Lemma lem:tamadd}: inflation-restriction with $\Hom(K,M)$ (\Cref{thm:cohomology:inf-res}), ``killed by $\#\Delta$'' (\Cref{thm:cohomology:order-kills}), restriction to a
Sylow subgroup (\Cref{prop:cohomology:res-inf-cor}(d), \Cref{ex:cohomology:sylow}), and $H^1$ of a free $\Z_3[C_3]$-module (\Cref{ex:cohomology:c3}).
\item \textbf{Lemma lem:dfg21 and the paragraph after it}: long exact sequences and connecting maps (\Cref{prop:cohomology:les-groups,prop:cohomology:les-cont}); Pontryagin
duals and lengths (\Cref{prop:cohomology:pontryagin}, \Cref{lem:cohomology:lengths}); ``$H^1(\R,T^D)$ is killed by $2$'' (\Cref{cor:cohomology:z2}, \Cref{pl:cohomology:real-place});
$\Hur(\Q_p,V^D)=(V^D)^{I_p}/(\Frob_p-1)$ and its vanishing (\Cref{prop:cohomology:unramified}); ``$H^1_f(\Q_p,T^D)$ is the torsion of $H^1(\Q_p,T^D)$''
(\Cref{prop:cohomology:les-cont}(iii)).
\item \textbf{Lemma lem:h0}: the $\ell^k$-torsion of $H^0(\Q,W^D)$ and the length of a cyclic $H^0$ (\Cref{lem:cohomology:h0-torsion}, \Cref{pl:cohomology:h0}).
\item \textbf{Proof of Theorem thm:main} and \textbf{Remark rem:ns}: the length bookkeeping along the sequence of Lemma lem:dfg21 (\Cref{lem:cohomology:lengths},
\Cref{pl:cohomology:lengths}).
\end{itemize}

\section{Conventions and plan of the chapter}\label{sec:cohomology:conventions}

This section fixes the conventions that this chapter and \Cref{ch:galcoh,ch:blochkato} share. The paper uses group cohomology in one
dense paragraph, the proof of Lemma lem:tamfactor, and in a handful of sentences around Lemma lem:dfg21, Lemma lem:h0 and the
definition of $\Tam_p$ in its Section 1.2. Every one of those sentences rests on a convention (which Frobenius, which conjugation action,
which topology on the coefficients), and a change of convention changes the formulas, though not the final answers. We fix them once.

\begin{convention}\label{conv:cohomology:conventions}
\begin{enumerate}[label=(\alph*)]
\item $\ell$ is a prime. Where $\ell$ must be odd we say so; this happens in exactly two places: the model computation of
\Cref{sec:cohomology:model}, where $2$ must be a unit of $\Zl$, and the vanishing of the cohomology of a group of order $2$ on $\ell$-adic modules
(\Cref{cor:cohomology:z2}).
\item Groups act on the left. A cochain complex has differentials raising the degree. The connecting homomorphism of a short exact sequence of
complexes is the one constructed in \Cref{thm:cohomology:les}.
\item All cohomology of a profinite group is \emph{continuous cochain cohomology} (\Cref{def:cohomology:continuous}). The coefficient modules are of three
types: (D) discrete modules, finite or $\ell$-primary torsion, such as $W=T\otimes\Qlmod$; (C) finitely generated $\Zl$-modules with the $\ell$-adic
topology and a continuous action; (V) finite-dimensional $\Ql$-vector spaces with a continuous action.
\item $\Frob$ is always an \emph{arithmetic} Frobenius: on $\overline\F_q$ it is $x\mapsto x^q$, and $\chi_\ell(\Frob_q)=q$ on $\Zl(1)=\varprojlim_k\mu_{\ell^k}$ for $q\ne\ell$.
For a prime $q\ne\ell$, $G_{\Q_q}/I_q=\Gal(\overline\F_q/\F_q)\simeq\widehat\Z$ is topologically generated by the image of $\Frob_q$.
\item The conjugation action of $G/H$ on $H^n(H,M)$ is
\[
(\sigma\cdot c)(h_1,\dots,h_n)=\sigma\,c(\sigma^{-1}h_1\sigma,\dots,\sigma^{-1}h_n\sigma),
\]
verbatim the paper's formula ``$(\sigma\cdot c)(i)=\sigma c(\sigma^{-1}i\sigma)$'' in the proof of Lemma lem:tamfactor (\Cref{prop:cohomology:conjugation}).
\item The tame relation, proved in \Cref{ch:local} and imported here: for a prime $q\ne\ell$ and a topological generator $t$ of the pro-$\ell$ tame inertia
quotient of $G_{\Q_q}$,
\[
\Frob\,t\,\Frob^{-1}=t^{q},\qquad\text{equivalently}\qquad \Frob^{-1}t\,\Frob=t^{q^{-1}},\quad q^{-1}\in\Zl^\times .
\]
\item $\ad T$ is the module of trace-zero endomorphisms of a representation $T$, with $g\cdot x=\rho(g)x\rho(g)^{-1}$. For $T$ free of rank $2$ with basis
$e_1,e_2$ its basis is $e=e_1\otimes e_2^\vee$, $h=e_1\otimes e_1^\vee-e_2\otimes e_2^\vee$, $f=e_2\otimes e_1^\vee$, as in the paper.
\item $A^*$ is the Pontryagin dual of a module $A$; a star after a \emph{map} of pairs, as in $(f,\varphi)^*$, is the induced map on cohomology. We
write $A_{\mathrm{tor}}$ for the torsion submodule of a $\Zl$-module, as the paper does, and $\len=\len_{\Zl}$.
\end{enumerate}
\end{convention}

\begin{warning}\label{warn:cohomology:frobenius}
\Cref{ch:blochkato,ch:lfunctions} compute $L$-factors with the \emph{geometric} Frobenius $\Frob^{-1}$:
$L_p(M,s)=\det(1-\Frob_p^{-1}p^{-s}\mid M^{I_p})^{-1}$. An eigenvalue $\lambda$ of the arithmetic Frobenius is the eigenvalue $\lambda^{-1}$ of the
geometric one. In this chapter only the arithmetic Frobenius occurs. The two conventions are both used in the paper: the proof of Lemma lem:tamfactor,
with $\rho(\Frob)=\sm{\varepsilon q}{c'}{0}{\varepsilon}$, is written with the arithmetic Frobenius, and Section 1.1 uses the geometric one. The paragraph after Lemma lem:dfg21 uses both: the quotient
$(V^D)^{I_p}/(\Frob_p-1)$ is written with the arithmetic Frobenius and the determinant ``$\det(1-\Frob_p^{-1}\mid(V^D)^{I_p})=L_p(A_f,1)^{-1}$'' with the geometric
one. The quotient does not depend on the choice, because $\Frob_p^{-1}-1=-\Frob_p^{-1}(\Frob_p-1)$ has the same image as $\Frob_p-1$; and the vanishing statements we prove
(\Cref{prop:cohomology:unramified}) do not depend on the convention, since $1$ is an eigenvalue of $\Frob$ if and only if it is an eigenvalue of $\Frob^{-1}$.
\end{warning}

The chapter proceeds from the general to the particular. Sections 2 and 3 recall the homological algebra (complexes, the long exact sequence,
derived functors in outline). Sections 4 to 6 define the cohomology of an abstract group by cochains and prove the facts the paper uses: the meaning
of $H^0$, $H^1$, $H^2$, functoriality, restriction, inflation and corestriction, the conjugation action, inflation-restriction, the vanishing for
coefficients on which the group order is invertible. Sections 7 to 9 pass to profinite groups with continuous cochains, which is the setting of Galois
cohomology, and compute the cohomology of procyclic groups; the identity $H^1(\langle t\rangle,M)=M/(t-1)M$ is \Cref{thm:cohomology:procyclic-compact}.
Section 10 is the group-cohomological content of Lemma lem:tamfactor, done line by line. Section 11 treats cup products and lengths in exact
sequences, Section 12 the notation of Galois cohomology and a preview of \Cref{ch:galcoh}.

\section{Complexes, cohomology and the long exact sequence}\label{sec:cohomology:complexes}

Every long exact sequence in the rest of the book (the Kummer sequences of \Cref{ch:galcoh}, the sequences relating $T$, $V$ and $W$, the exact
sequence of Lemma lem:dfg21 in \Cref{ch:blochkato}) is an instance of \Cref{thm:cohomology:les}. This section proves it, after recalling the snake
lemma from \Cref{ch:modules}. Throughout, $R$ is a commutative ring; in practice $R=\Z$ or $R=\Zl$.

\begin{definition}[cochain complexes]\label{def:cohomology:complex}
A \emph{cochain complex} of $R$-modules is a sequence $C^\bullet=(C^n,d^n)_{n\in\Z}$ of $R$-modules $C^n$ and $R$-linear maps $d^n\colon C^n\to C^{n+1}$
with $d^{n+1}d^n=0$ for all $n$. Its \emph{cocycles} are $Z^n=\ker d^n$, its \emph{coboundaries} $B^n=\im d^{n-1}$; since $d^nd^{n-1}=0$, $B^n\subseteq Z^n$,
and the \emph{cohomology} is $H^n(C^\bullet)=Z^n/B^n$. A \emph{morphism} $f\colon C^\bullet\to D^\bullet$ is a family of $R$-linear maps $f^n\colon C^n\to D^n$
with $d^nf^n=f^{n+1}d^n$. Two morphisms $f,g$ are \emph{homotopic} if there are $R$-linear maps $s^n\colon C^n\to D^{n-1}$ with
$f^n-g^n=d^{n-1}s^n+s^{n+1}d^n$ for all $n$. A sequence $A^\bullet\to B^\bullet\to C^\bullet$ of morphisms is \emph{exact} if it is exact in every degree.
\end{definition}

\begin{lemma}\label{lem:cohomology:complex-basics}
\begin{enumerate}[label=(\roman*)]
\item A morphism $f$ maps $Z^n(C)$ into $Z^n(D)$ and $B^n(C)$ into $B^n(D)$, hence induces an $R$-linear map $H^n(f)\colon H^n(C)\to H^n(D)$, and
$H^n(g\circ f)=H^n(g)\circ H^n(f)$, $H^n(\id)=\id$.
\item Homotopic morphisms induce the same map on cohomology.
\end{enumerate}
\end{lemma}

\begin{proof}
(i) If $dx=0$ then $d(fx)=f(dx)=0$, so $f(Z^n(C))\subseteq Z^n(D)$; and $f(dy)=d(fy)$, so $f(B^n(C))\subseteq B^n(D)$. Hence $H^n(f)[x]=[fx]$ is well
defined: if $x'=x+dy$ is another representative, $fx'=fx+d(fy)$ has the same class. On representatives, $H^n(g\circ f)[x]=[g(f(x))]=H^n(g)(H^n(f)[x])$ and
$H^n(\id)[x]=[x]$. (ii) For $x\in Z^n(C)$,
$f(x)-g(x)=d(s^nx)+s^{n+1}(dx)=d(s^nx)\in B^n(D)$, so $[f(x)]=[g(x)]$. Note that $s$ need not commute with $d$; only the combination $ds+sd$ enters.
\end{proof}

The definition of a homotopy may look unmotivated. Its role is exactly (ii): it is the weakest relation between $f$ and $g$ that forces equal maps on
cohomology by the one-line argument above, and it is the relation that the comparison theorem (\Cref{prop:cohomology:resolutions}) produces between two
lifts of the same map. This is why group cohomology, defined through a choice of resolution, does not depend on that choice.

We recall the snake lemma, proved in \Cref{ch:modules}, because the construction of its connecting map is the template for every connecting map that
follows.

\begin{lemma}[the snake lemma]\label{lem:cohomology:snake}
Let
\[
\begin{array}{ccccccc}
 & A' & \xrightarrow{\ i\ } & A & \xrightarrow{\ p\ } & A'' & \to 0\\
 & \big\downarrow{\scriptstyle f'} & & \big\downarrow{\scriptstyle f} & & \big\downarrow{\scriptstyle f''} & \\
0\to & B' & \xrightarrow{\ j\ } & B & \xrightarrow{\ r\ } & B'' &
\end{array}
\]
be a commutative diagram of $R$-modules with exact rows. For $a''\in\ker f''$ choose $a\in A$ with $p(a)=a''$; then $r(f(a))=f''(a'')=0$, so $f(a)=j(b')$
for a unique $b'\in B'$; put $\delta(a'')=b'+f'(A')$. Then $\delta\colon\ker f''\to\coker f'$ is well defined and $R$-linear, and
\[
\ker f'\to\ker f\to\ker f''\xrightarrow{\ \delta\ }\coker f'\to\coker f\to\coker f''
\]
is exact. If $i$ is injective the sequence starts with $0\to\ker f'$; if $r$ is surjective it ends with $\coker f''\to0$.
\end{lemma}

\begin{proof}
This is proved in \Cref{ch:modules}; see also \cite[Section 1.3]{Weibel}.
\end{proof}

\begin{theorem}[the long exact sequence]\label{thm:cohomology:les}
Let $0\to A^\bullet\xrightarrow{i}B^\bullet\xrightarrow{p}C^\bullet\to0$ be a short exact sequence of cochain complexes of $R$-modules.
\begin{enumerate}[label=(\roman*)]
\item For each $n$ there is an $R$-linear map $\delta^n\colon H^n(C^\bullet)\to H^{n+1}(A^\bullet)$ defined as follows. Represent a class by $c\in Z^n(C)$; lift it
to $b\in B^n$ with $p(b)=c$; then $p(db)=dc=0$, so $db=i(a)$ for a unique $a\in A^{n+1}$; this $a$ is a cocycle, and $\delta^n[c]=[a]$.
\item The sequence
\[
\cdots\to H^n(A)\xrightarrow{i_*}H^n(B)\xrightarrow{p_*}H^n(C)\xrightarrow{\delta^n}H^{n+1}(A)\xrightarrow{i_*}H^{n+1}(B)\xrightarrow{p_*}\cdots
\]
is exact.
\item (Naturality.) A commutative diagram of short exact sequences of complexes
$(0\to A\to B\to C\to0)\to(0\to\tilde A\to\tilde B\to\tilde C\to0)$, given by morphisms $\alpha,\beta,\gamma$, induces a commutative ladder between the two long
exact sequences; in particular $\delta^n\circ H^n(\gamma)=H^{n+1}(\alpha)\circ\delta^n$.
\end{enumerate}
\end{theorem}

\begin{proof}
(i) Since $i$ is injective and $i(da)=d(i(a))=ddb=0$, $da=0$. Independence of the choices: another lift is $b+i(x)$ with $x\in A^n$, which changes $a$ to
$a+dx$; another representative is $c+dy$ with $y\in C^{n-1}$; lift $y$ to $\tilde y\in B^{n-1}$, then $b+d\tilde y$ lifts $c+dy$ and has the same
differential $db$. So $[a]$ depends only on $[c]$. Linearity: if $(c_1,b_1,a_1)$ and $(c_2,b_2,a_2)$ are choices for $[c_1]$ and $[c_2]$, and
$\lambda\in R$, then $(c_1+\lambda c_2,\,b_1+\lambda b_2,\,a_1+\lambda a_2)$ is a valid choice for $[c_1]+\lambda[c_2]$, because $p$, $i$ and $d$ are $R$-linear; by
independence of the choices, $\delta^n([c_1]+\lambda[c_2])=[a_1]+\lambda[a_2]$.

(ii) A clean route is the snake lemma. For each of the three complexes $X^\bullet\in\{A^\bullet,B^\bullet,C^\bullet\}$ the differential induces a map
$\bar d\colon X^n/dX^{n-1}\to Z^{n+1}(X)$, $x+dX^{n-1}\mapsto dx$ (well defined since $dd=0$, and landing in $Z^{n+1}$ for the same reason). Its kernel is
$\{x:dx=0\}/dX^{n-1}=Z^n/B^n=H^n(X)$, and its cokernel is $Z^{n+1}(X)/dX^{n}=H^{n+1}(X)$. The rows
$A^n/dA^{n-1}\to B^n/dB^{n-1}\to C^n/dC^{n-1}\to0$ and $0\to Z^{n+1}(A)\to Z^{n+1}(B)\to Z^{n+1}(C)$ are exact. In the first row, surjectivity at the right end holds
because $p$ is surjective. For exactness in the middle, let $x\in B^n$ map to $0$ in $C^n/dC^{n-1}$, that is $p(x)=dy$ with $y\in C^{n-1}$; lift $y$ to $\tilde y\in B^{n-1}$, so
that $x-d\tilde y\in\ker p=i(A^n)$, and $x$ has the same class as $x-d\tilde y$, which comes from $A^n$. (The composite of the two maps is $0$ because $pi=0$.) In the
second row, injectivity is inherited from $i$; and a cocycle $z$ of $B^\bullet$ with $p(z)=0$ is $i(a)$, with $da=0$ because $i(da)=dz=0$ and $i$ is injective. The two
rows and the three vertical maps $\bar d$ form a commutative diagram, because $i$ and $p$ commute with $d$. The snake lemma applied to it gives the exact sequence
\[
H^n(A)\to H^n(B)\to H^n(C)\xrightarrow{\ \delta\ }H^{n+1}(A)\to H^{n+1}(B)\to H^{n+1}(C),
\]
and the snake construction of $\delta$ is exactly (i). Splicing these six-term sequences for all $n$ gives (ii).

(iii) The maps $i_*$, $p_*$ commute with $\alpha_*,\beta_*,\gamma_*$ by functoriality. For $\delta$: if $(c,b,a)$ are the choices in (i) for $[c]$, then
$(\gamma c,\beta b,\alpha a)$ are valid choices for $[\gamma c]$, since $\tilde p\beta=\gamma p$ and $\tilde i\alpha=\beta i$ and $\beta$ commutes with $d$:
indeed $\tilde p(\beta b)=\gamma(pb)=\gamma c$ and $d(\beta b)=\beta(db)=\beta(i(a))=\tilde i(\alpha a)$. So $\delta^n[\gamma c]=[\alpha a]=H^{n+1}(\alpha)\delta^n[c]$.
\end{proof}

The recipe of (i) is worth memorising in the form ``lift, differentiate, pull back'': every connecting map in the book, from the Kummer map of
\Cref{ch:galcoh} to the map $\delta\colon H^0(G,W)\to H^1(G,T)$ of \Cref{prop:cohomology:les-cont}, is computed this way. The class $\delta[c]$ is nonzero
exactly when the lift $b$ cannot be chosen to be a cocycle. \Cref{prop:cohomology:h0-h1}(iv) below is the first instance: there $\delta(m'')$ is the
coboundary of a lift $m$ of $m''$, a coboundary in $M$ that takes values in $M'$, where it is a cocycle but in general not a coboundary. For example, for
the group $G=\{1,\sigma\}$ acting on $M=\Z$ by $-1$ and the sequence $0\to2\Z\to\Z\to\Z/2\to0$, the element $1\in(\Z/2)^G$ lifts to $1\in\Z$, and
$\delta(1)$ is the class of $g\mapsto g\cdot1-1$, which is $\sigma\mapsto-2$; this is not of the form $\sigma\mapsto\sigma x-x=-2x$ with $x\in2\Z$, so $\delta(1)\ne0$.

\begin{example}\label{ex:cohomology:two-term}
Let $M$ be an $R$-module and $u$ an $R$-linear endomorphism of $M$. The complex $0\to M\xrightarrow{u-1}M\to0$ (in degrees $0$ and $1$) has
$H^0=\ker(u-1)$ and $H^1=M/(u-1)M$. We shall see in \Cref{thm:cohomology:procyclic-compact} that, for a procyclic group $\Gamma$ with a surjection onto $\Zl$
and generator $F$ acting by $u$ on a module $M$ of type (C), this complex computes $H^\bullet(\Gamma,M)$ in all degrees. For a finite cyclic group
$\langle t\rangle$ of order $m$ the analogous complex is not this one: it is the infinite periodic complex
$M\xrightarrow{t-1}M\xrightarrow{N}M\xrightarrow{t-1}M\xrightarrow{N}\cdots$ with $N=1+t+\dots+t^{m-1}$, which is why $H^1(\langle t\rangle,M)=\ker N/(t-1)M$
(\Cref{thm:cohomology:cyclic}). The difference between the two answers is the subject of \Cref{warn:cohomology:procyclic}.
\end{example}

\section[Resolutions, derived functors, Ext and Tor]{Resolutions, derived functors, $\Ext$ and $\Tor$ in outline}\label{sec:cohomology:derived}

This section is background. The paper never uses $\Ext$ or $\Tor$ by name; but $H^n(G,M)$ is a derived functor
(\Cref{thm:cohomology:as-ext}), this is what makes it independent of the cochains chosen to compute it, and $\Ext^1$ is the conceptual home of the
``cocycle in the corner'' of a reducible two-dimensional representation, which the paper uses for the Tate curve (\Cref{prop:cohomology:corner}). We
state the facts with references and compute the examples over $\Z$ and $\Zl$ that occur later.

\begin{definition}[projective and injective resolutions]\label{def:cohomology:resolutions}
\sloppy An $R$-module $P$ is \emph{projective} if $\Hom_R(P,-)$ is exact (equivalently, a direct summand of a free module); $I$ is \emph{injective} if
$\Hom_R(-,I)$ is exact. A \emph{projective resolution} of $M$ is an exact sequence $\cdots\to P_1\to P_0\to M\to0$ with every $P_i$ projective; an
\emph{injective resolution} is an exact sequence $0\to M\to I^0\to I^1\to\cdots$ with every $I^i$ injective.
\end{definition}

\begin{proposition}[existence and comparison]\label{prop:cohomology:resolutions}
\begin{enumerate}[label=(\roman*)]
\item Every $R$-module has a resolution by free modules, and an injective resolution.
\item (Comparison theorem.) A map $M\to N$ lifts to a morphism between given projective resolutions of $M$ and $N$, unique up to homotopy; dually for
injective resolutions. Hence any two projective (resp. injective) resolutions of $M$ are homotopy equivalent.
\item If $R$ is a principal ideal domain (for instance $\Z$ or $\Zl$), every module has a free resolution $0\to P_1\to P_0\to M\to0$ of length one
(submodules of free modules are free), and the divisible modules, such as $\Q/\Z$, $\Qlmod$ and $\Ql$, are injective.
\end{enumerate}
\end{proposition}

\begin{proof}
We do not prove this; see \cite[Sections 2.2, 2.3]{Weibel}. For (i) with free modules, map the free module on the elements of $M$ onto $M$ and repeat with
the kernel.
\end{proof}

\begin{definition}[$\Ext$ and $\Tor$]\label{def:cohomology:ext-tor}
For $R$-modules $M$, $N$, a projective resolution $P_\bullet\to M$ and an injective resolution $N\to I^\bullet$,
\[
\Ext^n_R(M,N)=H^n\bigl(\Hom_R(P_\bullet,N)\bigr)=H^n\bigl(\Hom_R(M,I^\bullet)\bigr),\qquad \Tor_n^R(M,N)=H_n(P_\bullet\otimes_RN).
\]
\end{definition}

\begin{proposition}[properties of $\Ext$ and $\Tor$]\label{prop:cohomology:ext-tor}
\begin{enumerate}[label=(\roman*)]
\item The two descriptions of $\Ext$ agree and are independent of the resolutions; $\Ext^0=\Hom$ and $\Tor_0=\otimes$.
\item A short exact sequence in either variable gives a long exact sequence of $\Ext$ (resp. $\Tor$).
\item $\Ext^n_R(M,N)=0$ for $n\ge1$ if $M$ is projective or $N$ is injective. Over a principal ideal domain $\Ext^n=\Tor_n=0$ for $n\ge2$.
\item $\Ext^1_\Z(\Z/m,N)\simeq N/mN$, $\Tor_1^\Z(\Z/m,N)\simeq N[m]$, and $\Ext^1_{\Zl}(\Z/\ell^a,\Z/\ell^b)\simeq\Z/\ell^{\min(a,b)}$.
\end{enumerate}
\end{proposition}

\begin{proof}
(i) to (iii) are stated; see \cite[Sections 2.4 to 2.7 and 3.1 to 3.3]{Weibel}. The vanishing for $n\ge2$ over a principal ideal domain follows from
resolutions of length one. (iv) Use the resolution $0\to\Z\xrightarrow{m}\Z\to\Z/m\to0$. Applying $\Hom(-,N)$ gives the complex
$N\xrightarrow{m}N$ in degrees $0$, $1$, whose $H^1$ is $N/mN$; tensoring with $N$ gives $N\xrightarrow{m}N$ in degrees $1$, $0$, whose $H_1$ is $N[m]$. Over
$\Zl$ the resolution $0\to\Zl\xrightarrow{\ell^a}\Zl\to\Z/\ell^a\to0$ gives $\Ext^1_{\Zl}(\Z/\ell^a,N)=N/\ell^aN$, and $(\Z/\ell^b)/\ell^a=\Z/\ell^{\min(a,b)}$.
\end{proof}

\begin{theorem}[$\Ext^1$ classifies extensions]\label{thm:cohomology:ext1}
$\Ext^1_R(M,N)$ is in natural bijection with the set of equivalence classes of extensions $0\to N\to E\to M\to0$ of $R$-modules, two extensions being
equivalent if there is an isomorphism $E\to E'$ inducing the identity on $N$ and on $M$. The split extension $N\oplus M$ corresponds to $0$, and the
group law corresponds to the Baer sum. The class of an extension is the image of $\id_M$ under the connecting map $\Hom_R(M,M)\to\Ext^1_R(M,N)$ of the
long exact sequence of $\Hom_R(M,-)$.
\end{theorem}

\begin{proof}
We do not prove this; see \cite[Section 3.4]{Weibel}. \Cref{prop:cohomology:corner} below proves by hand the case that the paper uses.
\end{proof}

A first example shows the bijection at work. By \Cref{prop:cohomology:ext-tor}(iv), $\Ext^1_\Z(\Z/2,\Z/2)=(\Z/2)/2(\Z/2)=\Z/2$ has two elements, and there are
exactly two extensions $0\to\Z/2\to E\to\Z/2\to0$ up to equivalence: $E$ is a group of order $4$, so $E\simeq\Z/2\times\Z/2$ (the split extension, class $0$) or
$E\simeq\Z/4$ (with the sub $2\Z/4\Z$, the nonzero class). The extension class records whether a lift of the generator of the quotient can be chosen to be
killed by $2$. In the same way $\Ext^1_{\Zl}(\Z/\ell^a,\Z/\ell^b)=\Z/\ell^{\min(a,b)}$ counts the ways of gluing a cyclic group of order $\ell^a$ on top of one of
order $\ell^b$ ($a,b\ge1$); the cyclic extensions $E\simeq\Z/\ell^{a+b}$ are those whose class is a generator. The case the paper needs is the
equivariant analogue: the module is $\Zl^2$, the extension is one of Galois modules, and the class lives in a group cohomology $H^1$.

\begin{remark}\label{rem:cohomology:dvr-degree-one}
Over a discrete valuation ring such as $\Zl$ the homological algebra of modules stops in degree one. This is one reason why the cohomology of a
procyclic group with $\ell$-adic coefficients is so simple: in \Cref{sec:cohomology:procyclic} the two-term complex of \Cref{ex:cohomology:two-term}
computes it. Group cohomology does not stop in degree one in general: a finite cyclic group has nonzero cohomology in infinitely many degrees.
\end{remark}

\section[Group cohomology by cochains]{Group cohomology by cochains: $H^0$, $H^1$, $H^2$}\label{sec:cohomology:cochains}

This section defines $H^n(G,M)$ for an abstract group $G$ and proves what the paper uses about the low degrees. Lemma lem:tamfactor calls $\kappa$
``the Kummer cocycle'' and reads a reducible two-dimensional representation as a $1$-cocycle (\Cref{prop:cohomology:corner}); Lemma lem:h0 computes an
$H^0$ through its $\ell^k$-torsion (\Cref{lem:cohomology:h0-torsion}); Lemma lem:dfg21 is a sequence of connecting maps
(\Cref{prop:cohomology:les-groups}). The profinite version, with continuous cochains, is \Cref{sec:cohomology:profinite}; every definition and proof of this
section carries over to it with the word ``continuous'' added, and we say so there.

\begin{definition}[$G$-modules]\label{def:cohomology:gmodule}
Let $G$ be a group. A \emph{$G$-module} is an abelian group $M$ with a left action of $G$ by group automorphisms, that is, a left module over the group
ring $\Z[G]$. The \emph{invariants} are $M^G=\{m:gm=m\text{ for all }g\}$; the \emph{coinvariants} are $M_G=M/I_GM$, where the augmentation ideal $I_G$ is
the kernel of $\Z[G]\to\Z$, $\sum a_gg\mapsto\sum a_g$, generated by the elements $g-1$. $M$ is \emph{trivial} if $gm=m$ for all $g,m$. For finite $G$ the
\emph{norm} is $N_G=\sum_{g\in G}g$, and $N_GM\subseteq M^G$. For $G$-modules $M$, $N$, the group $\Hom_\Z(M,N)$ is a $G$-module by
$(g\varphi)(m)=g\varphi(g^{-1}m)$, and $\Hom_G(M,N)=\Hom_\Z(M,N)^G$. A map of $G$-modules is a $\Z[G]$-linear map.
If $M$ is a $\Zl$-module and $G$ acts $\Zl$-linearly, $M$ is a $\Zl[G]$-module, and every cohomology group defined below is a $\Zl$-module, every map
$\Zl$-linear.
\end{definition}

The invariants and the coinvariants are the largest submodule and the largest quotient on which $G$ acts trivially; group cohomology measures how far
$M\mapsto M^G$ is from being exact. A first example: let $G=\langle\sigma\rangle$ have order $2$ and act on $\Z^2$ by swapping the coordinates. Then
$M^G=\{(a,a)\}$, $I_GM=(\sigma-1)M=\{(b,-b)\}$, $M_G=\Z^2/\{(b,-b)\}\simeq\Z$ via $(x,y)\mapsto x+y$, and $N_GM=\{(a+b,a+b)\}=M^G$. For the sign module $\Z^-$ (with
$\sigma$ acting by $-1$) instead, $M^G=0$, $M_G=\Z/2$ and $N_G=0$. The examples of the paper are $\Zl$-modules such as $T_\ell E$, $\ad T$ and $\Zl(1)$, with $G$ a
Galois group acting through a representation.

\begin{definition}[inhomogeneous cochains]\label{def:cohomology:cochains}
For $n\ge0$ let $C^n(G,M)$ be the group of all maps $G^n\to M$ ($C^0(G,M)=M$), and define $d^n\colon C^n(G,M)\to C^{n+1}(G,M)$ by
\begin{multline*}
(d^n\varphi)(g_1,\dots,g_{n+1})=g_1\varphi(g_2,\dots,g_{n+1})\\
+\sum_{i=1}^{n}(-1)^i\varphi(g_1,\dots,g_ig_{i+1},\dots,g_{n+1})+(-1)^{n+1}\varphi(g_1,\dots,g_n).
\end{multline*}
In low degrees:
\begin{align*}
(d^0m)(g)&=gm-m,\\
(d^1\varphi)(g,h)&=g\varphi(h)-\varphi(gh)+\varphi(g),\\
(d^2\psi)(g,h,k)&=g\psi(h,k)-\psi(gh,k)+\psi(g,hk)-\psi(g,h).
\end{align*}
The \emph{cohomology} of $G$ with coefficients in $M$ is $H^n(G,M)=H^n(C^\bullet(G,M))$ (with $C^n=0$ for $n<0$). Elements of $Z^n(G,M)=\ker d^n$ are
\emph{$n$-cocycles}, of $B^n(G,M)=\im d^{n-1}$ \emph{$n$-coboundaries}.
\end{definition}

The formula for $d^n$ can be read off from the low degrees, which are the ones the book uses. In degree $0$, $d^0m=0$ says that $m$ is invariant. In degree $1$,
$d^1\varphi=0$ says $\varphi(gh)=\varphi(g)+g\varphi(h)$: $\varphi$ is a homomorphism ``twisted by the action'', the rule satisfied by the corner entry of a triangular
representation (\Cref{prop:cohomology:corner}) and by the function $g\mapsto gm-m$. In degree $2$, $d^2\psi=0$ is the associativity constraint on the factor set of a
group extension (\Cref{thm:cohomology:h2-ext}). The general $d^n$ alternates between applying the action to the first variable, multiplying adjacent variables,
and dropping the last variable; this pattern is what makes $d\circ d=0$, by a cancellation in pairs. The adjective ``inhomogeneous'' distinguishes these cochains
from the homogeneous cochains, the $G$-equivariant maps $G^{n+1}\to M$, which compute the same groups (\Cref{thm:cohomology:as-ext}; the two are related by
$\varphi(g_1,\dots,g_n)=\Phi(1,g_1,g_1g_2,\dots,g_1\cdots g_n)$).

\begin{lemma}\label{lem:cohomology:dd}
$d^{n+1}d^n=0$ for all $n$, so $C^\bullet(G,M)$ is a cochain complex. A map of $G$-modules $f\colon M\to N$ induces a morphism of complexes
$\varphi\mapsto f\circ\varphi$, hence maps $f_*\colon H^n(G,M)\to H^n(G,N)$; so $H^n(G,-)$ is a covariant additive functor.
\end{lemma}

\begin{proof}
In degree $0$: $(d^1d^0m)(g,h)=g(hm-m)-(ghm-m)+(gm-m)=0$. In degree $1$, expanding the four terms of $(d^2d^1\varphi)(g,h,k)$:
\begin{multline*}
g\bigl(h\varphi(k)-\varphi(hk)+\varphi(h)\bigr)-\bigl(gh\varphi(k)-\varphi(ghk)+\varphi(gh)\bigr)\\
+\bigl(g\varphi(hk)-\varphi(ghk)+\varphi(g)\bigr)-\bigl(g\varphi(h)-\varphi(gh)+\varphi(g)\bigr)=0,
\end{multline*}
every term cancelling against another. Degree $2$ is \Cref{ex:cohomology:dd3}. In general the identity follows from \Cref{thm:cohomology:as-ext}: $C^\bullet(G,M)$
is $\Hom_{\Z[G]}(F_\bullet,M)$ for a complex $F_\bullet$ (the bar resolution) whose differential squares to zero by a direct computation
\cite[Chapter III, Section 1]{Brown}, \cite[Chapter I, Section 2]{NSW}. Since $f$ is $G$-linear, $f(g\varphi(\dots))=gf(\varphi(\dots))$, so $f\circ-$ commutes
with $d$.
\end{proof}

\begin{theorem}[cohomology as a derived functor]\label{thm:cohomology:as-ext}
Let $\Z$ have the trivial $G$-action. There is a resolution $F_\bullet\to\Z$ of $\Z$ by free $\Z[G]$-modules, $F_n$ free on $G^n$ (the bar resolution), with
$\Hom_{\Z[G]}(F_\bullet,M)\simeq C^\bullet(G,M)$ as complexes. Hence $H^n(G,M)\simeq\Ext^n_{\Z[G]}(\Z,M)$, $H^n(G,-)$ is the $n$-th right derived functor of
$M\mapsto M^G=\Hom_{\Z[G]}(\Z,M)$, and any projective resolution of $\Z$ over $\Z[G]$ computes it.
\end{theorem}

\begin{proof}
We do not prove this; see \cite[Chapter I, Section 5 and Chapter III, Section 1]{Brown} and \cite[Section 6.5]{Weibel}.
\end{proof}

\begin{proposition}[$H^0$ and $H^1$]\label{prop:cohomology:h0-h1}
Let $M$ be a $G$-module.
\begin{enumerate}[label=(\roman*)]
\item $H^0(G,M)=M^G$.
\item $Z^1(G,M)$ is the group of \emph{crossed homomorphisms}, the maps $\varphi\colon G\to M$ with $\varphi(gh)=\varphi(g)+g\varphi(h)$, and $B^1(G,M)$ the group of
\emph{principal} ones, $g\mapsto gm-m$ for some $m\in M$. A $1$-cocycle has $\varphi(1)=0$ and $\varphi(g^{-1})=-g^{-1}\varphi(g)$.
\item If $G$ acts trivially on $M$, then $B^1=0$ and $H^1(G,M)=\Hom(G,M)=\Hom(G^{\mathrm{ab}},M)$. In particular $H^1(G,M)=0$ if $G$ is finite, $M$ is
torsion-free and the action is trivial.
\item For an exact sequence $0\to M'\to M\xrightarrow{p}M''\to0$ of $G$-modules the sequence
\[
0\to M'^G\to M^G\to M''^G\xrightarrow{\ \delta\ }H^1(G,M')\to H^1(G,M)\to H^1(G,M'')
\]
is exact, where $\delta(m'')$ is the class of $g\mapsto gm-m$ for any $m\in M$ with $p(m)=m''$ (this cocycle takes values in $M'$).
\end{enumerate}
\end{proposition}

\begin{proof}
(i) $C^{-1}=0$ and $(d^0m)(g)=gm-m$, so $H^0=\ker d^0=M^G$. (ii) $d^1\varphi=0$ is the identity $\varphi(gh)=\varphi(g)+g\varphi(h)$. Putting $g=h=1$ gives
$\varphi(1)=2\varphi(1)$, so $\varphi(1)=0$; putting $h=g^{-1}$ gives $0=\varphi(g)+g\varphi(g^{-1})$. (iii) With trivial action $gm-m=0$ and the cocycle
identity reads $\varphi(gh)=\varphi(g)+\varphi(h)$; a homomorphism to an abelian group factors through $G^{\mathrm{ab}}$. If $G$ is finite every $g$ has finite
order $k$, so $k\varphi(g)=\varphi(g^k)=0$ and $\varphi(g)=0$ when $M$ is torsion-free. (iv) is the beginning of the long exact sequence of
\Cref{prop:cohomology:les-groups}; the formula for $\delta$ is the recipe of \Cref{thm:cohomology:les}(i): lift $m''$ to $m\in C^0(G,M)=M$ and apply $d^0$.
Concretely, $p(gm-m)=gm''-m''=0$ because $m''$ is invariant, so $gm-m$ lies in $M'$; the cochain $g\mapsto gm-m$ is a cocycle (it is $d^0m$ in $M$, and $d^1d^0=0$);
another lift $m+m'$ changes it by $d^0m'$, a coboundary in $M'$. The class $\delta(m'')$ is $0$ if and only if $gm-m=gm'-m'$ for some $m'\in M'$, that is, if and only if
$m''$ has the invariant lift $m-m'$; this is exactness at $M''^G$.
\end{proof}

\begin{remark}\label{rem:cohomology:cocycle-words}
The words come from topology. A reader who has met projective representations in quantum mechanics has seen $2$-cocycles: a map
$U\colon G\to\mathrm U(n)$ with $U(g)U(h)=\omega(g,h)U(gh)$ forces $\omega$ to be a $2$-cocycle with values in the circle group (trivial action), and rescaling
$U(g)$ by phases changes $\omega$ by a coboundary. Crossed homomorphisms appear in the same way as the off-diagonal entry of a triangular representation
(\Cref{prop:cohomology:corner}).
\end{remark}

\begin{remark}[non-abelian $H^1$]\label{rem:cohomology:nonabelian}
If $A$ is a group, not necessarily abelian, on which $G$ acts by automorphisms, one still defines $Z^1(G,A)$ as the set of maps with
$\varphi(gh)=\varphi(g)\,{}^g\varphi(h)$, and $H^1(G,A)$ as its quotient by the relation $\varphi\sim\psi$ if $\psi(g)=a^{-1}\varphi(g)\,{}^ga$ for some $a\in A$. This is a
pointed set, not a group, and there is no $H^2$ in general. For an algebraic group $G_0$ over a field $k$ one writes $H^1(k,G_0)=H^1(\Gal(k^{\mathrm{sep}}/k),G_0(k^{\mathrm{sep}}))$;
it classifies principal homogeneous spaces of $G_0$ over $k$ \cite[Chapter I, Section 5 and Chapter III]{SerreGaloisCoh}. This is the sense of
``$H^1(k,G)=0$'' (Lang's theorem for a connected group over a finite field) in \Cref{ch:reduction}. Nothing else in the book uses non-abelian cohomology.
\end{remark}

\begin{proposition}[the cocycle in the corner]\label{prop:cohomology:corner}
Let $R$ be a commutative ring, $G$ a group, $\eta_1,\eta_2\colon G\to R^\times$ characters, and let $R(\eta)$ denote $R$ with $g$ acting by multiplication by
$\eta(g)$.
\begin{enumerate}[label=(\roman*)]
\item A map $\rho\colon G\to\GL_2(R)$ of the form $\rho(g)=\Mtwo{\eta_1(g)}{\kappa(g)}{0}{\eta_2(g)}$ is a homomorphism if and only if
$b(g)=\kappa(g)\eta_2(g)^{-1}$ is a $1$-cocycle with values in $R(\eta_1\eta_2^{-1})$:
\[
b(gh)=b(g)+\eta_1\eta_2^{-1}(g)\,b(h).
\]
\item Replacing the basis vector $e_2$ by $e_2+re_1$ ($r\in R$) replaces $b$ by $b+d^0r$, where $(d^0r)(g)=(\eta_1\eta_2^{-1}(g)-1)r$. Scaling $e_1$ or
$e_2$ by units multiplies $b$ by a unit.
\item $\rho$ is split, that is, $Re_1$ has a $G$-stable complement, if and only if the class of $b$ in $H^1(G,R(\eta_1\eta_2^{-1}))$ is $0$.
\item Every class arises. Hence the isomorphism classes of extensions $0\to R(\eta_1)\to E\to R(\eta_2)\to0$ of $R[G]$-modules (with $E$ free of rank two over
$R$, and isomorphisms inducing the identity on the sub and the quotient) are in bijection with $H^1(G,R(\eta_1\eta_2^{-1}))$; this is the case of
\Cref{thm:cohomology:ext1} that we need.
\item (The paper's shape.) Let $\rho(g)=\Mtwo{\chi_\ell(g)\varepsilon(g)}{\varepsilon(g)\kappa(g)}{0}{\varepsilon(g)}$ with $\varepsilon\colon G\to\{\pm1\}$ and $\chi_\ell\colon G\to\Zl^\times$
characters, $R=\Zl$. Then $\rho$ is a homomorphism if and only if $\kappa$ is a $1$-cocycle with values in $\Zl(\chi_\ell)$, written $\Zl(1)$ when
$\chi_\ell$ is the cyclotomic character. On a subgroup $I$ on which $\chi_\ell$ and $\varepsilon$ are trivial, $\kappa|_I$ is a homomorphism $I\to\Zl$, and
the corner of $\rho|_I$ is $\kappa|_I$.
\end{enumerate}
\end{proposition}

\begin{proof}
(i) The upper right entry of $\rho(g)\rho(h)$ is $\eta_1(g)\kappa(h)+\kappa(g)\eta_2(h)$; the diagonal entries are always multiplicative. So $\rho$ is a
homomorphism if and only if $\kappa(gh)=\eta_1(g)\kappa(h)+\kappa(g)\eta_2(h)$. Dividing by $\eta_2(gh)=\eta_2(g)\eta_2(h)$ gives the cocycle identity for $b$.

(ii) With $P=\sm1r01$, the matrix of $g$ in the basis $e_1,e_2+re_1$ is
\[
P^{-1}\rho(g)P=\Mtwo{\eta_1(g)}{\kappa(g)+r(\eta_1(g)-\eta_2(g))}{0}{\eta_2(g)},
\]
so $b$ becomes $b+r(\eta_1\eta_2^{-1}-1)$. Scaling is clear.

(iii) A complement $L$ of $Re_1$ in $R^2$ maps isomorphically onto $R^2/Re_1\simeq R$, so it is free of rank one, spanned by the unique $v\in L$ mapping to
the class of $e_2$; and $v=e_2+re_1$ for some $r\in R$. It is $G$-stable if and only if the corner vanishes in the basis $e_1,e_2+re_1$, that is, $b+d^0r=0$ by
(ii). So $\rho$ is split if and only if $b$ is a coboundary.

(iv) Given a cocycle $b$, put $\kappa=b\eta_2$; by (i), $\rho$ is a homomorphism, and $E=R^2$ with this action is an extension of the stated kind with class $[b]$.
Conversely, an extension $0\to R(\eta_1)\xrightarrow{\iota}E\to R(\eta_2)\to0$ has the basis $e_1=\iota(1)$, $e_2$ a lift of $1\in R(\eta_2)$. In this basis
$g$ fixes the line $Re_1$, acting on it by $\eta_1(g)$, and acts on $E/Re_1$ by $\eta_2(g)$, so its matrix has the shape of (i). The only freedom is the lift $e_2$,
which can be replaced by $e_2+re_1$; by (ii) this changes $b$ by the coboundary $d^0r$, so $[b]$ depends only on $E$. If $E$, $E'$ with such bases have cocycles
$b$ and $b'=b+d^0r$, then in the basis $e_1,e_2+re_1$ of $E$ the cocycle is $b'$ by (ii), so the $R$-linear map $E'\to E$ with $e_1'\mapsto e_1$, $e_2'\mapsto e_2+re_1$
has the same matrices on both sides: it is $G$-equivariant, and it induces the identity on the sub and on the quotient. An isomorphism inducing the identity on
the sub and the quotient maps $e_1$ to $e'_1$ and $e_2$ to a lift of $1$, so it is of this form; hence the bijection.

(v) Here $\eta_1=\chi_\ell\varepsilon$, $\eta_2=\varepsilon$, so $\eta_1\eta_2^{-1}=\chi_\ell$ and $b=\varepsilon\kappa\cdot\varepsilon^{-1}=\kappa$. On $I$ the action on $\Zl(\chi_\ell)$ is
trivial, so the cocycle $\kappa|_I$ is a homomorphism by \Cref{prop:cohomology:h0-h1}(iii).
\end{proof}

\begin{paperlink}\label{pl:cohomology:corner}
The proof of Lemma lem:tamfactor in the revised paper writes the matrices of $G_{\Q_q}$ on $T_\ell E$ as
$\sm{\chi_\ell\varepsilon}{\varepsilon\kappa}{0}{\varepsilon}$ ``with $\varepsilon$ unramified of order at most $2$ and $\kappa$ the Kummer cocycle of $q_E$'', and adds that
``$\varepsilon\kappa$, not $\kappa$, satisfies the cocycle relation for the twisted matrices''. \Cref{prop:cohomology:corner}(i) and (v) make this precise: the
corner divided by the lower character, $\varepsilon\kappa/\varepsilon=\kappa$, is the cocycle, with values in $\Zl(1)$, and its restriction to inertia is a
homomorphism. (With the corner written $\kappa$, as in the original version, the cocycle would be $\kappa/\varepsilon=\varepsilon\kappa$.) The uniqueness of the stable line in a non-split
extension, used in Lemma lem:old, is proved in \Cref{ch:reps}; (iii) above is its cohomological form. That $\kappa$ is a Kummer cocycle is
\Cref{ch:galcoh}'s.
\end{paperlink}

\begin{theorem}[$H^2$ and group extensions]\label{thm:cohomology:h2-ext}
Let $M$ be a $G$-module. The equivalence classes of group extensions $1\to M\to\mathcal E\xrightarrow{\pi}G\to1$ in which the conjugation action of $\mathcal E$
on the abelian normal subgroup $M$ induces the given action of $G$ are in bijection with $H^2(G,M)$. If $s\colon G\to\mathcal E$ is a set-theoretic section
with $s(1)=1$, the class is that of the $2$-cocycle $\psi(g,h)=s(g)s(h)s(gh)^{-1}\in M$ (written additively in $M$). The semidirect product $M\rtimes G$
corresponds to $0$.
\end{theorem}

\begin{proof}[Sketch of proof]
Since $\pi(s(g)s(h))=\pi(s(gh))$, $\psi(g,h)$ lies in $M$. Write $\mathcal E$ multiplicatively, so that $s(g)s(h)=\psi(g,h)s(gh)$ and $s(g)\,m\,s(g)^{-1}=gm$ for
$m\in M$. Expanding $s(g)s(h)s(k)$ in two ways,
\begin{align*}
\bigl(s(g)s(h)\bigr)s(k)&=\psi(g,h)\psi(gh,k)\,s(ghk),\\
s(g)\bigl(s(h)s(k)\bigr)&=s(g)\psi(h,k)s(g)^{-1}\cdot s(g)s(hk)=(g\psi(h,k))\,\psi(g,hk)\,s(ghk),
\end{align*}
and comparing, in additive notation in $M$, $g\psi(h,k)+\psi(g,hk)=\psi(g,h)+\psi(gh,k)$, which is $d^2\psi=0$. Another section is $s'(g)=\varphi(g)s(g)$ with
$\varphi\colon G\to M$, and
\[
s'(g)s'(h)s'(gh)^{-1}=\varphi(g)\,s(g)\varphi(h)s(g)^{-1}\cdot s(g)s(h)s(gh)^{-1}\cdot\varphi(gh)^{-1},
\]
which is $\varphi(g)+g\varphi(h)+\psi(g,h)-\varphi(gh)$ additively (the factors lie in the abelian group $M$): it changes $\psi$ by $d^1\varphi$. Conversely a cocycle $\psi$ (normalised so that $\psi(1,g)=\psi(g,1)=0$, which one can arrange by a
coboundary) defines a group law $(m,g)(m',g')=(m+gm'+\psi(g,g'),gg')$ on $M\times G$. Details are in \cite[Chapter IV, Section 3]{Brown}.
\end{proof}

\begin{proposition}[functoriality in $M$ and the long exact sequence]\label{prop:cohomology:les-groups}
Let $G$ be a group and $0\to M'\xrightarrow{i}M\xrightarrow{p}M''\to0$ an exact sequence of $G$-modules. Then
$0\to C^\bullet(G,M')\to C^\bullet(G,M)\to C^\bullet(G,M'')\to0$ is an exact sequence of complexes, and there is a long exact sequence
\begin{multline*}
0\to H^0(G,M')\to H^0(G,M)\to H^0(G,M'')\xrightarrow{\delta}H^1(G,M')\\
\to H^1(G,M)\to H^1(G,M'')\xrightarrow{\delta}H^2(G,M')\to\cdots,
\end{multline*}
natural with respect to morphisms of short exact sequences. If the modules are $\Zl[G]$-modules and the maps $\Zl$-linear, every map in it is $\Zl$-linear.
\end{proposition}

\begin{proof}
For each $n$, $C^n(G,-)=\operatorname{Map}(G^n,-)$ takes the exact sequence to an exact sequence: injectivity and exactness in the middle are pointwise, and a
map $G^n\to M''$ lifts pointwise because $p$ is surjective. The maps $i\circ-$, $p\circ-$ are morphisms of complexes (\Cref{lem:cohomology:dd}).
Now apply \Cref{thm:cohomology:les}. Naturality is \Cref{thm:cohomology:les}(iii). $\Zl$-linearity: for $\lambda\in\Zl$, multiplication by $\lambda$ on $M'$, $M$,
$M''$ is $G$-linear (the action is $\Zl$-linear) and commutes with $i$ and $p$, so it is an endomorphism of the short exact sequence. It induces multiplication
by $\lambda$ on each $C^n$, hence on each $H^n$, and naturality gives $\delta\circ\lambda=\lambda\circ\delta$. The maps $i_*$, $p_*$ are $\Zl$-linear for the same reason.
\end{proof}

\begin{lemma}[$H^0$ and torsion]\label{lem:cohomology:h0-torsion}
\begin{enumerate}[label=(\alph*)]
\item \sloppy For a $G$-module $W$ and an integer $m\ge1$, $H^0(G,W[m])=H^0(G,W)[m]$.
\item \sloppy Let $T$ be a $\Zl$-module without $\ell$-torsion on which $G$ acts $\Zl$-linearly, and $W=T\otimes_{\Zl}\Qlmod$. Then $W[\ell^k]\simeq T/\ell^kT$ as
$G$-modules, by $t\otimes\ell^{-k}\mapsto t\bmod\ell^k$. Hence $H^0(G,W)[\ell^k]\simeq H^0(G,T/\ell^kT)$, and $H^0(G,W)=\varinjlim_kH^0(G,T/\ell^kT)$ along the maps
induced by multiplication by $\ell$, $T/\ell^kT\to T/\ell^{k+1}T$.
\item A finite $\ell$-primary abelian group $A$ with $\#A[\ell]\le\ell$ is cyclic, $A\simeq\Z/\ell^h$, where $h=\len_{\Zl}A$ is the largest $k$ such that $A$ has an
element of order $\ell^k$.
\item Suppose in (b) that $H^0(G,W)$ is finite and $\#H^0(G,T/\ell T)\le\ell$. Then $H^0(G,W)$ is cyclic, and $\len H^0(G,W)$ is the largest $k$ for which
$H^0(G,T/\ell^kT)$ contains an element of exact order $\ell^k$, that is, an invariant element of $T/\ell^kT$ not in $\ell T/\ell^kT$.
\item If $T$ in (b) is free of rank $r$ and $\#H^0(G,T/\ell T)\le\ell$, then $H^0(G,W)$ is either finite cyclic or isomorphic to $\Qlmod$.
\end{enumerate}
\end{lemma}

\begin{proof}
(a) Both sides are the set of $w\in W$ with $mw=0$ and $gw=w$ for all $g$.

(b) Since $T$ has no $\ell$-torsion, $T\to T\otimes\Ql=T[1/\ell]$ is injective and $W=T[1/\ell]/T$. Then $W[\ell^k]=\ell^{-k}T/T$, and multiplication by $\ell^k$ is a
$G$-isomorphism $\ell^{-k}T/T\to T/\ell^kT$. Apply (a). As $W=\bigcup_kW[\ell^k]$, also $H^0(G,W)=\bigcup_kH^0(G,W[\ell^k])$; the inclusion
$W[\ell^k]\subseteq W[\ell^{k+1}]$ becomes multiplication by $\ell$ under the isomorphisms.

(c) By the structure theorem $A\simeq\bigoplus_{i=1}^s\Z/\ell^{a_i}$ with all $a_i\ge1$, and $A[\ell]\simeq(\Z/\ell)^s$, so $s\le1$. Then $A\simeq\Z/\ell^h$ with $h=\len A$,
and $\Z/\ell^h$ has an element of order $\ell^k$ exactly for $k\le h$.

(d) By (a), (b) the $\ell$-torsion of $H^0(G,W)$ is $H^0(G,T/\ell T)$, of order at most $\ell$, so (c) applies, and $H^0(G,W)$ has an element of order $\ell^k$ if
and only if its subgroup $H^0(G,W)[\ell^k]\simeq H^0(G,T/\ell^kT)$ does. An element $\bar x\in T/\ell^kT$ has order $\ell^k$ if and only if $\ell^{k-1}x\notin\ell^kT$, if and
only if $x\notin\ell T$, because $T$ has no $\ell$-torsion.

(e) $W\simeq(\Qlmod)^r$, and $X=H^0(G,W)$ is a submodule with $\#X[\ell]\le\ell$ (by (a), (b), $X[\ell]=H^0(G,T/\ell T)$). By Pontryagin duality (\Cref{ch:modules}),
dualising the inclusion $X\hookrightarrow W$ makes $X^*$ a quotient of $((\Qlmod)^r)^*=\Zl^r$, hence finitely generated. Dualising the exact sequence
$0\to X[\ell]\to X\xrightarrow{\ell}X$ gives the exact sequence $X^*\xrightarrow{\ell}X^*\to X[\ell]^*\to0$ (the dual of multiplication by $\ell$ is multiplication by $\ell$),
so $X^*/\ell X^*\simeq(X[\ell])^*$ has order at most $\ell$. By Nakayama's lemma $X^*$ is generated by one element: $X^*\simeq\Zl$ or $\Z/\ell^h$, and
$X\simeq X^{**}$ is $\Qlmod$ or $\Z/\ell^h$.
\end{proof}

Two examples show the two cases of (e). If $G$ acts trivially on $T=\Zl$, then every $T/\ell^kT$ is invariant and $H^0(G,W)=W=\Qlmod$: infinite. If $G$ acts on
$T=\Zl$ through a character $\chi$ with $\chi(G)\subseteq1+\ell^h\Zl$, $\chi(G)\not\subseteq1+\ell^{h+1}\Zl$, then the invariant elements of $T/\ell^kT$ are all of it for
$k\le h$ and lie in $\ell T/\ell^kT$ for $k>h$, so $H^0(G,W)\simeq\Z/\ell^h$ (\Cref{ex:cohomology:h0-length}). Part (d) is the tool for reading off $h$: find the
largest level $k$ at which a primitive vector of $T/\ell^kT$ is invariant.

\begin{paperlink}\label{pl:cohomology:h0}
\sloppy In the proof of Lemma lem:h0: ``The $\ell^k$-torsion of $H^0(\Q,W^D)$ is $H^0(\Q,\ad(T/\ell^k)(1))$'' is \Cref{lem:cohomology:h0-torsion}(b) for the lattice
$\ad T(1)$, since $W^D=\ad T(1)\otimes\Qlmod$. The next step, that $H^0(\Q,W^D)$ ``a submodule of $(\Q_\ell/\Z_\ell)^3$ with $\ell$-torsion of rank at most one, is
cyclic'', is (e): such a module is finite cyclic or isomorphic to $\Qlmod$, and the paper's determinant argument, which shows that every element has
order at most $\ell$, excludes $\Qlmod$. The identification of the length with ``the largest $k$ such that $\rho\bmod\ell^k\simeq(\rho\otimes\chi_\ell)\bmod\ell^k$''
is (d): an element of exact order $\ell^k$ is an invariant map $T/\ell^k\to(T/\ell^k)(1)$ not divisible by $\ell$. The lemma itself is proved in \Cref{ch:galrep}.
\end{paperlink}

\section{Changing the group}\label{sec:cohomology:change}

The proof of Lemma lem:tamfactor uses three operations that change the group: restriction from $G_{\Q_q}$ to the inertia group $I_q$, inflation from
the quotient $I_q/K$ to $I_q$, and the conjugation action of $G_{\Q_q}/I_q$ on $H^1(I_q,M)$. The paper's ``by inflation-restriction
$H^1(I,M')=H^1(\langle t\rangle,M')$'' is \Cref{thm:cohomology:inf-res} with \Cref{cor:cohomology:pro-l-quotient}. Corestriction enters only through
$\Cor\circ\Res=[G:H]$, which gives \Cref{thm:cohomology:order-kills}, used for the real place after Lemma lem:dfg21. All of them are instances of
one construction.

\begin{proposition}[compatible pairs]\label{prop:cohomology:compatible}
Let $f\colon G'\to G$ be a group homomorphism, $M$ a $G$-module, $M'$ a $G'$-module, and $\varphi\colon M\to M'$ a homomorphism of abelian groups with
$\varphi(f(g')m)=g'\varphi(m)$ for all $g'\in G'$, $m\in M$ (a \emph{compatible pair}). Then $c\mapsto\varphi\circ c\circ f^{\times n}$, that is
\[
c\longmapsto\bigl((g'_1,\dots,g'_n)\mapsto\varphi(c(f(g'_1),\dots,f(g'_n)))\bigr),
\]
is a morphism of complexes $C^\bullet(G,M)\to C^\bullet(G',M')$. It induces maps $(f,\varphi)^*\colon H^n(G,M)\to H^n(G',M')$, compatible with the connecting
maps of short exact sequences of compatible pairs, and $(f,\varphi)^*\circ(f_1,\varphi_1)^*=(f_1\circ f,\varphi\circ\varphi_1)^*$ for composable pairs.
\end{proposition}

\begin{proof}
Write $c'=\varphi\circ c\circ f^{\times n}$. In $(d c')(g'_1,\dots,g'_{n+1})$ the first term is $g'_1\varphi(c(f(g'_2),\dots))=\varphi\bigl(f(g'_1)c(f(g'_2),\dots)\bigr)$ by
compatibility; the middle terms are $\varphi(c(\dots,f(g'_i)f(g'_{i+1}),\dots))$ because $f$ is a homomorphism; the last is $\varphi(c(f(g'_1),\dots,f(g'_n)))$. The sum is
$\varphi\bigl((dc)(f(g'_1),\dots,f(g'_{n+1}))\bigr)$, so $dc'=(dc)'$. Compatibility with connecting maps follows from \Cref{thm:cohomology:les}(iii), applied to
the morphism of short exact sequences of complexes given by the three pairs. For the composition rule, let $(f_1\colon G\to G_1,\varphi_1\colon M_1\to M)$ and
$(f\colon G'\to G,\varphi\colon M\to M')$ be compatible. Then $(f_1\circ f,\varphi\circ\varphi_1)$ is compatible, since
$\varphi\varphi_1(f_1f(g')m)=\varphi(f(g')\varphi_1(m))=g'\varphi\varphi_1(m)$, and on a cochain $c$ both sides send $c$ to
$(g'_1,\dots)\mapsto\varphi(\varphi_1(c(f_1f(g'_1),\dots)))$.
\end{proof}

The definition of a compatible pair encodes the one requirement that makes ``compose with $f$, then apply $\varphi$'' commute with $d$: the first term of
$d$ involves the action, so $\varphi$ must convert the action of $f(g')$ on $M$ into the action of $g'$ on $M'$. The variance is mixed (the group map goes from
$G'$ to $G$, the module map from $M$ to $M'$), which is why restriction ($G'$ a subgroup) and inflation ($G$ a quotient, $M=M^H$ a submodule) both appear as
special cases below.

\begin{proposition}[restriction, inflation, corestriction]\label{prop:cohomology:res-inf-cor}
Let $H\le G$ and $M$ a $G$-module.
\begin{enumerate}[label=(\alph*)]
\item \emph{Restriction} $\Res\colon H^n(G,M)\to H^n(H,M)$ is $(H\hookrightarrow G,\id_M)^*$: restrict cochains to $H^n$.
\item If $H$ is normal, \emph{inflation} $\Inf\colon H^n(G/H,M^H)\to H^n(G,M)$ is $(G\to G/H,\ M^H\hookrightarrow M)^*$: compose with the projection.
\item If $[G:H]=r<\infty$, there are \emph{corestriction} maps $\Cor\colon H^n(H,M)\to H^n(G,M)$ for all $n$, compatible with the long exact sequences. In
degree $0$, $\Cor(m)=\sum_{i=1}^rg_im$ for left coset representatives $g_1,\dots,g_r$ ($G=\bigsqcup g_iH$), independent of the choice. In degree $1$, writing
$gg_i=g_{\tau(i)}h_i(g)$ with $\tau=\tau_g$ a permutation and $h_i(g)\in H$,
\[
(\Cor\varphi)(g)=\sum_{i=1}^rg_{\tau(i)}\,\varphi(h_i(g)).
\]
\item $\Cor\circ\Res=r$ on $H^n(G,M)$ for all $n$. If $H$ is normal, $\Res\circ\Cor=\sum_{\sigma\in G/H}\sigma_*$ on $H^n(H,M)$, where $\sigma_*$ is the
conjugation action of \Cref{prop:cohomology:conjugation}.
\end{enumerate}
\end{proposition}

\begin{proof}
(a), (b): the pairs are compatible: for (b), $m\in M^H$ and $\pi(g)m=gm$. Apply \Cref{prop:cohomology:compatible}.

(c) Degree $0$: for $m\in M^H$, $g_ihm=g_im$, so $\sum g_im$ does not depend on the representatives; it is $G$-invariant because $g$ permutes the cosets:
$g\sum_ig_im=\sum_ig_{\tau(i)}h_i(g)m=\sum_ig_{\tau(i)}m=\sum_jg_jm$.
Degree $1$: first $h_i(gg')=h_{\tau_{g'}(i)}(g)h_i(g')$ and $\tau_{gg'}=\tau_g\tau_{g'}$, since $gg'g_i=gg_{\tau_{g'}(i)}h_i(g')=g_{\tau_g\tau_{g'}(i)}h_{\tau_{g'}(i)}(g)h_i(g')$. For
$\varphi\in Z^1(H,M)$,
\[
(\Cor\varphi)(gg')=\sum_ig_{\tau_g\tau_{g'}(i)}\bigl[\varphi(h_{\tau_{g'}(i)}(g))+h_{\tau_{g'}(i)}(g)\,\varphi(h_i(g'))\bigr].
\]
Reindexing by $j=\tau_{g'}(i)$, the first part is $(\Cor\varphi)(g)$; in the second, $g_{\tau_g(j)}h_j(g)=gg_j$, so it is $g\sum_ig_{\tau_{g'}(i)}\varphi(h_i(g'))=g(\Cor\varphi)(g')$.
So $\Cor\varphi$ is a cocycle. If $\varphi(h)=hm-m$, then $(\Cor\varphi)(g)=\sum_i(g_{\tau(i)}h_i(g)m-g_{\tau(i)}m)=\sum_i(gg_im-g_{\tau(i)}m)=gm'-m'$ with $m'=\sum_ig_im$, a
coboundary. So $\Cor$ is defined on $H^1$. In higher degrees we do not prove the existence or the compatibility with long exact sequences; see
\cite[Chapter III, Section 9]{Brown} and \cite[Chapter I, Section 5]{NSW}.

(d) Degree $0$: $\Cor(\Res m)=\sum g_im=rm$ for $m\in M^G$. Degree $1$: for $c\in Z^1(G,M)$, $h_i(g)=g_{\tau(i)}^{-1}gg_i$, and the cocycle identity gives
$g_{\tau(i)}c(g_{\tau(i)}^{-1}gg_i)=-c(g_{\tau(i)})+c(g)+gc(g_i)$: indeed $c(x^{-1}y)=c(x^{-1})+x^{-1}c(y)=x^{-1}(c(y)-c(x))$ for $x=g_{\tau(i)}$ and $y=gg_i$, by
\Cref{prop:cohomology:h0-h1}(ii), and $c(gg_i)=c(g)+gc(g_i)$. Summing over $i$ (and using that $\tau$ is a permutation, so $\sum_ic(g_{\tau(i)})=\sum_ic(g_i)$),
$(\Cor\Res c)(g)=rc(g)+gm-m$ with $m=\sum_ic(g_i)$: so $\Cor\Res c=rc+d^0m$. If $H$ is normal and $h\in H$, then $hg_iH=g_i(g_i^{-1}hg_i)H=g_iH$, so $\tau_h=\id$ and
$h_i(h)=g_i^{-1}hg_i$; hence $(\Res\Cor\varphi)(h)=\sum_ig_i\varphi(g_i^{-1}hg_i)=\sum_i(g_i\cdot\varphi)(h)$. In degrees $n\ge2$ (d) is stated; see the references
for (c).
\end{proof}

The next theorem is the reason a group of order $2$ is invisible to $\ell$-adic coefficients when $\ell$ is odd. We prove it in every degree with an
explicit cochain, so that it does not depend on the corestriction in degrees $n\ge2$.

\begin{theorem}[the order kills the cohomology]\label{thm:cohomology:order-kills}
Let $G$ be a finite group of order $N$ and $M$ a $G$-module. Then $N\cdot H^n(G,M)=0$ for every $n\ge1$. More precisely, for $\varphi\in Z^n(G,M)$ put
$\psi(g_1,\dots,g_{n-1})=\sum_{g\in G}\varphi(g_1,\dots,g_{n-1},g)$; then $N\varphi=(-1)^nd^{n-1}\psi$. Consequently:
\begin{enumerate}[label=(\roman*)]
\item if multiplication by $N$ is bijective on $M$, then $H^n(G,M)=0$ for $n\ge1$;
\item if $M$ is finitely generated as an abelian group, $H^n(G,M)$ is finite for $n\ge1$;
\item if $\ell\nmid N$ and $M$ is an $\ell$-primary torsion group or a $\Zl$-module, then $H^n(G,M)=0$ for $n\ge1$, and $M^G$ is a direct summand of $M$ as a
$G$-module, cut out by the projector $N^{-1}N_G$.
\end{enumerate}
\end{theorem}

\begin{proof}
Write the cocycle identity $(d^n\varphi)(g_1,\dots,g_n,g)=0$ and sum over $g\in G$:
\begin{multline*}
g_1\psi(g_2,\dots,g_n)+\sum_{i=1}^{n-1}(-1)^i\psi(\dots,g_ig_{i+1},\dots)\\
+(-1)^n\sum_{g}\varphi(g_1,\dots,g_{n-1},g_ng)+(-1)^{n+1}N\varphi(g_1,\dots,g_n)=0.
\end{multline*}
Here the first term comes from $g_1\varphi(g_2,\dots,g_n,g)$; the middle sum from the products $g_ig_{i+1}$ with $i\le n-1$, in which the last variable $g$ is
untouched; the third from the product $g_ng$ in position $n$; and the last from dropping the last variable $g$, which leaves $\varphi(g_1,\dots,g_n)$ for each of the
$N$ values of $g$. As $g$ runs over $G$ so does $g_ng$, so the third sum is $(-1)^n\psi(g_1,\dots,g_{n-1})$. Comparing with the definition
\[
(d^{n-1}\psi)(g_1,\dots,g_n)=g_1\psi(g_2,\dots,g_n)+\sum_{i=1}^{n-1}(-1)^i\psi(\dots,g_ig_{i+1},\dots)+(-1)^n\psi(g_1,\dots,g_{n-1}),
\]
the first three terms add up to $(d^{n-1}\psi)(g_1,\dots,g_n)$, so $d^{n-1}\psi+(-1)^{n+1}N\varphi=0$, that is
$N\varphi=(-1)^nd^{n-1}\psi$. (For $n=1$: with $m=\sum_g\varphi(g)$, summing $\varphi(\sigma g)=\varphi(\sigma)+\sigma\varphi(g)$ gives $m=N\varphi(\sigma)+\sigma m$.)

(i) Multiplication by $N$ on $H^n(G,M)$ is induced by the automorphism $N$ of $M$ (functoriality), so it is bijective; it is also $0$. (ii) $C^n(G,M)=M^{G^n}$ is
finitely generated, so $H^n(G,M)$ is a finitely generated abelian group killed by $N$. (iii) An integer prime to $\ell$ acts bijectively on an $\ell$-primary
torsion group and on a $\Zl$-module (it is a unit of $\Zl$), so (i) applies. The map $e=N^{-1}N_G$ is $G$-linear, takes values in $M^G$ and is the identity
on $M^G$; so $M=M^G\oplus\ker e$ as $G$-modules.
\end{proof}

\begin{proposition}[the conjugation action]\label{prop:cohomology:conjugation}
Let $H\trianglelefteq G$ and $M$ a $G$-module. For $\sigma\in G$ the pair $(h\mapsto\sigma^{-1}h\sigma,\ m\mapsto\sigma m)$ is compatible and induces
$\sigma_*\colon H^n(H,M)\to H^n(H,M)$, given on cochains by
\[
(\sigma\cdot\varphi)(h_1,\dots,h_n)=\sigma\,\varphi(\sigma^{-1}h_1\sigma,\dots,\sigma^{-1}h_n\sigma).
\]
\begin{enumerate}[label=(\roman*)]
\item This is a left action of $G$ on $C^\bullet(H,M)$ by morphisms of complexes, hence on $H^n(H,M)$; it is $\Zl$-linear when $M$ is a $\Zl[G]$-module, so it
preserves $H^n(H,M)_{\mathrm{tor}}$, and $(H^n(H,M)_{\mathrm{tor}})^{G}$ is defined.
\item $H$ acts trivially on $H^n(H,M)$; so the action factors through $G/H$, and the map induced by a class of $G/H$ does not depend on its representative.
For $n=1$ and $h_0\in H$, $\varphi\in Z^1(H,M)$: $(h_0\cdot\varphi)-\varphi=d^0(\varphi(h_0))$.
\item The restriction $H^n(G,M)\to H^n(H,M)$ lands in $H^n(H,M)^{G/H}$. For $n=1$ and $c\in Z^1(G,M)$: $\sigma\cdot\Res c-\Res c=d^0(c(\sigma))$ on $H$.
\end{enumerate}
\end{proposition}

\begin{proof}
Compatibility: $\sigma\bigl((\sigma^{-1}h\sigma)m\bigr)=h\sigma m$. (i) The cochain formula is the one of \Cref{prop:cohomology:compatible}. For the action,
$(\sigma\cdot(\tau\cdot\varphi))(h)=\sigma\tau\,\varphi(\tau^{-1}\sigma^{-1}h\sigma\tau)=((\sigma\tau)\cdot\varphi)(h)$, and similarly in several variables. Since $\varphi\mapsto\sigma\cdot\varphi$ is
$\Zl$-linear, it maps torsion classes to torsion classes.

(iii) and (ii) for $n=1$. Let $c$ be a $1$-cocycle on a group containing $H$ and $\sigma$ (that is, on $G$ for (iii); on $H$, with $\sigma=h_0\in H$, for (ii)).
By the cocycle identity and $c(\sigma^{-1})=-\sigma^{-1}c(\sigma)$,
\[
c(\sigma^{-1}h\sigma)=c(\sigma^{-1})+\sigma^{-1}\bigl(c(h)+hc(\sigma)\bigr)=\sigma^{-1}\bigl(c(h)+hc(\sigma)-c(\sigma)\bigr),
\]
so $(\sigma\cdot c)(h)=c(h)+(h-1)c(\sigma)=c(h)+(d^0c(\sigma))(h)$ for $h\in H$. Here $c(\sigma)\in M$, so $d^0c(\sigma)$ is a coboundary on $H$. For $n=0$ both statements are
clear ($\sigma\cdot m=\sigma m$, and $m\in M^H$, resp. $M^G$). For $n\ge2$ the triviality of the action of $H$ and the invariance of the image of restriction
are stated; see \cite[Chapter III, Section 8]{Brown} and \cite[Chapter I, Section 5]{NSW}. (They follow from the case $n=0$ by dimension shifting, since
both sides are $\delta$-functors.) The independence of the representative follows from the triviality on $H$.
\end{proof}

\begin{theorem}[inflation-restriction]\label{thm:cohomology:inf-res}
Let $H\trianglelefteq G$ and $M$ a $G$-module.
\begin{enumerate}[label=(\alph*)]
\item The sequence $0\to H^1(G/H,M^H)\xrightarrow{\Inf}H^1(G,M)\xrightarrow{\Res}H^1(H,M)^{G/H}$ is exact.
\item \sloppy (Five-term exact sequence of Hochschild and Serre.) There is a homomorphism $\operatorname{tg}\colon H^1(H,M)^{G/H}\to H^2(G/H,M^H)$, the
\emph{transgression}, such that
\[
0\to H^1(G/H,M^H)\xrightarrow{\Inf}H^1(G,M)\xrightarrow{\Res}H^1(H,M)^{G/H}\xrightarrow{\operatorname{tg}}H^2(G/H,M^H)\xrightarrow{\Inf}H^2(G,M)
\]
is exact.
\item If $H^i(H,M)=0$ for $1\le i\le n-1$, then $0\to H^n(G/H,M^H)\xrightarrow{\Inf}H^n(G,M)\xrightarrow{\Res}H^n(H,M)$ is exact.
\item If $H^1(H,M)=0$, then $\Inf\colon H^1(G/H,M^H)\to H^1(G,M)$ is an isomorphism. In particular, if $H$ acts trivially on $M$ and $\Hom(H,M)=0$, then
$\Inf\colon H^1(G/H,M)\to H^1(G,M)$ is an isomorphism.
\end{enumerate}
These hold for abstract groups; and, with the conventions of \Cref{sec:cohomology:profinite}, for profinite $G$, closed normal $H$ and continuous
cochains with coefficients of type (D), (C) or (V) (in (b) and (c) for type (D); see \Cref{rem:cohomology:inf-res-continuous}). For $G/H\simeq\widehat\Z$, (b) holds
also for types (C) and (V), with $\operatorname{tg}=0$ (\Cref{prop:cohomology:inf-res-zhat}).
\end{theorem}

\begin{proof}
(a) Write $\pi\colon G\to G/H$. By \Cref{prop:cohomology:conjugation}(iii) the image of $\Res$ is $G/H$-invariant.

$\Inf$ is injective. Let $c\in Z^1(G/H,M^H)$ with $c\circ\pi=d^0m$ for some $m\in M$, that is $c(\pi(g))=gm-m$. For $h\in H$, $0=c(1)=hm-m$, so $m\in M^H$, and then
$c=d^0m$ in $C^1(G/H,M^H)$.

$\Res\circ\Inf=0$: $(c\circ\pi)(h)=c(1)=0$.

$\ker\Res\subseteq\im\Inf$. Let $c\in Z^1(G,M)$ with $c(h)=hm-m$ for $h\in H$. Then $c'=c-d^0m$ is cohomologous to $c$ and vanishes on $H$. For $g\in G$,
$h\in H$: $c'(gh)=c'(g)+gc'(h)=c'(g)$, so $c'$ is constant on the cosets $gH$; and $c'(hg)=c'(h)+hc'(g)=hc'(g)$, while $hg=g(g^{-1}hg)$ gives
$c'(hg)=c'(g)$. So $hc'(g)=c'(g)$: $c'$ takes values in $M^H$, and $c'=\bar c\circ\pi$ with $\bar c\colon G/H\to M^H$, which is a cocycle because $\pi$ is surjective:
for $g,g'\in G$, $\bar c(\pi(g)\pi(g'))=c'(gg')=c'(g)+gc'(g')=\bar c(\pi(g))+\pi(g)\bar c(\pi(g'))$, the action of $G/H$ on $M^H$ being that of any lift.

In words: a class on $G$ that dies on $H$ can be normalised to vanish on $H$, and a cocycle vanishing on a normal subgroup is automatically constant on its cosets
with values in $M^H$, so it comes from the quotient. The sequence does not continue with a surjection in general; the obstruction to extending an invariant class
from $H$ to $G$ is the transgression of (b). When $G/H\simeq\widehat\Z$ there is no obstruction (\Cref{prop:cohomology:inf-res-zhat}).

(d) follows from (a), since $H^1(H,M)=0$ makes the invariants $0$; with trivial action $M^H=M$ and $H^1(H,M)=\Hom(H,M)$ (\Cref{prop:cohomology:h0-h1}(iii)).

(b) and (c) are stated; see \cite[Chapter VII, Section 6]{SerreLocalFields}, \cite[Chapter I, Section 6 and Chapter II, Section 4]{NSW} and \cite[Section 6.8]{Weibel}.
The transgression is constructed on cocycles: a class in $H^1(H,M)^{G/H}$ is represented by a cocycle $c$ on $H$ that extends to a cochain $\tilde c$ on $G$
with $\tilde c(gh)=\tilde c(g)+g c(h)$ and $\tilde c(hg)=c(h)+h\tilde c(g)$; then $d^1\tilde c$ is a $2$-cocycle on $G$ that is inflated from a $2$-cocycle on $G/H$ with values in $M^H$,
and $\operatorname{tg}[c]$ is its class (\Cref{ex:cohomology:transgression}).
\end{proof}

\begin{remark}[continuous cochains]\label{rem:cohomology:inf-res-continuous}
The proof of (a) uses only cochain manipulations, so it holds verbatim for continuous cochains (\Cref{def:cohomology:continuous}) once one checks that
$\bar c$ is continuous when $c'$ is. This holds because the projection $\pi\colon G\to G/H$ of topological groups is open: for an open set $U\subseteq M^H$,
$\bar c^{-1}(U)=\pi(c'^{-1}(U))$ is open. Also $M^H$ is closed in $M$ (an intersection of closed sets $\{m:hm=m\}$, $M$ being Hausdorff), so for $M$ of type (C)
or (V) it is again of type (C) or (V). This is the form in which \Cref{ch:galcoh} uses (a): $\Hur(\Q_p,M)=H^1(G_{\Q_p}/I_p,M^{I_p})$ is the kernel of the restriction
to $I_p$ (\Cref{conv:cohomology:galois}).
\end{remark}

\begin{theorem}[the Hochschild-Serre spectral sequence]\label{thm:cohomology:hochschild-serre}
For $H\trianglelefteq G$ and a $G$-module $M$ there is a first-quadrant spectral sequence
\[
E_2^{p,q}=H^p(G/H,H^q(H,M))\Longrightarrow H^{p+q}(G,M),
\]
with $G/H$ acting on $H^q(H,M)$ by conjugation; its edge maps are inflation and restriction, and its five-term exact sequence is
\Cref{thm:cohomology:inf-res}(b). The same holds for profinite $G$, closed normal $H$ and $M$ of type (D). Consequently, if $H^q(H,M)=0$ for all $q\ge1$,
then $\Inf\colon H^p(G/H,M^H)\to H^p(G,M)$ is an isomorphism for every $p$.
\end{theorem}

\begin{proof}
We do not prove this; see \cite[Chapter VII, Section 6]{Brown}, \cite[Section 6.8]{Weibel} and \cite[Chapter II, Section 4]{NSW}. The consequence holds because
the spectral sequence then has a single nonzero row $q=0$.
\end{proof}

\section{Finite groups}\label{sec:cohomology:finite}

Two facts about finite groups reach the paper. The first is the formula for $H^1$ of a finite cyclic group, which is the model for the procyclic
formula of \Cref{sec:cohomology:procyclic} and also its main counterexample. The second is that the cohomology of $\Gal(\C/\R)$, a group of order $2$, is
killed by $2$: the paper's ``At $\infty$, $H^1(\R,T^D)$ is killed by $2$, so $H^1_f(\R,T^D)=0$ for $\ell$ odd'' (the paragraph after Lemma lem:dfg21). Tate
cohomology is recalled briefly, for orientation.

\begin{theorem}[cohomology of a finite cyclic group]\label{thm:cohomology:cyclic}
Let $G=\langle t\rangle$ be cyclic of order $m$, $N=1+t+\dots+t^{m-1}$, and $M$ a $G$-module.
\begin{enumerate}[label=(\roman*)]
\item The map $\varphi\mapsto\varphi(t)$ induces an isomorphism $H^1(G,M)\simeq\ker(N\colon M\to M)/(t-1)M$. A cocycle is determined by $a=\varphi(t)$ through
$\varphi(t^k)=(1+t+\dots+t^{k-1})a$ for $0\le k<m$, and $a$ occurs if and only if $Na=0$.
\item $H^2(G,M)\simeq M^G/NM$, and $H^{n+2}(G,M)\simeq H^n(G,M)$ for all $n\ge1$.
\item For the trivial action on $\Z$: $H^1(G,\Z)=0$, $H^2(G,\Z)=\Z/m$. For the trivial action on $\Z/m$: $H^1=H^2=\Z/m$.
\end{enumerate}
\end{theorem}

\begin{proof}
(i) Put $S_k=1+t+\dots+t^{k-1}$ ($S_0=0$). By induction from $\varphi(t^{k+1})=\varphi(t^k)+t^k\varphi(t)$, a cocycle satisfies $\varphi(t^k)=S_ka$ for all $k\ge0$; at
$k=m$, $0=\varphi(1)=S_ma=Na$. Conversely, given $a$ with $Na=0$, define $\varphi(t^k)=S_ka$ for $0\le k<m$. For $0\le j,k<m$ we have $S_{j+k}=S_j+t^jS_k$, and if
$j+k\ge m$ then $S_{j+k}a=S_ma+t^mS_{j+k-m}a=S_{j+k-m}a$; so $\varphi(t^jt^k)=\varphi(t^j)+t^j\varphi(t^k)$ in all cases. Finally $(d^0b)(t)=(t-1)b$, so the coboundaries
correspond to $(t-1)M$.

(ii) is stated; see \cite[Chapter VIII, Section 4]{SerreLocalFields} and \cite[Chapter IV, Section 8]{CasselsFrohlich}. It comes from the periodic free resolution
$\cdots\xrightarrow{N}\Z[G]\xrightarrow{t-1}\Z[G]\xrightarrow{N}\Z[G]\xrightarrow{t-1}\Z[G]\to\Z\to0$ and \Cref{thm:cohomology:as-ext}: applying $\Hom_{\Z[G]}(-,M)$ gives
$M\xrightarrow{t-1}M\xrightarrow{N}M\xrightarrow{t-1}M\xrightarrow{N}\cdots$.

(iii) For $\Z$ with trivial action $N=m$ is injective and $t-1=0$, so $H^1=0$ and $H^2=\Z/m\Z$. For $\Z/m$, $N=m=0$ and $t-1=0$.
\end{proof}

\begin{definition}[Tate cohomology]\label{def:cohomology:tate}
For a finite group $G$ and a $G$-module $M$ put $\hat H^n(G,M)=H^n(G,M)$ for $n\ge1$, $\hat H^0(G,M)=M^G/N_GM$, $\hat H^{-1}(G,M)=\ker(N_G)/I_GM$, and
$\hat H^{-n}(G,M)=H_{n-1}(G,M)$ (group homology) for $n\ge2$.
\end{definition}

\begin{proposition}[properties of Tate cohomology]\label{prop:cohomology:tate}
Let $G$ be finite.
\begin{enumerate}[label=(\roman*)]
\item A short exact sequence of $G$-modules gives a long exact sequence of the $\hat H^n$, infinite in both directions.
\item $\#G$ kills $\hat H^n(G,M)$ for every $n\in\Z$.
\item Induced modules $\Z[G]\otimes_\Z A$ have $\hat H^n=0$ for every $n$.
\item For $G$ cyclic, $\hat H^n(G,M)\simeq\hat H^{n+2}(G,M)$ for every $n$; so $\hat H^{-1}=\ker N/(t-1)M\simeq H^1$ and $\hat H^0=M^G/NM\simeq H^2$.
\end{enumerate}
\end{proposition}

\begin{proof}
We do not prove this; see \cite[Chapter VIII, Sections 1 to 4]{SerreLocalFields} and \cite[Chapter VI, Section 4]{Brown}. For $n\ge1$, (ii) is
\Cref{thm:cohomology:order-kills}, and for $n=0,-1$ it follows from $N_G\cdot m=\#G\,m$ for $m\in M^G$ and from $\#G\,m=N_Gm-\sum_g(g-1)m$.
\end{proof}

\begin{corollary}[groups of order $2$ and $\ell$-adic coefficients]\label{cor:cohomology:z2}
Let $G=\{1,\sigma\}$ have order $2$, for instance $\Gal(\C/\R)$ generated by complex conjugation.
\begin{enumerate}[label=(\roman*)]
\item For every $G$-module $M$ and $n\ge1$, $H^n(G,M)$ is killed by $2$.
\item Let $\ell$ be odd and $M$ a $\Zl$-module with $\Zl$-linear action (a module of type (D) that is $\ell$-primary, or of type (C) or (V)). Then $H^n(G,M)=0$
for all $n\ge1$, $M=M^+\oplus M^-$ with $M^\pm=\{m:\sigma m=\pm m\}$, and $H^0(G,M)=M^+$.
\end{enumerate}
In the notation $H^n(\R,M):=H^n(\Gal(\C/\R),M)$: $H^1(\R,M)$ is a group of exponent dividing $2$, and it is $0$ when $M$ is an $\ell$-adic module with $\ell$ odd.
\end{corollary}

\begin{proof}
(i) is \Cref{thm:cohomology:order-kills} with $N=2$. (ii) $2$ is a unit of $\Zl$, so \Cref{thm:cohomology:order-kills}(iii) applies; the projectors are
$\frac12(1\pm\sigma)$.
\end{proof}

\begin{example}\label{ex:cohomology:z2-examples}
Let $\sigma$ act on $\Z$ by $-1$ (the module $\Z^-$). Then $N=1+\sigma=0$, so by \Cref{thm:cohomology:cyclic}, $H^1(G,\Z^-)=\Z/(\sigma-1)\Z=\Z/2\Z\ne0$, and
$H^2(G,\Z^-)=(\Z^-)^G/N\Z^-=0$. On $\Zl^-$ with $\ell$ odd, $H^1(G,\Zl^-)=\Zl/2\Zl=0$, as \Cref{cor:cohomology:z2} predicts. On $\Z_2^-$ it is $\Z/2$: the
hypothesis $\ell$ odd cannot be dropped.
\end{example}

\begin{paperlink}\label{pl:cohomology:real-place}
After Lemma lem:dfg21: ``At $\infty$, $H^1(\R,T^D)$ is killed by $2$, so $H^1_f(\R,T^D)=0$ for $\ell$ odd.'' The first clause is
\Cref{cor:cohomology:z2}(i); since $H^1(\R,T^D)$ is a $\Zl$-module (\Cref{def:cohomology:gmodule}) and $2\in\Zl^\times$, it is $0$, and so is its submodule
$H^1_f(\R,T^D)$ (the local condition at $\infty$ is defined in \Cref{ch:blochkato}). This is also the argument of \Cref{ch:modules}. Consequently the term at
$\infty$ in the exact sequence of Lemma lem:dfg21 vanishes, and \Cref{lem:cohomology:lengths} below has no contribution from the real place.
\end{paperlink}

\section{Profinite groups and continuous cohomology}\label{sec:cohomology:profinite}

Every cohomology group in the paper is the continuous cohomology of a profinite group: $H^1(G,\Zl(1))$ and $H^1(I_q,\ad T)$ in Lemma lem:tamfactor,
$H^1(\Q_p,T^D)$ in Lemma lem:dfg21, $H^1_f(\Q,W)$ in Theorem thm:main. This section defines it for the three types of coefficients, proves that the long
exact sequences exist in the cases the book needs, and records what can go wrong. Profinite groups and the Krull topology are in \Cref{ch:groups,ch:galois}.

\begin{definition}[discrete modules]\label{def:cohomology:discrete}
Let $G$ be a profinite group; its open normal subgroups form a basis of neighbourhoods of $1$. A \emph{discrete $G$-module} is a $G$-module $M$ such that the
action $G\times M\to M$ is continuous for the discrete topology on $M$. Modules of \emph{type (D)} are the discrete modules that are finite, or $\ell$-primary
torsion, such as $W=T\otimes\Qlmod$.
\end{definition}

\begin{lemma}\label{lem:cohomology:discrete}
For a $G$-module $M$ with the discrete topology the following are equivalent: (i) the action is continuous; (ii) every stabiliser
$G_m=\{g:gm=m\}$ is open; (iii) $M=\bigcup_UM^U$, the union over the open normal subgroups $U$.
\end{lemma}

\begin{proof}
(i)$\Rightarrow$(ii): $G_m$ is the preimage of $\{m\}$ under the continuous map $g\mapsto gm$. (ii)$\Rightarrow$(i): for $(g,m)$, the open set $gG_m\times\{m\}$
is mapped to $gm$. (ii)$\Leftrightarrow$(iii): an open subgroup has finite index, so it contains the intersection of its finitely many conjugates, an open
normal subgroup; conversely $M^U\ni m$ means $U\subseteq G_m$.
\end{proof}

\begin{definition}[continuous cohomology]\label{def:cohomology:continuous}
Let $G$ be profinite and $M$ a Hausdorff topological abelian group with a continuous action $G\times M\to M$ by automorphisms. Let
$C^n_{\mathrm{cont}}(G,M)\subseteq C^n(G,M)$ be the group of continuous maps $G^n\to M$. The differential of \Cref{def:cohomology:cochains} maps continuous cochains to
continuous cochains, because the action and the group law of $M$ are continuous. The \emph{continuous cohomology} is
$H^n(G,M)=H^n(C^\bullet_{\mathrm{cont}}(G,M))$. The coefficient modules of the book are:
\begin{enumerate}
\item[(D)] discrete modules of type (D);
\item[(C)] finitely generated $\Zl$-modules $M$ with the $\ell$-adic topology (the topology with basis of neighbourhoods of $0$ the $\ell^kM$) and a continuous
$\Zl$-linear action;
\item[(V)] finite-dimensional $\Ql$-vector spaces with their $\ell$-adic topology and a continuous linear action.
\end{enumerate}
\emph{Convention: all cohomology of profinite groups in this book is continuous cochain cohomology.} For a finite group with the discrete topology every
cochain is continuous, and the definition agrees with \Cref{def:cohomology:cochains}.
\end{definition}

Why continuity? The Galois groups of the paper are profinite and its representations are continuous, and the cohomology must see exactly the continuous
objects. The corner of a continuous triangular representation $\rho\colon G\to\GL_2(\Zl)$ is a continuous cocycle, and conversely a continuous cocycle gives a
continuous $\rho$ by \Cref{prop:cohomology:corner}; so continuous $H^1$ classifies continuous extensions, which is what the paper's Kummer cocycle $\kappa$ is. Abstract
cochains would also count discontinuous objects, which have no arithmetic meaning: for the profinite group $P=\prod_{n\ge1}\Z/2$ the continuous homomorphisms
$P\to\Z/2$ are those factoring through finitely many coordinates, a countable set, while the abstract ones form the dual of an $\F_2$-vector space of
uncountable dimension, a much larger set. The three types match the three modules that occur together in the paper: a lattice $T$ (type (C)), the vector space $V=T\otimes\Ql$ (type (V)) and the
discrete module $W=V/T$ (type (D)).

\begin{proposition}[basic properties]\label{prop:cohomology:continuous}
Let $G$ be profinite.
\begin{enumerate}[label=(\roman*)]
\item $H^0(G,M)=M^G$; $H^1(G,M)$ is the group of continuous crossed homomorphisms modulo the principal ones; if $G$ acts trivially, $H^1(G,M)=\Hom_{\mathrm{cont}}(G,M)$.
\item (Type (D).) For a discrete $G$-module $M$,
\[
H^n(G,M)=\varinjlim_UH^n(G/U,M^U),
\]
the limit over the open normal subgroups $U$, with the inflation maps $H^n(G/U,M^U)\to H^n(G/U',M^{U'})$ for $U'\subseteq U$ as transition maps.
\item (Type (C).) For $M$ finitely generated over $\Zl$, a $\Zl$-linear action is continuous if and only if each kernel $U_k$ of $G\to\Aut(M/\ell^kM)$ is open,
if and only if $G\to\Aut_{\Zl}(M)=\varprojlim_k\Aut(M/\ell^kM)$ is continuous.
\item (Type (V).) A continuous action of $G$ on $V$ of type (V) preserves a $\Zl$-lattice $T$ (a finitely generated $\Zl$-submodule spanning $V$), which is then of
type (C), with $V=T\otimes\Ql$.
\end{enumerate}
\end{proposition}

\begin{proof}
(i) As in \Cref{prop:cohomology:h0-h1}.

(ii) Let $\varphi\colon G^n\to M$ be continuous. Its image is compact in the discrete space $M$, hence finite, and its fibres are open and closed. A clopen subset
of the profinite space $G^n$ is a finite union of cosets of a subgroup $U^n$ with $U$ open normal (it is open, so a union of such cosets, and compact, so
finitely many suffice; then take the intersection of the finitely many $U$). So $\varphi$ factors through $(G/U)^n$ for some $U$. Each of the finitely many values
is fixed by an open normal subgroup (\Cref{lem:cohomology:discrete}); shrinking $U$, $\varphi$ takes values in $M^U$. Hence
$C^\bullet_{\mathrm{cont}}(G,M)=\bigcup_UC^\bullet(G/U,M^U)$, the union of subcomplexes along the inflation maps, which are injective on cochains. Cohomology
commutes with filtered unions (direct limits are exact), which gives (ii).

(iii) If the action is continuous, the stabiliser of each of the finitely many elements of $M/\ell^kM$ is open, so $U_k$ is open. Conversely, given $(g_0,m_0)$
and $k$, for $g=g_0u$ with $u\in U_k$ and $m\in m_0+\ell^kM$ we have $gm-g_0m_0=g(m-m_0)+g_0(um_0-m_0)\in\ell^kM$. The last equivalence is the definition of the
inverse-limit topology.

(iv) Choose a basis, so that $G\to\GL_n(\Ql)$ is continuous. $\GL_n(\Zl)$ is open in $\GL_n(\Ql)$, so its preimage $H$ is an open subgroup, of finite index.
The standard lattice $T_0=\Zl^n$ is $H$-stable, so $gT_0$ depends only on the coset $gH$, and $T=\sum_{g\in G/H}gT_0$ is a finite sum of lattices, hence a
lattice; it is $G$-stable.
\end{proof}

\begin{proposition}[long exact sequences for continuous cohomology]\label{prop:cohomology:les-cont}
Let $G$ be profinite and $0\to M'\xrightarrow{i}M\xrightarrow{p}M''\to0$ an exact sequence of topological $G$-modules with continuous $G$-maps, $i$ a
homeomorphism onto its closed image, and $p$ admitting a \emph{continuous section}: a continuous map $s\colon M''\to M$ with $ps=\id$, not necessarily additive.
\begin{enumerate}[label=(\roman*)]
\item $0\to C^\bullet_{\mathrm{cont}}(G,M')\to C^\bullet_{\mathrm{cont}}(G,M)\to C^\bullet_{\mathrm{cont}}(G,M'')\to0$ is exact, and there is a natural long exact sequence
\[
\cdots\to H^n(G,M')\to H^n(G,M)\to H^n(G,M'')\xrightarrow{\delta}H^{n+1}(G,M')\to\cdots.
\]
\item A continuous section exists in each of the following cases, which are all the cases used in this book:
\begin{enumerate}[label=(\alph*)]
\item $M''$ is discrete;
\item $p$ is a surjection of finitely generated $\Zl$-modules with their $\ell$-adic topologies;
\item $0\to T\to V\to W\to0$, with $T$ of type (C) without $\ell$-torsion, $V=T\otimes\Ql$ and $W=V/T=T\otimes\Qlmod$;
\item $0\to T\xrightarrow{\ell^k}T\to T/\ell^kT\to0$ with $T$ of type (C) without $\ell$-torsion.
\end{enumerate}
\item \sloppy In case (c), $\ker\bigl(H^1(G,T)\to H^1(G,V)\bigr)=H^1(G,T)_{\mathrm{tor}}$, so the connecting map $\delta\colon H^0(G,W)\to H^1(G,T)$ has image $H^1(G,T)_{\mathrm{tor}}$
and kernel the image of $H^0(G,V)$.
\end{enumerate}
\end{proposition}

\begin{proof}
(i) Injectivity is clear. If $\varphi\colon G^n\to M$ is continuous and $p\varphi=0$, then $\varphi$ takes values in $i(M')$, and $i^{-1}\circ\varphi$ is continuous because $i$ is
a homeomorphism onto its image. If $\psi\colon G^n\to M''$ is continuous, $s\circ\psi$ is a continuous lift. Now apply \Cref{thm:cohomology:les}.

(ii)(a) Every map out of a discrete space is continuous. (b) Write $M''=\bigoplus_j\Zl y_j$ with each summand free or $\simeq\Z/\ell^{a_j}$, choose $x_j\in M$ with
$p(x_j)=y_j$, and put $s(\sum_j\lambda_jy_j)=\sum_j\tilde\lambda_jx_j$, where $\tilde\lambda_j=\lambda_j$ on a free summand and $\tilde\lambda_j\in\{0,1,\dots,\ell^{a_j}-1\}\subset\Zl$ is the
representative on a summand $\Z/\ell^{a_j}$. The coordinates are continuous for the $\ell$-adic topology, and a map out of the finite discrete set
$\Z/\ell^{a_j}$ is continuous. The inclusion $i$ is an injective continuous map from a compact space to a Hausdorff one, hence a homeomorphism onto its
closed image. (c) $T$ is open in $V$, so $W=V/T$ is discrete and (a) applies; $T$ is compact, so its image is closed. (d) $T/\ell^kT$ is finite and
discrete, and $\ell^k\colon T\to\ell^kT$ is a continuous bijection of compact Hausdorff spaces.

(iii) A class in $H^1(G,V)$ is $0$ if and only if it is $d^0v$ for some $v\in V$. If $c\in Z^1(G,T)$ and $c=d^0v$, choose $a$ with $\ell^av\in T$; then
$\ell^ac=d^0(\ell^av)$ is a coboundary in $C^\bullet_{\mathrm{cont}}(G,T)$, so $[c]$ is torsion. Conversely, if $\ell^a[c]=0$, then $\ell^ac=d^0t$ with $t\in T$, and $c=d^0(\ell^{-a}t)$ in
$V$. (Torsion in a $\Zl$-module means killed by a power of $\ell$, since other integers are units.) The rest is the exactness of
$H^0(G,V)\to H^0(G,W)\to H^1(G,T)\to H^1(G,V)$.
\end{proof}

\begin{theorem}[Tate: inverse limits]\label{thm:cohomology:tate-limits}
Let $G$ be profinite and $T=\varprojlim_kT_k$ an inverse limit of finite discrete $G$-modules with surjective transition maps, with the inverse-limit
topology (for example $T$ of type (C) and $T_k=T/\ell^kT$).
\begin{enumerate}[label=(\roman*)]
\item For every $n\ge0$ there is a natural exact sequence $0\to\varprojlim_k^1H^{n-1}(G,T_k)\to H^n(G,T)\to\varprojlim_kH^n(G,T_k)\to0$.
\item In particular $H^1(G,T)=\varprojlim_kH^1(G,T_k)$ always (the $H^0(G,T_k)$ are finite, so the Mittag-Leffler condition holds and $\varprojlim^1=0$), and
$H^n(G,T)=\varprojlim_kH^n(G,T_k)$ whenever all the $H^{n-1}(G,T_k)$ are finite.
\item For $T$ of type (C) and $V=T\otimes\Ql$: $H^n(G,V)=H^n(G,T)\otimes_{\Zl}\Ql$ for all $n$.
\end{enumerate}
\end{theorem}

\begin{proof}
(i) is stated; see \cite[Section 2]{Tate95}, \cite[Chapter II, Section 7]{NSW} and \cite[Appendix B]{Rubin}. The idea: a continuous map into $\varprojlim T_k$ is a
compatible family of continuous maps into the $T_k$, so $C^\bullet_{\mathrm{cont}}(G,T)=\varprojlim_kC^\bullet(G,T_k)$, and the transition maps of these complexes are
surjective (lift values through the finite discrete $T_{k+1}\to T_k$ by a section). For an inverse system of complexes with surjective transition maps,
the cohomology of the limit sits in the stated $\varprojlim^1$ sequence \cite[Theorem 3.5.8]{Weibel}.

(ii) We prove the case $n=1$ directly. \emph{Surjectivity.} Given compatible classes $[c_k]\in H^1(G,T_k)$, choose the cocycles inductively: if $c_{k+1}$ reduces
to $c_k+d^0m$ with $m\in T_k$, lift $m$ to $\tilde m\in T_{k+1}$ and replace $c_{k+1}$ by $c_{k+1}-d^0\tilde m$. The compatible family $(c_k)$ is a continuous cocycle
$c\colon G\to T$. \emph{Injectivity.} Let $c\in Z^1_{\mathrm{cont}}(G,T)$ have every reduction $c_k$ a coboundary. The set $X_k=\{m\in T_k:c_k=d^0m\}$ is nonempty and finite (a
coset of $T_k^G$), and reduction maps $X_{k+1}$ into $X_k$. An inverse limit of nonempty finite sets is nonempty (\Cref{ch:galois}), so there is $m\in T$ with
$c=d^0m$.

(iii) A continuous cochain $G^n\to V$ has compact image, covered by the increasing open sets $\ell^{-a}T$, hence contained in one of them. So
$C^\bullet_{\mathrm{cont}}(G,V)=\bigcup_a\ell^{-a}C^\bullet_{\mathrm{cont}}(G,T)=C^\bullet_{\mathrm{cont}}(G,T)\otimes_{\Zl}\Ql$, and cohomology commutes with the exact functor $-\otimes_{\Zl}\Ql$.
\end{proof}

\begin{warning}[pitfalls of continuous cohomology]\label{warn:cohomology:pitfalls}
\begin{enumerate}[label=(\roman*)]
\item For $M$ of type (C), $H^n(G,M)$ is \emph{not} $\varinjlim_UH^n(G/U,M^U)$. Example: $H^1(\Zl,\Zl)$ with trivial action is $\Hom_{\mathrm{cont}}(\Zl,\Zl)=\Zl$ (the
homomorphisms $a\mapsto\lambda a$), while every $H^1(\Zl/\ell^k\Zl,\Zl)=\Hom(\Z/\ell^k,\Zl)$ is $0$, so the direct limit is $0$. Equivalently: a continuous
homomorphism from a compact group to $\Zl$ with the discrete topology has finite, hence zero, image.
\item The long exact sequence needs the continuous section of \Cref{prop:cohomology:les-cont}; it exists in all the cases we use.
\item The $\varprojlim^1$ terms of \Cref{thm:cohomology:tate-limits} can be nonzero in general; they vanish under the finiteness assumptions stated there.
\item $H^n(G,V)$ for $V$ of type (V) is a $\Ql$-vector space and sees no torsion; the torsion of $H^1(G,T)$ is the image of $H^0(G,W)$
(\Cref{prop:cohomology:les-cont}(iii)).
\item For a finite group with the discrete topology continuous and abstract cohomology coincide.
\item The formula $H^1(\Gamma,M)=M/(F-1)M$ for a procyclic group has precise hypotheses (\Cref{warn:cohomology:procyclic}): it fails for finite cyclic groups
(where it is $\ker N/(F-1)M$) and for discrete modules that are not torsion: $H^1(\widehat\Z,\Z)=\Hom_{\mathrm{cont}}(\widehat\Z,\Z)=0$ by the argument in (i), while
$\Z/(F-1)\Z=\Z$.
\end{enumerate}
\end{warning}

\begin{remark}[topological Nakayama, recalled]\label{rem:cohomology:top-nakayama}
We use the following form of Nakayama's lemma, proved in \Cref{ch:modules}: \emph{let $X$ be a compact (Hausdorff) $\Zl$-module, for instance a compact pro-$\ell$ abelian
group; if $X/\ell X$ is finite, then $X$ is a finitely generated $\Zl$-module, and any elements whose images generate $X/\ell X$ generate $X$.} It applies to
$H^n(G,T)$ for $T$ of type (C) whenever this group is an inverse limit of finite groups (\Cref{thm:cohomology:tate-limits}), and is used in
\Cref{prop:cohomology:h1-fg} and in \Cref{ch:galcoh}.
\end{remark}

\begin{proposition}[finiteness of $H^1$ for type (C)]\label{prop:cohomology:h1-fg}
Let $G$ be profinite and $T$ of type (C).
\begin{enumerate}[label=(\roman*)]
\item $H^1(G,T)=\varprojlim_kH^1(G,T/\ell^kT)$. If the groups $H^1(G,T/\ell^kT)$ are finite, $H^1(G,T)$ is a compact $\Zl$-module, an inverse limit of finite
$\ell$-groups.
\item Suppose $H^1(G,A)$ is finite for every finite $G$-module $A$ of $\ell$-power order. Then $H^1(G,T)$ is a finitely generated $\Zl$-module,
$H^1(G,V)=H^1(G,T)\otimes\Ql$ is finite-dimensional for $V=T\otimes\Ql$, and $H^1(G,T)_{\mathrm{tor}}$ is finite.
\item The hypothesis of (ii) holds for finite groups, for procyclic groups admitting a continuous surjection onto $\Zl$ (\Cref{thm:cohomology:procyclic-compact}
gives $H^1(\Gamma,T)=T/(F-1)T$ directly), for the inertia group $I_q$ at a prime $q\ne\ell$ (\Cref{prop:cohomology:tame-inertia}(iv)), and for $G_{\Q_p}$ and $G_{\Q,S}$
(\Cref{ch:galcoh}).
\end{enumerate}
\end{proposition}

\begin{proof}
(i) The equality is \Cref{thm:cohomology:tate-limits}(ii), proved there. An inverse limit of finite groups is compact.

(ii) First let $T$ have no $\ell$-torsion. The sequence $0\to T\xrightarrow{\ell}T\to T/\ell T\to0$ has a continuous section
(\Cref{prop:cohomology:les-cont}(d)), and its long exact sequence embeds $H^1(G,T)/\ell H^1(G,T)$ into the finite group $H^1(G,T/\ell T)$. By (i) $H^1(G,T)$ is
compact, and a compact $\Zl$-module $A$ with $A/\ell A$ finite is finitely generated (\Cref{rem:cohomology:top-nakayama}). In general $T_{\mathrm{tor}}$ is
finite and $T/T_{\mathrm{tor}}$ has no $\ell$-torsion; the sequence $0\to T_{\mathrm{tor}}\to T\to T/T_{\mathrm{tor}}\to0$ has a continuous section
(\Cref{prop:cohomology:les-cont}(b)), and its long exact sequence contains $H^1(G,T_{\mathrm{tor}})\xrightarrow{\alpha}H^1(G,T)\xrightarrow{\beta}H^1(G,T/T_{\mathrm{tor}})$, exact
at $H^1(G,T)$. The image of $\alpha$ is a quotient of the finite group $H^1(G,T_{\mathrm{tor}})$ (finite by hypothesis), and $H^1(G,T)/\im\alpha\simeq\im\beta$ is a submodule of
the finitely generated module $H^1(G,T/T_{\mathrm{tor}})$ (by the first case), hence finitely generated since $\Zl$ is noetherian. An extension of a finitely
generated module by a finite one is finitely generated (lift generators of the quotient and add the finitely many elements of the sub), so $H^1(G,T)$ is
finitely generated. The statement about $V$ is \Cref{thm:cohomology:tate-limits}(iii), and the torsion of a finitely generated $\Zl$-module is finite.

(iii) For a finite group, $H^1(G,A)$ is finite because the cochain groups are. For $I_q$, the proof of \Cref{prop:cohomology:tame-inertia}(iv) applies verbatim to a finite
$\ell$-primary $A$, with \Cref{thm:cohomology:prime-to-l}(a) and \Cref{thm:cohomology:procyclic-discrete} for $I/K\simeq\Zl$, and gives $H^1(I_q,A)\simeq A^K/(t-1)A^K$, a finite group.
The other cases are proved where cited.
\end{proof}

\section[Groups of order prime to l and pro-l quotients]{Groups of order prime to $\ell$ and pro-$\ell$ quotients}\label{sec:cohomology:prime-to-l}

This section justifies the sentence of the proof of Lemma lem:tamfactor ``the kernel of $I\to\langle t\rangle$ has trivial pro-$\ell$ quotient, so by
inflation-restriction $H^1(I,M')=H^1(\langle t\rangle,M')$''. Two properties of a closed subgroup $K$ are in play: ``order prime to $\ell$'', under which
all higher cohomology with $\ell$-adic coefficients vanishes, and the weaker ``trivial maximal pro-$\ell$ quotient'', under which $\Hom_{\mathrm{cont}}(K,M)=0$
for trivial $\ell$-adic $M$. The paper needs only the weaker one, together with \Cref{thm:cohomology:inf-res}(d). Both hold for the kernel of the tame
character.

\begin{definition}[supernatural orders and pro-$\ell$ groups]\label{def:cohomology:supernatural}
A \emph{supernatural number} is a formal product $\prod_pp^{n_p}$ over all primes with $n_p\in\{0,1,2,\dots,\infty\}$. The \emph{order} of a profinite group $G$ is
$\#G=\lcm_U\#(G/U)$ over the open normal subgroups $U$. $G$ is \emph{pro-$\ell$} if $\#G$ is a power of $\ell$ (equivalently, every finite continuous quotient
is an $\ell$-group), and has \emph{order prime to $\ell$} if $\ell\nmid\#G$. The \emph{maximal pro-$\ell$ quotient} of $G$ is $G^{(\ell)}=G/K_\ell$, where $K_\ell$ is
the intersection of the open normal subgroups of $\ell$-power index.
\end{definition}

\begin{lemma}\label{lem:cohomology:pro-l}
Let $G$ be profinite.
\begin{enumerate}[label=(\roman*)]
\item $G^{(\ell)}$ is pro-$\ell$, and every continuous homomorphism from $G$ to a pro-$\ell$ group factors through $G^{(\ell)}$. If $\ell\nmid\#G$ then $G^{(\ell)}=1$;
the converse fails (for $\ell=2$, the alternating group $A_4$ has order divisible by $2$ but no normal subgroup of index $2$ or $4$, so $A_4^{(2)}=1$).
\item The kernel $K_\ell$ of $G\to G^{(\ell)}$ has $K_\ell^{(\ell)}=1$.
\item A closed subgroup of a pro-$\ell$ group is pro-$\ell$. A finitely generated $\Zl$-module is a pro-$\ell$ group, and so is the group
$\{\sm1a01:a\in\Zl\}\simeq\Zl$ of unipotent upper triangular matrices.
\item (Stated.) Sylow pro-$\ell$ subgroups exist and are conjugate; see \cite[Chapter I, Section 1]{SerreGaloisCoh}.
\end{enumerate}
\end{lemma}

\begin{proof}
(i) The open normal subgroups of $\ell$-power index are closed under finite intersection (the quotient by an intersection embeds in a product of
$\ell$-groups), so $G/K_\ell$ is the inverse limit of the $\ell$-groups $G/U$; its open normal subgroups are images of open normal subgroups of $G$ containing
some such $U$, so its finite quotients are $\ell$-groups. If $f\colon G\to P$ is continuous with $P$ pro-$\ell$, then for every open normal $V\subseteq P$,
$f^{-1}(V)$ is open normal in $G$ with $G/f^{-1}(V)$ embedded in the $\ell$-group $P/V$, so $f^{-1}(V)\supseteq K_\ell$; as $\bigcap V=1$, $f(K_\ell)=1$. If $\ell\nmid\#G$, the
only open normal subgroup of $\ell$-power index is $G$. For $A_4$: its normal subgroups are $1$, $V_4$, $A_4$, of index $12$, $3$, $1$.

(ii) Suppose $L_0$ is an open normal subgroup of $K_\ell$ of index $\ell^a$. Since $K_\ell$ is closed, $L_0\supseteq K_\ell\cap U_0$ for some open normal subgroup
$U_0$ of $G$. Every conjugate $gL_0g^{-1}$ is again normal of index $\ell^a$ in $K_\ell$ (as $K_\ell$ is normal in $G$) and contains $K_\ell\cap U_0$, so there are only
finitely many, corresponding to subgroups of the finite group $K_\ell/(K_\ell\cap U_0)$. Their intersection $L$ is normal in $G$, and $K_\ell/L$ embeds in a finite
product of groups of order $\ell^a$, so it is a finite $\ell$-group. In the profinite group $G/L$ choose an open normal subgroup meeting the finite subgroup
$K_\ell/L$ trivially, and let $U$ be its preimage: $U$ is open normal in $G$ and $U\cap K_\ell=L$. Then $\#(G/U)=\#(UK_\ell/U)\cdot\#(G/UK_\ell)$, where
$UK_\ell/U\simeq K_\ell/L$ is an $\ell$-group and $G/UK_\ell$ is a finite quotient of the pro-$\ell$ group $G/K_\ell$. So $G/U$ is an $\ell$-group, hence $U\supseteq K_\ell$ by the
definition of $K_\ell$, so $L=K_\ell$ and $L_0=K_\ell$.

(iii) If $H\le P$ is closed, the open normal subgroups $H\cap V$ ($V$ open normal in $P$) form a basis at $1$ in $H$, and $H/(H\cap V)$ embeds in $P/V$. A finitely
generated $\Zl$-module is $\varprojlim M/\ell^kM$, with finite $\ell$-groups $M/\ell^kM$. The last group is isomorphic to $(\Zl,+)$ by $a\mapsto\sm1a01$.
\end{proof}

\begin{theorem}[vanishing for groups of order prime to $\ell$]\label{thm:cohomology:prime-to-l}
Let $G$ be profinite.
\begin{enumerate}[label=(\alph*)]
\item If $\ell\nmid\#G$ and $M$ is a discrete $\ell$-primary torsion module, then $H^n(G,M)=0$ for all $n\ge1$.
\item If $\ell\nmid\#G$ and $M$ is of type (C) or (V), then $H^1(G,M)=0$; and $H^n(G,M)=0$ for all $n\ge1$.
\item If $G$ acts trivially on $M$ of type (C), (V), or discrete $\ell$-primary, and $G^{(\ell)}=1$, then $H^1(G,M)=\Hom_{\mathrm{cont}}(G,M)=0$.
\item If $\ell\nmid\#G$ and $M$ is as in (a) or (b), then $M^G$ is a direct summand of $M$ as a $G$-module.
\end{enumerate}
\end{theorem}

\begin{proof}
(a) By \Cref{prop:cohomology:continuous}(ii), $H^n(G,M)=\varinjlim_UH^n(G/U,M^U)$. Each $\#(G/U)$ is prime to $\ell$ and $M^U$ is $\ell$-primary, so
$H^n(G/U,M^U)=0$ by \Cref{thm:cohomology:order-kills}(iii).

(b) For $M$ of type (C), $H^1(G,M)=\varprojlim_kH^1(G,M/\ell^kM)=0$ by \Cref{thm:cohomology:tate-limits}(ii) (proved for $n=1$) and (a). For $n\ge2$ the same argument
uses \Cref{thm:cohomology:tate-limits}(ii) in degree $n$, whose hypothesis holds because $H^{n-1}(G,M/\ell^kM)$ is finite (for $n=2$ it is $H^1$, which is $0$ by (a)); this part rests on the stated
\Cref{thm:cohomology:tate-limits}(i). For type (V) use \Cref{thm:cohomology:tate-limits}(iii) with a stable lattice.

(c) A continuous homomorphism $f\colon G\to M$ has compact image $f(G)$. For $M$ of type (C), $f(G)$ is a closed subgroup of $M$, hence pro-$\ell$ by
\Cref{lem:cohomology:pro-l}(iii); so $f$ factors through $G^{(\ell)}=1$. For type (V), $f(G)$ lies in some $\ell^{-a}T$ and the same argument applies. For discrete
$\ell$-primary $M$, $f(G)$ is a finite abelian $\ell$-group.

(d) For an open normal $U$ acting trivially on a finite or $\ell$-primary module $A$, put $e_U=\#(G/U)^{-1}\sum_{g\in G/U}g$ (the integer $\#(G/U)$ is prime to $\ell$,
hence invertible on $A$). If $U'\subseteq U$ also acts trivially, $\sum_{g\in G/U'}g=\#(U/U')\sum_{g\in G/U}g$ on $A$, so $e_{U'}=e_U$: the projector $e$ does not depend on $U$. It is
$G$-linear, has image in $A^G$ and is the identity on $A^G$ (\Cref{thm:cohomology:order-kills}(iii)). For $M$ of type (D) define $e$ on each element $m$ using any $U$
fixing $m$. For $M$ of type (C) the projectors on the $M/\ell^kM$ are compatible with reduction, and their limit is a projector on $M=\varprojlim M/\ell^kM$. For
type (V), $e$ on a stable lattice extends $\Ql$-linearly. In each case $M=M^G\oplus\ker e$.
\end{proof}

\begin{corollary}[cohomology through the pro-$\ell$ quotient]\label{cor:cohomology:pro-l-quotient}
Let $G$ be profinite, $K\trianglelefteq G$ closed, and $M$ of type (D) ($\ell$-primary), (C) or (V), with $K$ acting trivially on $M$.
\begin{enumerate}[label=(\alph*)]
\item If $K^{(\ell)}=1$, then $\Inf\colon H^1(G/K,M)\to H^1(G,M)$ is an isomorphism.
\item If $\ell\nmid\#K$, then $\Inf\colon H^n(G/K,M)\to H^n(G,M)$ is an isomorphism for every $n$ (for $M$ of type (D)).
\item If the image of $G$ in the automorphism group of $M$ is pro-$\ell$ (for example, unipotent upper triangular in a basis of $M\simeq\Zl^2$), the action
factors through $G^{(\ell)}$, and $H^1(G,M)=H^1(G^{(\ell)},M)$ by inflation.
\end{enumerate}
\end{corollary}

\begin{proof}
(a) By \Cref{thm:cohomology:prime-to-l}(c), $H^1(K,M)=\Hom_{\mathrm{cont}}(K,M)=0$; apply \Cref{thm:cohomology:inf-res}(d) and \Cref{rem:cohomology:inf-res-continuous}.
(b) By \Cref{thm:cohomology:prime-to-l}(a), $H^q(K,M)=0$ for $q\ge1$; apply \Cref{thm:cohomology:hochschild-serre}. This is stated, since the spectral sequence is.
(c) The action is a continuous homomorphism to a pro-$\ell$ group, so it factors through $G^{(\ell)}$ by \Cref{lem:cohomology:pro-l}(i); then (a) applies with
$K=K_\ell$, by \Cref{lem:cohomology:pro-l}(ii).
\end{proof}

\begin{proposition}[the inertia group at $q\ne\ell$]\label{prop:cohomology:tame-inertia}
Let $q\ne\ell$ be primes, $I=I_q\subset G_{\Q_q}$ the inertia group, and $t_\ell\colon I\to\Zl(1)$ the tame character of \Cref{ch:local}, with kernel $K$. Imported
from \Cref{ch:local}: $t_\ell$ is surjective, it identifies $I^{(\ell)}$ with $\Zl(1)$, and every finite quotient of $K$ has order prime to $\ell$. Let $t\in I$ map
to a topological generator of $\Zl(1)$. Then:
\begin{enumerate}[label=(\roman*)]
\item $\ell\nmid\#K$, so $K^{(\ell)}=1$, and $I/K\simeq\Zl(1)\simeq\Zl$ is procyclic, generated by the image of $t$.
\item For every $M$ of type (C) or (V) on which $I$ acts through a pro-$\ell$ group (for instance through unipotent matrices), $K$ acts trivially, and
$c\mapsto c(t)$ induces
\[
H^1(I,M)\ \simeq\ H^1(I/K,M)\ \simeq\ M/(t-1)M .
\]
\item (Given \Cref{thm:cohomology:hochschild-serre} and \Cref{thm:cohomology:tate-limits}(i), both stated.) $H^n(I,M)=0$ for $n\ge2$, for $M$ as in (ii) of type (C).
\item For \emph{every} $M$ of type (C) (resp. (V)), with no hypothesis on the image of $I$: $H^1(K,M)=0$, $M^K$ is of type (C) (resp. (V)), and
\[
\Inf\colon H^1(I/K,M^K)\xrightarrow{\ \sim\ }H^1(I,M),\qquad H^1(I,M)\simeq M^K/(t-1)M^K\quad\text{via }[c]\mapsto c(t).
\]
In particular $H^1(I,M)$ is a finitely generated $\Zl$-module (resp. a finite-dimensional $\Ql$-space), and $H^1(I,M)_{\mathrm{tor}}$ is finite.
\end{enumerate}
\end{proposition}

\begin{proof}
(i) is the imported statement. (ii) The image of $K$ in the pro-$\ell$ image of $I$ is a quotient of $K$, so its finite quotients have order prime to $\ell$
and are $\ell$-groups: it is trivial. So $K$ acts trivially, \Cref{cor:cohomology:pro-l-quotient}(a) gives the first isomorphism, and
\Cref{thm:cohomology:procyclic-compact} applied to $\Gamma=I/K\simeq\Zl$ gives the second. Composing, a cocycle $c$ on $I$ that is inflated from $I/K$ is sent to
$c(t)$; an arbitrary cocycle is cohomologous to an inflated one, and changing $c$ by $d^0m$ changes $c(t)$ by $(t-1)m$, so the map $[c]\mapsto c(t)\bmod(t-1)M$ is
defined on all of $H^1(I,M)$ by the same formula. (iii) By \Cref{thm:cohomology:prime-to-l}(a), $H^q(K,M/\ell^kM)=0$ for $q\ge1$, so by
\Cref{thm:cohomology:hochschild-serre} $H^n(I,M/\ell^kM)=H^n(I/K,M/\ell^kM)$, which is $0$ for $n\ge2$ (\Cref{thm:cohomology:procyclic-discrete}). Passing to the limit
with \Cref{thm:cohomology:tate-limits} (the $H^1(I,M/\ell^kM)$ being finite) gives $H^n(I,M)=0$.

(iv) $H^1(K,M)=0$ by \Cref{thm:cohomology:prime-to-l}(b) (its proof for $H^1$ uses only \Cref{thm:cohomology:tate-limits}(ii) in degree $1$, which is proved). $M^K$ is closed in $M$
(\Cref{rem:cohomology:inf-res-continuous}), so it is a $\Zl$-submodule of type (C) (resp. a subspace, of type (V)), on which $I/K$ acts continuously. By
\Cref{thm:cohomology:inf-res}(d), valid for continuous cochains by \Cref{rem:cohomology:inf-res-continuous}, $\Inf$ is an isomorphism. By (i), $I/K\simeq\Zl$ is procyclic,
generated by the image of $t$, so \Cref{thm:cohomology:procyclic-compact} gives $H^1(I/K,M^K)\simeq M^K/(t-1)M^K$, $[c]\mapsto c(t)$, and the finiteness statements of its last
paragraph. The composite is $[c]\mapsto c(t)$ as in (ii).
\end{proof}

\section{Procyclic groups}\label{sec:cohomology:procyclic}

The identity the paper uses, ``$H^1(I,M')=H^1(\langle t\rangle,M')=M'/(t-1)M'$ for every $\Z_\ell$-module $M'$ of finite type with this action'' (proof of Lemma
lem:tamfactor), is \Cref{thm:cohomology:procyclic-compact}. The same theorem, for the procyclic group $G_{\Q_p}/I_p\simeq\widehat\Z$, gives
$H^1(G/I,\Zl(1)^I)=\Zl(1)/(q-1)$ (Lemma lem:tamfactor) and $\Hur(\Q_p,V^D)=(V^D)^{I_p}/(\Frob_p-1)$ (after Lemma lem:dfg21). The formula is false for finite
cyclic groups, and it needs a different proof for discrete and for compact coefficients; we give both, with exact hypotheses.

\begin{lemma}[procyclic groups]\label{lem:cohomology:procyclic}
A profinite group $\Gamma$ is \emph{procyclic} if it is topologically generated by one element $F$ (the subgroup $\langle F\rangle$ is dense).
\begin{enumerate}[label=(\roman*)]
\item $\Gamma$ is abelian, and $k\mapsto F^k$ extends to a continuous surjection $\widehat\Z\to\Gamma$, where $\widehat\Z=\varprojlim_n\Z/n\simeq\prod_p\Zp$. So the procyclic groups are the
quotients of $\widehat\Z$, the groups $\prod_p\Z/p^{a_p}$ with $a_p\in\{0,1,\dots,\infty\}$ and $\Z/p^\infty:=\Zp$. Examples: $\widehat\Z=\Gal(\overline\F_q/\F_q)$ with
generator the Frobenius $x\mapsto x^q$; $\Zl$; $\Zl(1)$; the tame inertia group $I_q/P_q\simeq\prod_{p\ne q}\Zp(1)$.
\item If $\pi\colon\Gamma\to\Zl$ is a continuous surjection, then $\pi(F)\in\Zl^\times$, and $k\mapsto F^k$ is injective on $\Z$. Replacing $\pi$ by $\pi(F)^{-1}\pi$ we may and do
assume $\pi(F)=1$; then $\pi(F^k)=k$, and for $a\in\Zl$ we write $F^a$ for any element with $\pi(F^a)=a$ when $\pi$ is an isomorphism.
\item Let $M$ be of type (C). A continuous action of $\Gamma$ on $M$ is determined by $u=F|_M\in\Aut_{\Zl}(M)$. For $\Gamma=\widehat\Z$ every $u$ arises. For $\Gamma=\Zl$ exactly
the $u$ with $u^{\ell^N}\to1$ in $\Aut(M)$ as $N\to\infty$ arise; every unipotent $u$ (that is, $u-1$ nilpotent) is one of them.
\end{enumerate}
\end{lemma}

\begin{proof}
(i) $\langle F\rangle$ is abelian and dense, so $\Gamma$ is abelian. For each open normal $U$, $\Gamma/U$ is a finite cyclic group generated by the image of $F$,
of order $n_U$; the maps $\Z/n_U\to\Gamma/U$, $1\mapsto F$, are compatible and give a continuous $\widehat\Z\to\varprojlim\Gamma/U=\Gamma$, whose image is compact and dense, hence
everything. A quotient of $\widehat\Z=\prod_p\Zp$ by a closed subgroup is the product of the quotients of the $\Zp$, whose closed subgroups are $0$ and $p^a\Zp$.

(ii) $\pi(\Gamma)$ is the closure of $\pi(F)\Z$, which is $\pi(F)\Zl$; it is $\Zl$ only if $\pi(F)$ is a unit. Then $\pi(F^k)=k\pi(F)\ne0$ for $k\ne0$.

(iii) $u$ determines the action on the dense subgroup $\langle F\rangle$, hence everywhere. $\Aut_{\Zl}(M)=\varprojlim\Aut(M/\ell^kM)$ is profinite, so the closure of
$\langle u\rangle$ is procyclic and $1\mapsto u$ extends to $\widehat\Z$ by (i). If $\Zl$ acts continuously then $u^{\ell^N}=\rho(\ell^N)\to\rho(0)=1$. Conversely, if
$u^{\ell^N}\to1$, the map $k\mapsto u^k$ on $\Z$ is uniformly continuous for the $\ell$-adic topology (if $u^{\ell^N}$ acts trivially on $M/\ell^rM$, so does every power), and
it extends to $\Zl$ by completeness. If $u=1+X$ with $X^r=0$, then $u^{\ell^N}=\sum_{i<r}\binom{\ell^N}{i}X^i$ and $v_\ell\binom{\ell^N}{i}=N-v_\ell(i)\to\infty$ for $1\le i<r$.
\end{proof}

\begin{lemma}[cocycles on powers]\label{lem:cohomology:powers}
Let $\Gamma$ be procyclic with generator $F$, $M$ a Hausdorff topological $\Gamma$-module and $c$ a continuous $1$-cocycle. Put $S_k=1+F+\dots+F^{k-1}$.
\begin{enumerate}[label=(\roman*)]
\item $c(F^k)=S_k\,c(F)$ for $k\ge0$, and $c(F^{-k})=-F^{-k}c(F^k)$. So $c$ is determined by $c(F)$.
\item Suppose $\pi\colon\Gamma\to\Zl$ is a continuous surjection with $\pi(F)=1$, and $M$ is of type (C), (V), or finite of $\ell$-power order. Then $(F-1)M$ is
closed, and for every $\gamma\in\Gamma$
\[
c(\gamma)\equiv\pi(\gamma)\,c(F)\pmod{(F-1)M},
\]
where $\pi(\gamma)\in\Zl$ acts on the $\Zl$-module $M/(F-1)M$. In particular $c(F^a)\equiv a\,c(F)$ modulo $(F-1)M$ for every $a\in\Zl$ when $\pi$ is an isomorphism.
\end{enumerate}
\end{lemma}

\begin{proof}
(i) By induction, $c(F^{k+1})=c(F^k)+F^kc(F)$; and $c(F^{-k})=-F^{-k}c(F^k)$ by \Cref{prop:cohomology:h0-h1}(ii). By continuity $c$ is determined on the
closure of $\langle F\rangle$, which is $\Gamma$.

(ii) $(F-1)M$ is the image of a $\Zl$-linear endomorphism: a finitely generated, hence compact, submodule for type (C), a subspace for type (V). So
$M/(F-1)M$ is Hausdorff, and both sides of the congruence are continuous functions of $\gamma$ into it. On $M/(F-1)M$ the element $F$ acts as $1$, so
$S_k\equiv k$: explicitly $S_k-k=\sum_{i=0}^{k-1}(F^i-1)$ and each $F^i-1=(F-1)(1+F+\dots+F^{i-1})$ maps $M$ into $(F-1)M$. By (i) the two sides agree at
$\gamma=F^k$ for all $k\ge0$, and for $k<0$ too: $c(F^{k})=-F^{k}c(F^{-k})\equiv-(-k)c(F)=k\,c(F)$, since $F^k$ also acts as $1$ on $M/(F-1)M$. The $F^k$, $k\in\Z$,
form a dense subset of $\Gamma$, so the two sides agree everywhere.
\end{proof}

The congruence holds only modulo $(F-1)M$, and the paper needs no more. For $q^{-1}\in\Zl$, which is not an integer, there is no exact formula for
$c(F^{q^{-1}})$ as a finite sum; what (ii) says is that its class in $M/(F-1)M$ is $q^{-1}$ times that of $c(F)$, as if $F$ were $1$. For instance, for $\ell=5$ and
$q=11$ one has $q^{-1}\equiv91\pmod{125}$, so $c(F^{q^{-1}})\equiv91\,c(F)$ modulo $(F-1)M+125M$.

\begin{theorem}[Serre: procyclic groups, discrete coefficients]\label{thm:cohomology:procyclic-discrete}
Let either $\Gamma=\widehat\Z$ and $M$ a discrete torsion $\Gamma$-module, or $\Gamma=\Zl$ and $M$ a discrete $\ell$-primary torsion $\Gamma$-module; let $F$ be the
generator $1$. Then $H^0(\Gamma,M)=\ker(F-1)$, and $c\mapsto c(F)$ induces an isomorphism
\[
H^1(\Gamma,M)\ \simeq\ M/(F-1)M .
\]
If $M$ is finite, $\#H^1(\Gamma,M)=\#H^0(\Gamma,M)$. Moreover $H^n(\Gamma,M)=0$ for $n\ge2$.
\end{theorem}

\begin{proof}
$H^0$: an element fixed by $F$ is fixed by the closure of $\langle F\rangle$, since stabilisers are closed.

Injectivity of $c\mapsto c(F)$ modulo coboundaries. If $c(F)=(F-1)m$, then $c-d^0m$ vanishes at $F$, hence on $\langle F\rangle$ by \Cref{lem:cohomology:powers}(i), hence
everywhere, being continuous into a discrete (Hausdorff) space. Coboundaries map onto $(F-1)M$.

Surjectivity. Let $m\in M$. Its stabiliser is open, so it contains $n\Gamma$ for some $n\ge1$ (an $\ell$-power if $\Gamma=\Zl$): $F^nm=m$. Let $e$ be the order of
$m$ (an $\ell$-power in the second case). Then, with $S_k$ as in \Cref{lem:cohomology:powers},
\[
S_{ne}\,m=S_n\bigl(1+F^n+F^{2n}+\dots+F^{(e-1)n}\bigr)m=S_n(em)=0 ,
\]
using the factorisation $S_{ne}(F)=S_n(F)\,S_e(F^n)$ (group the $ne$ powers of $F$ into $e$ blocks of length $n$) and $F^{jn}m=m$. Define $c(F^k)=S_km$ for
$k\ge0$. Since $S_{k+j}=S_k+F^kS_j$, $c(F^{k+ne})=S_km+F^kS_{ne}m=c(F^k)$, so $c$ is a well-defined function on $\Z/ne$, a quotient of $\Gamma$ (the quotient of $\widehat\Z$
or of $\Zl$ by $ne$; in the second case $ne$ is a power of $\ell$); composed with $\Gamma\to\Z/ne$ it is continuous, since $\Z/ne$ is finite and discrete. The cocycle identity holds on $\Z/ne$, by the identity for $S$, hence on $\Gamma$. And $c(F)=m$.

If $M$ is finite, the exact sequence $0\to M^{F=1}\to M\xrightarrow{F-1}M\to M/(F-1)M\to0$ gives $\#H^0=\#H^1$. The vanishing for $n\ge2$ is stated; see
\cite[Chapter XIII, Section 1]{SerreLocalFields} and \cite[Chapter I, Sections 3 and 4]{SerreGaloisCoh} (the groups $\widehat\Z$ and $\Zl$ have cohomological dimension $1$).
\end{proof}

The compact case needs one more ingredient: geometric sums of an automorphism of a finitely generated $\Zl$-module become $\ell$-adically small along
suitable arithmetic progressions.

\begin{lemma}[geometric sums]\label{lem:cohomology:geometric-sums}
Let $M$ be a finitely generated $\Zl$-module, $u\in\Aut_{\Zl}(M)$, and $a_0$ the order of the image of $u$ in the finite group $\Aut(M/\ell M)$. For $j\ge1$ put
$S_j(u)=1+u+\dots+u^{j-1}$. Then $S_j(u)M\subseteq\ell^NM$ for every $N\ge0$ and every $j$ divisible by $a_0\ell^N$.
\end{lemma}

\begin{proof}
We use $S_{ab}(x)=S_a(x)\,S_b(x^a)$, which holds because both sides are $\sum_{i<a,\,j<b}x^{i+aj}$. First, if $v=1+Y$ with $YM\subseteq\ell M$, then
$(1+Y)^i=1+Y\cdot(\text{a polynomial in }Y)$, so $S_\ell(v)=\sum_{i=0}^{\ell-1}(1+Y)^i=\ell+YP(Y)$ for a polynomial $P$ with integer coefficients. Hence $S_\ell(v)$ maps
$\ell^iM$ into $\ell^{i+1}M$ for every $i$: $\ell\cdot\ell^iM=\ell^{i+1}M$, and $YP(Y)\ell^iM=\ell^iYP(Y)M\subseteq\ell^iYM\subseteq\ell^{i+1}M$, since $Y$ is $\Zl$-linear. Also
$v^\ell=1+Y'$ with $Y'=(1+Y)^\ell-1=Y\cdot(\ell+\dots)$, so again $Y'M\subseteq\ell M$. Now $w=u^{a_0}$ satisfies $(w-1)M\subseteq\ell M$, because $u^{a_0}$ is trivial on
$M/\ell M$ by the definition of $a_0$; so
\[
S_{\ell^N}(w)=S_\ell(w)\,S_\ell(w^\ell)\,S_\ell(w^{\ell^2})\cdots S_\ell(w^{\ell^{N-1}})
\]
maps $M$ into $\ell^NM$ (each of the $N$ factors $S_\ell(w^{\ell^i})$ raises the level by one, by the first step applied to $v=w^{\ell^i}$). Hence so does
$S_{a_0\ell^N}(u)=S_{a_0}(u)S_{\ell^N}(w)$, and so does $S_{a_0\ell^Nj'}(u)=S_{a_0\ell^N}(u)\,S_{j'}(u^{a_0\ell^N})$ for every $j'\ge1$ (all these operators are polynomials in $u$
and commute).
\end{proof}

For $M=\Zl$ the lemma is a statement about numbers: $u\in\Zl^\times$, and for $u\ne1$, $S_j(u)=(u^j-1)/(u-1)$. Take $\ell=5$ and $u=2$. The order of $2$ modulo $5$ is
$a_0=4$, and the lemma says $5^N\mid(2^j-1)/(2-1)$ whenever $4\cdot5^N\mid j$. In fact $v_5(2^{4\cdot5^N}-1)=N+1$ for $N=0,1,2,3$ (computed with PARI/GP), one more than the lemma
guarantees, because $v_5(2^4-1)=v_5(15)=1$ already. The lemma needs $a_0$: for $j=5^N$ alone, $2^{5^N}\equiv2\pmod5$, so $2^{5^N}-1$ is prime to $5$.

\begin{theorem}[procyclic groups, compact coefficients]\label{thm:cohomology:procyclic-compact}
Let $\Gamma$ be a procyclic group admitting a continuous surjection $\pi$ onto $\Zl$ (in $\Gamma\simeq\prod_p\Z/p^{a_p}$: $a_\ell=\infty$; examples $\Zl$, $\Zl(1)$, $\widehat\Z$,
and $I_q/P_q$ for $q\ne\ell$), $F$ a topological generator, and $M$ of type (C) or (V). Then:
\begin{enumerate}[label=(\roman*)]
\item $H^0(\Gamma,M)=\ker(F-1)$.
\item Every $m\in M$ is $c(F)$ for a unique continuous $1$-cocycle $c$, and $c\mapsto c(F)$ induces an isomorphism of $\Zl$-modules
\[
H^1(\Gamma,M)\ \simeq\ M/(F-1)M .
\]
\item $H^n(\Gamma,M)=0$ for $n\ge2$.
\end{enumerate}
Consequently $H^1(\Gamma,M)$ is a finitely generated $\Zl$-module (resp. a finite-dimensional $\Ql$-vector space), $H^1(\Gamma,M)_{\mathrm{tor}}=(M/(F-1)M)_{\mathrm{tor}}$ is
finite, and for $V$ of type (V)
\[
H^1(\Gamma,V)=V/(F-1)V=0\iff1\text{ is not an eigenvalue of }F\text{ on }V\iff\det(1-F\mid V)\ne0 .
\]
\end{theorem}

\begin{proof}
(i) As in \Cref{thm:cohomology:procyclic-discrete}: the set of elements fixed by $F$ is fixed by the closure of $\langle F\rangle$, as $M$ is Hausdorff.

(ii), injectivity and uniqueness. The map $Z^1_{\mathrm{cont}}(\Gamma,M)\to M$, $c\mapsto c(F)$, is $\Zl$-linear, and it is injective by \Cref{lem:cohomology:powers}(i): a
continuous cocycle vanishing at $F$ vanishes on the dense subgroup $\langle F\rangle$, hence everywhere. Coboundaries $d^0m$ map onto $(F-1)M$. So we get an injection
$H^1(\Gamma,M)\to M/(F-1)M$, and it remains to show that every $m\in M$ is $c(F)$ for some continuous cocycle $c$.

(ii), surjectivity, type (C). The idea: on the dense subgroup $\langle F\rangle\simeq\Z$ there is only one candidate, $F^k\mapsto S_k(u)m$, and the work is to
show that it is uniformly continuous for the topology induced from $\Gamma$, so that it extends to $\Gamma$. \Cref{lem:cohomology:geometric-sums} is exactly the
needed estimate: $F^j$ close to $1$ forces $S_j(u)m$ close to $0$. Normalise $\pi(F)=1$ (\Cref{lem:cohomology:procyclic}(ii)); then $k\mapsto F^k$ is injective, and we may define
$c_0\colon\langle F\rangle\to M$ by $c_0(F^k)=S_k(u)m$ for $k\ge0$ and $c_0(F^{-k})=-F^{-k}c_0(F^k)$, where $u=F|_M$. The identity $S_{k+j}=S_k+u^kS_j$ (for all $k,j\in\Z$, with
$S_{-k}=-u^{-k}S_k$) says that $c_0$ is a cocycle on $\langle F\rangle$. Let $a_0=m_0\ell^s$ ($\ell\nmid m_0$) be as in \Cref{lem:cohomology:geometric-sums} and, for $N\ge1$,
\[
U_N=\ker\bigl(\Gamma\to\Aut(M/\ell^NM)\bigr)\cap\pi^{-1}(\ell^{s+N}\Zl),
\]
an open subgroup of $\Gamma$ (\Cref{prop:cohomology:continuous}(iii)), decreasing in $N$. If $F^j\in U_N$, then $u^j$ is trivial on $M/\ell M$, so $a_0\mid j$, and
$\ell^{s+N}\mid\pi(F^j)=j$; so $a_0\ell^N\mid j$ and $c_0(F^j)\in\ell^NM$ by \Cref{lem:cohomology:geometric-sums} (for $j<0$ use $c_0(F^{j})=-F^{j}c_0(F^{-j})$ and $u(\ell^NM)=\ell^NM$).
Hence, for $F^k,F^{k'}$ in the same coset of $U_N$,
\[
c_0(F^{k'})-c_0(F^k)=F^k\,c_0(F^{k'-k})\in\ell^NM .
\]
Now define $c\colon\Gamma\to M$: for $\gamma\in\Gamma$ and $N\ge1$ the coset $\gamma U_N$ meets the dense subgroup $\langle F\rangle$; let $c(\gamma)\bmod\ell^NM$ be $c_0(F^k)\bmod\ell^NM$ for
any $F^k\in\gamma U_N$. This is independent of the choice by the display, compatible as $N$ grows because $U_{N+1}\subseteq U_N$, and so defines
$c(\gamma)\in\varprojlim_NM/\ell^NM=M$. It extends $c_0$, and $c(\gamma U_N)\subseteq c(\gamma)+\ell^NM$, so $c$ is continuous. For the cocycle identity, given $\gamma,\delta$ and $N$,
choose $F^k\in\gamma U_N$, $F^{k'}\in\delta U_N$; then $F^{k+k'}\in\gamma\delta U_N$ and
\[
c(\gamma\delta)\equiv c_0(F^{k+k'})=c_0(F^k)+F^kc_0(F^{k'})\equiv c(\gamma)+\gamma\,c(\delta)\pmod{\ell^NM},
\]
because $F^k\gamma^{-1}\in U_N$ acts trivially on $M/\ell^NM$. As $N$ is arbitrary, $c$ is a continuous cocycle with $c(F)=m$.

(ii), type (V). Let $T$ be a $\Gamma$-stable lattice (\Cref{prop:cohomology:continuous}(iv)) and $m\in V$; choose $a$ with $\ell^am\in T$, let $c_T$ be the cocycle with
$c_T(F)=\ell^am$ just constructed, and put $c=\ell^{-a}c_T$.

(iii) is stated; see \cite[Chapter I, Sections 3 and 4]{SerreGaloisCoh} ($\operatorname{cd}_\ell\Gamma\le1$) and \cite[Section 2]{Tate95}; it follows from
\Cref{thm:cohomology:procyclic-discrete} and \Cref{thm:cohomology:tate-limits} when $\Gamma$ is $\widehat\Z$ or $\Zl$. For a general $\Gamma\simeq\prod_p\Z/p^{a_p}$ with $a_\ell=\infty$,
first remove the prime-to-$\ell$ part: $\Gamma=\Gamma_\ell\times\Gamma'$ with $\Gamma_\ell\simeq\Zl$ and $\ell\nmid\#\Gamma'$, so $H^q(\Gamma',A)=0$ for $q\ge1$ and finite $\ell$-primary $A$
(\Cref{thm:cohomology:prime-to-l}(a)), and \Cref{thm:cohomology:hochschild-serre} gives $H^n(\Gamma,A)=H^n(\Gamma_\ell,A^{\Gamma'})$, which vanishes for $n\ge2$ by
\Cref{thm:cohomology:procyclic-discrete}; then pass to the limit over $A=M/\ell^kM$ with \Cref{thm:cohomology:tate-limits}.

The consequences: $M/(F-1)M$ is a quotient of a finitely generated module; for $V$, $F-1$ is surjective if and only if it is injective, if and only if
$\det(F-1)=\pm\det(1-F)\ne0$.
\end{proof}

\begin{example}\label{ex:cohomology:procyclic-examples}
Let $\Gamma$ and $F$ be as in \Cref{thm:cohomology:procyclic-compact}.
\begin{enumerate}[label=(\roman*)]
\item Trivial action on $M$ of type (C): $H^1(\Zl,M)=\Hom_{\mathrm{cont}}(\Zl,M)=M$, the homomorphism with value $m$ at $1$ being $a\mapsto am$; and
$H^1(\widehat\Z,M)=M$, every continuous homomorphism $\widehat\Z\to M$ factoring through $\widehat\Z\to\Zl$.
\item $\Gamma=\widehat\Z=\Gal(\overline\F_q/\F_q)$ acting on $\Zl(1)$, $q\ne\ell$, so that $F$ acts by $q$. Then $H^0=0$ (as $q\ne1$) and
$H^1(\widehat\Z,\Zl(1))=\Zl/(q-1)\Zl\simeq\Z/\ell^{v_\ell(q-1)}$. On $\mu_{\ell^k}$, by \Cref{thm:cohomology:procyclic-discrete}, $H^0=\{\zeta:\zeta^{q-1}=1\}$ and
$H^1=\mu_{\ell^k}/\mu_{\ell^k}^{q-1}$ both have order $\ell^{\min(k,v_\ell(q-1))}$. This is the group $H^1(G/I,\Zl(1)^I)=\Zl(1)/(q-1)\simeq\F_q^\times\otimes\Zl$ of Lemma
lem:tamfactor, for $G=G_{\Q_q}$, $I=I_q$ (on which $\chi_\ell$ is trivial).
\item $M=\Zl e_1\oplus\Zl e_2$ with $F=\sm1{n'}01$, $n'\ne0$: $(F-1)M=n'\Zl e_1$, so $H^1=(\Zl/n')\bar e_1\oplus\Zl\bar e_2$, with torsion $\Z/\ell^{v_\ell(n')}$.
\item $\Gamma=\widehat\Z$ and $F$ acting on $M$ of type (C) through a finite quotient of order $m$ prime to $\ell$: then $M=M^F\oplus(F-1)M$ (the projector is
$m^{-1}(1+F+\dots+F^{m-1})$, whose kernel is $(F-1)M$), and $H^1(\widehat\Z,M)=M/(F-1)M\simeq M^F$.
\item $V$ of type (V): $H^1(\Gamma,V)=0$ if and only if $\det(1-F\mid V)\ne0$.
\end{enumerate}
\end{example}

\begin{proof}[Verification]
(i) and (v) are special cases of \Cref{thm:cohomology:procyclic-compact}. (ii) $(F-1)\Zl(1)=(q-1)\Zl(1)$, and $\Zl/(q-1)$ has length $v_\ell(q-1)$; the finite case is the
order of the kernel and cokernel of $\zeta\mapsto\zeta^{q-1}$ on the cyclic group $\mu_{\ell^k}$. (iii) $(F-1)e_1=0$, $(F-1)e_2=n'e_1$. (iv) $N=1+F+\dots+F^{m-1}$
satisfies $N(F-1)=0$ and $N\equiv m$ modulo $F-1$; so $e=m^{-1}N$ is a projector with image $M^F$, and $1-e=m^{-1}(m-N)$ is divisible by $F-1$, which shows $\ker e=(F-1)M$.
\end{proof}

\begin{warning}[the hypotheses of the procyclic formula]\label{warn:cohomology:procyclic}
The identity $H^1(\Gamma,M)=M/(F-1)M$ holds: for $\Gamma$ procyclic with a surjection onto $\Zl$ and $M$ of type (C) or (V); for $\Gamma=\widehat\Z$ and $M$ discrete torsion;
for $\Gamma=\Zl$ and $M$ discrete $\ell$-primary. It fails in each of the following situations.
\begin{enumerate}[label=(\roman*)]
\item $\Gamma$ finite cyclic: $H^1=\ker N/(F-1)M$ (\Cref{thm:cohomology:cyclic}); with trivial action on $\Zl$ this is $0$, not $\Zl$.
\item $\Gamma=\Zp$ with $p\ne\ell$, acting on a $\Zl$-module of type (C): $\Gamma$ has order prime to $\ell$, so $H^1(\Gamma,M)=0$ by \Cref{thm:cohomology:prime-to-l}, while
$M/(F-1)M$ can be nonzero (\Cref{ex:cohomology:zp-not-zl}).
\item Discrete modules that are not torsion: $H^1(\widehat\Z,\Z)=0$ (\Cref{warn:cohomology:pitfalls}(vi)).
\end{enumerate}
The hypothesis of the paper, ``every $\Z_\ell$-module $M'$ of finite type with this action'', is the first case: $\Gamma=I/K\simeq\Zl$, $M'$ of type (C).
\end{warning}

The procyclic computations also give the five-term inflation-restriction sequence for compact and $\Ql$-coefficients in the one case the book needs,
$G_{\Q_p}\to G_{\Q_p}/I_p\simeq\widehat\Z$: \Cref{thm:cohomology:inf-res}(b) was stated only for discrete coefficients. The proof uses that $\widehat\Z$ is free on one
generator, so that $G\to G/H$ splits.

\begin{proposition}[inflation-restriction over $\widehat\Z$]\label{prop:cohomology:inf-res-zhat}
Let $G$ be profinite, $H\trianglelefteq G$ closed with $G/H\simeq\widehat\Z$, and $M$ of type (D) (discrete torsion), (C) or (V).
\begin{enumerate}[label=(\roman*)]
\item $\Res\colon H^1(G,M)\to H^1(H,M)^{G/H}$ is surjective, so
\[
0\to H^1(G/H,M^H)\xrightarrow{\Inf}H^1(G,M)\xrightarrow{\Res}H^1(H,M)^{G/H}\to0
\]
is exact.
\item (Given \Cref{thm:cohomology:procyclic-discrete} and \Cref{thm:cohomology:procyclic-compact}(iii) for $n=2$, stated.) $H^2(G/H,M^H)=0$, so the five-term sequence of
\Cref{thm:cohomology:inf-res}(b) holds for $M$, with $\operatorname{tg}=0$.
\end{enumerate}
\end{proposition}

\begin{proof}
(i) Exactness on the left and in the middle is \Cref{thm:cohomology:inf-res}(a) with \Cref{rem:cohomology:inf-res-continuous}. \emph{A splitting.} Choose $F\in G$ mapping to the
generator $1$ of $G/H=\widehat\Z$, and let $\Gamma_0$ be the closure of $\langle F\rangle$. By \Cref{lem:cohomology:procyclic}(i) there is a continuous surjection
$\widehat\Z\to\Gamma_0$, $1\mapsto F$; composed with $\Gamma_0\to G/H$ it is a continuous endomorphism of $\widehat\Z$ fixing $1$, hence the identity (it fixes the dense subgroup
$\Z$). So $\Gamma_0\to G/H$ is an isomorphism, $H\cap\Gamma_0=1$, $G=H\Gamma_0$, and $(h,\gamma)\mapsto h\gamma$ is a continuous bijection $H\times\Gamma_0\to G$ of compact
Hausdorff spaces, hence a homeomorphism.

\emph{Extending a cocycle.} Let $c\in Z^1(H,M)$ have $G/H$-invariant class. Then $F\cdot c-c=d^0m$ for some $m\in M$, that is $Fc(F^{-1}hF)-c(h)=hm-m$ for $h\in H$; replacing $h$
by $FhF^{-1}$,
\[
\gamma\,c(h)-c(\gamma h\gamma^{-1})=(\gamma h\gamma^{-1}-1)\,\tilde c(\gamma)\qquad(h\in H)\tag{$*$}
\]
holds for $\gamma=F$ and $\tilde c(F)=m$. Let $\tilde c\colon\Gamma_0\to M$ be the continuous cocycle with $\tilde c(F)=m$: it exists by \Cref{thm:cohomology:procyclic-compact}(ii) for
types (C), (V) and by \Cref{thm:cohomology:procyclic-discrete} for type (D), since $\Gamma_0\simeq\widehat\Z$. The set of $\gamma\in\Gamma_0$ satisfying $(*)$ for all $h$ is closed
(both sides are continuous in $\gamma$, and $M$ is Hausdorff) and stable under products: if $(*)$ holds for $\gamma$ and $\delta$, then
\[
\gamma\delta\,c(h)-c(\gamma\delta h\delta^{-1}\gamma^{-1})=\gamma(\delta h\delta^{-1}-1)\tilde c(\delta)+(\gamma\delta h\delta^{-1}\gamma^{-1}-1)\tilde c(\gamma)
=(\gamma\delta h\delta^{-1}\gamma^{-1}-1)\bigl(\tilde c(\gamma)+\gamma\tilde c(\delta)\bigr),
\]
which is $(*)$ for $\gamma\delta$ since $\tilde c(\gamma\delta)=\tilde c(\gamma)+\gamma\tilde c(\delta)$. It contains $F$, hence every $F^k$, $k\ge1$, hence their closure, which is $\Gamma_0$
($F^{-1}$ is the limit of $F^{k!-1}$, since $F^{k!}\to1$ in the profinite group $\Gamma_0$). Now define $c^\sharp(h\gamma)=c(h)+h\tilde c(\gamma)$ for $h\in H$, $\gamma\in\Gamma_0$. It is
continuous by the homeomorphism above, and it restricts to $c$ on $H$. For the cocycle identity, write $h_1\gamma_1h_2\gamma_2=h_1(\gamma_1h_2\gamma_1^{-1})\cdot\gamma_1\gamma_2$; then
\[
c^\sharp(h_1\gamma_1h_2\gamma_2)-c^\sharp(h_1\gamma_1)-h_1\gamma_1c^\sharp(h_2\gamma_2)
=h_1\Bigl(c(\gamma_1h_2\gamma_1^{-1})+(\gamma_1h_2\gamma_1^{-1}-1)\tilde c(\gamma_1)-\gamma_1c(h_2)\Bigr),
\]
after expanding $\tilde c(\gamma_1\gamma_2)=\tilde c(\gamma_1)+\gamma_1\tilde c(\gamma_2)$ and cancelling the terms $h_1\gamma_1h_2\tilde c(\gamma_2)$; the bracket is $0$ by $(*)$ for $\gamma_1$. So
$c^\sharp\in Z^1_{\mathrm{cont}}(G,M)$ and $\Res[c^\sharp]=[c]$.

(ii) $M^H$ is of the same type as $M$ (\Cref{rem:cohomology:inf-res-continuous}; a submodule of a discrete torsion module is discrete torsion), and
$H^2(\widehat\Z,M^H)=0$ by the cited results. With $H^2(G/H,M^H)=0$ the five-term sequence reduces to (i) followed by $0\to0\to H^2(G,M)$.
\end{proof}

\begin{remark}\label{rem:cohomology:inf-res-zhat-use}
For $G=G_{\Q_p}$ and $H=I_p$ this is the exact sequence $0\to\Hur(\Q_p,M)\to H^1(\Q_p,M)\to H^1(I_p,M)^{G_{\Q_p}/I_p}\to0$ for $M$ of type (C) or (V), at every prime $p$ (the
proof does not use $p\ne\ell$). Part (i) is proved here without any stated input; \Cref{ch:galcoh} reaches the same surjectivity through part (ii).
\end{remark}

\section[The model computation]{The model computation: inertia, Frobenius and the torsion of $H^1(I,\ad T)$}\label{sec:cohomology:model}

This section is the group cohomology of Lemma lem:tamfactor, done line by line. We isolate the abstract setting (a profinite group with a normal subgroup
whose $\ell$-adic part is procyclic, and an element normalising it by $t\mapsto t^q$), compute the conjugation action of Frobenius on $H^1$ of the inertia
group, and then evaluate everything for the adjoint of a two-dimensional representation of the Tate shape. What is not done here: the identification of
the datum $c'$ with the unit part of the Tate parameter, which is Kummer theory (\Cref{ch:galcoh}), and the definition of $\Tam_q$ and its use in Theorem
thm:main (\Cref{ch:blochkato}).

\begin{construction}[the model setting]\label{constr:cohomology:model}
The data are:
\begin{itemize}
\item a profinite group $G$, a closed normal subgroup $I\trianglelefteq G$, and a closed subgroup $K\subseteq I$, normal in $G$, of order prime to $\ell$, with $I/K\simeq\Zl$
topologically generated by the image of an element $t\in I$; we fix the isomorphism $\pi\colon I/K\to\Zl$ with $\pi(t)=1$ and write $t^a$ for an element of $I$ with
$\pi(t^aK)=a$, $a\in\Zl$;
\item an element $\Frob\in G$ whose image topologically generates $G/I\simeq\widehat\Z$;
\item \sloppy a positive integer $q$ prime to $\ell$ with $\Frob\,t\,\Frob^{-1}\equiv t^q\pmod K$, equivalently $\Frob^{-1}t\,\Frob\equiv t^{q^{-1}}\pmod K$ with $q^{-1}\in\Zl^\times$;
\item a finitely generated $\Zl$-module $M$ with a continuous action $\rho_M\colon G\to\Aut_{\Zl}(M)$ that is trivial on $K$.
\end{itemize}
The two forms of the relation are equivalent: conjugation by $\Frob^{-1}$ is a continuous automorphism $\alpha$ of $I/K\simeq\Zl$, hence multiplication by a unit;
the first form says $\alpha^{-1}(t)=t^q$, so $\alpha$ is multiplication by $q^{-1}$.
We write $\rho_M$, not $\rho$, because in the application $\rho$ is the two-dimensional representation on $T$ and $M$ is often $\ad T$, on which $g$ acts by
$\rho_M(g)=\Ad\rho(g)\colon x\mapsto\rho(g)x\rho(g)^{-1}$; for $M=T$, $\rho_M=\rho$. On the cohomology we keep the dot notation $\sigma\cdot c$ of
\Cref{prop:cohomology:conjugation}, and inside $M$ we also write $tm$ for $\rho_M(t)m$.

\emph{Dictionary with the paper.} $G=G_{\Q_q}$ for a prime $q\ne\ell$, $I=I_q$, $K=\ker(t_\ell\colon I_q\to\Zl(1))$ (of order prime to $\ell$ by \Cref{prop:cohomology:tame-inertia}), $t$
a lift of a generator of $\Zl(1)$, $\Frob$ a lift of the arithmetic Frobenius; the relation is the tame relation of \Cref{ch:local}. With $T=T_\ell E$, the module is
$M=\ad T$ for $E$ with multiplicative or potentially multiplicative reduction at $q$, and $M=T$ only for multiplicative reduction. In these cases the action of $I_q$ on
$M$ is unipotent in a suitable basis, hence has pro-$\ell$ image and is trivial on $K$ (\Cref{ch:galrep}; this is a hypothesis of the model, not proved here). At an
additive potentially multiplicative prime $T$ itself does not qualify: $E$ is a twist of a Tate curve by a ramified quadratic character, which is nontrivial on $I_q$ and,
for $\ell$ odd, nontrivial on $K$ (its restriction to $I_q$ has order $2$, prime to $\ell$, so it does not factor through the pro-$\ell$ quotient $I_q/K$); it cancels in
$\ad T$. The source of the definition of $\Tam_q$ that
uses this computation is \cite[I.4.2.2]{FPR94}, as recalled in \cite[p.~708]{DFG04}.
\end{construction}

\begin{proposition}[Frobenius on $H^1$ of inertia]\label{prop:cohomology:frob-h1}
In the model setting:
\begin{enumerate}[label=(\alph*)]
\item $\Inf\colon H^1(I/K,M)\to H^1(I,M)$ is an isomorphism, and $[c]\mapsto c(t)\bmod(t-1)M$ is an isomorphism $H^1(I,M)\simeq M/(t-1)M$.
\item $\rho_M(\Frob)\bigl((t-1)M\bigr)\subseteq(t-1)M$, and under the isomorphism of (a) the conjugation action of $G/I$ on $H^1(I,M)$ is given by
\[
\overline{\Frob}\cdot\bar m=q^{-1}\rho_M(\Frob)\,\bar m ;
\]
this does not depend on the lift $\Frob$ of the generator of $G/I$. For $M=\ad T$ the operator is $q^{-1}\Ad\rho(\Frob)$.
\item $H^1(I,M)_{\mathrm{tor}}=(M/(t-1)M)_{\mathrm{tor}}$ is finite and $G/I$-stable, and
\[
\bigl(H^1(I,M)_{\mathrm{tor}}\bigr)^{G/I}=\bigl\{\bar m\in(M/(t-1)M)_{\mathrm{tor}}:\ \rho_M(\Frob)\bar m=q\,\bar m\bigr\}.
\]
\end{enumerate}
\end{proposition}

\begin{proof}
(a) $K$ acts trivially and $K^{(\ell)}=1$ (\Cref{lem:cohomology:pro-l}(i)), so \Cref{cor:cohomology:pro-l-quotient}(a) gives the first isomorphism, and
\Cref{thm:cohomology:procyclic-compact} for $\Gamma=I/K\simeq\Zl$ the second; the composite is $[c]\mapsto c(t)$, as in the proof of \Cref{prop:cohomology:tame-inertia}(ii).

(b) Since $K$ acts trivially, $\rho_M(\Frob)\rho_M(t)\rho_M(\Frob)^{-1}=\rho_M(\Frob\,t\,\Frob^{-1})=\rho_M(t)^q$, so
\[
\rho_M(\Frob)(t-1)=(\rho_M(t)^q-1)\rho_M(\Frob)=(t-1)(1+t+\dots+t^{q-1})\rho_M(\Frob),
\]
and $\rho_M(\Frob)(t-1)M\subseteq(t-1)M$. (We do not need, and do not claim, equality.) Let $c$ be a cocycle on $I$ inflated from $I/K$. By
\Cref{prop:cohomology:conjugation}, $(\Frob\cdot c)(i)=\rho_M(\Frob)\,c(\Frob^{-1}i\,\Frob)$, again inflated from $I/K$ since $K$ is normal in $G$. At $i=t$:
$\Frob^{-1}t\,\Frob\equiv t^{q^{-1}}\pmod K$, so by \Cref{lem:cohomology:powers}(ii), applied to $\Gamma=I/K$ with $\pi(t)=1$,
\[
(\Frob\cdot c)(t)=\rho_M(\Frob)\,c(t^{q^{-1}})\equiv\rho_M(\Frob)\bigl(q^{-1}c(t)+x\bigr),\qquad x\in(t-1)M,
\]
and $\rho_M(\Frob)x\in(t-1)M$. So $\Frob\cdot[c]$ corresponds to $q^{-1}\rho_M(\Frob)c(t)$ modulo $(t-1)M$. By \Cref{prop:cohomology:conjugation}(ii), $I$ acts trivially on
$H^1(I,M)$, so the action of $\Frob i$ ($i\in I$) is that of $\Frob$.

(c) The isomorphism of (a) is $\Zl$-linear, so it identifies the torsion submodules; $(M/(t-1)M)_{\mathrm{tor}}$ is finite because $M/(t-1)M$ is finitely generated. The
action of $G/I$ is $\Zl$-linear, so it preserves the torsion. It remains to see that invariance under $G/I$ is invariance under $\Frob$. The same computation as
in (b) applies to any $\sigma\in G$: conjugation by $\sigma^{-1}$ acts on $I/K\simeq\Zl$ by multiplication by a unit $a(\sigma)$, so $\sigma t\sigma^{-1}\equiv t^{b}\pmod K$ with
$b=a(\sigma)^{-1}\in\Zl^\times$, and $\rho_M(\sigma)(t-1)=(\rho_M(t)^b-1)\rho_M(\sigma)$. Here $(\rho_M(t)^b-1)M\subseteq(t-1)M$: for $b\in\Z$ the polynomial $X^b-1$ is divisible by
$X-1$ (for $b<0$ use $X^b-1=-X^b(X^{-b}-1)$); for $b\in\Zl$ take integers $b_j\to b$, so that $\rho_M(t)^{b_j}m\to\rho_M(t)^bm$ by continuity of the action of $I/K\simeq\Zl$, and use
that $(t-1)M$ is closed (\Cref{lem:cohomology:powers}(ii)). So $\rho_M(\sigma)(t-1)M\subseteq(t-1)M$, and as in (b), $\sigma$ acts on $M/(t-1)M$ by
$a(\sigma)\rho_M(\sigma)$. Since $\sigma\mapsto a(\sigma)\in\Zl^\times$ and $\rho_M$ are continuous, the stabiliser of an element of the Hausdorff space $M/(t-1)M$ is a
closed subgroup of $G$. If it contains $\Frob$ it contains $I$ (which acts trivially) and the closure of the group generated by $I$ and $\Frob$, which is $G$
because $\Frob$ generates $G/I$ topologically. Finally $q^{-1}\rho_M(\Frob)\bar m=\bar m$ if and only if $\rho_M(\Frob)\bar m=q\bar m$.
\end{proof}

The factor $q^{-1}$ is the point of (b), and two sanity checks show where it comes from. If $G$ acts trivially on $M=\Zl$, then $H^1(I,\Zl)=\Hom_{\mathrm{cont}}(I,\Zl)$,
and $\Frob$ acts on a homomorphism $c$ by $c\mapsto c(\Frob^{-1}\cdot\Frob)$; since $\Frob^{-1}t\Frob\equiv t^{q^{-1}}$, the value at $t$ is multiplied by $q^{-1}$, in
accordance with $q^{-1}\rho_M(\Frob)=q^{-1}$. So $H^1(I,\Zl)\simeq\Zl(-1)$ as a $G/I$-module: inertia's $\ell$-part is $\Zl(1)$, and homomorphisms out of it form its
dual. If instead $M=\Zl(1)$, with $\Frob$ acting by $q$, the operator is $q^{-1}q=1$, and $H^1(I,\Zl(1))=\Hom(\Zl(1),\Zl(1))\simeq\Zl$ with trivial Frobenius action. Only
the second module has nonzero Frobenius-invariants; the same effect, the eigenvalue $q$ of $\rho_M(\Frob)$ being the one that survives, drives the invariants
in \Cref{prop:cohomology:model}.

\begin{remark}[the other convention]\label{rem:cohomology:other-convention}
Some texts let $\sigma$ act by $c\mapsto\sigma^{-1}c(\sigma\,\cdot\,\sigma^{-1})$. With that convention the computation of (b) gives $q\,\rho_M(\Frob)^{-1}$ instead of $q^{-1}\rho_M(\Frob)$.
The invariants are the same, since one operator is the inverse of the other; the formulas are not. We use the paper's convention throughout.
\end{remark}

\begin{lemma}[the Tate shape and its adjoint]\label{lem:cohomology:tate-shape}
In the model setting let $T=\Zl e_1\oplus\Zl e_2$ and
\[
\rho(g)=\Mtwo{\chi(g)\varepsilon(g)}{b(g)}{0}{\varepsilon(g)},
\]
where $\varepsilon\colon G\to\{\pm1\}$ and $\chi\colon G\to\Zl^\times$ are continuous characters trivial on $I$ with $\chi(\Frob)=q$, and $b|_I\colon I\to\Zl$ is a
continuous homomorphism, trivial on $K$, with $b(t)=n'\ne0$. (In the paper $b=\varepsilon\kappa$, so $b|_I=\kappa|_I$, by \Cref{prop:cohomology:corner}(v).) Write
$\varepsilon=\varepsilon(\Frob)$ and $c'=b(\Frob)$, so that
\[
\rho(t)=\Mtwo{1}{n'}{0}{1},\qquad \rho(\Frob)=\Mtwo{\varepsilon q}{c'}{0}{\varepsilon}.
\]
Let $M=\ad T$, with $g$ acting by $\rho_M(g)=\Ad\rho(g)\colon x\mapsto\rho(g)x\rho(g)^{-1}$ (written $g\cdot x$), and basis $e=E_{12}$, $h=E_{11}-E_{22}$, $f=E_{21}$. Then
\begin{align*}
(t-1)e&=0, & (t-1)h&=-2n'e, & (t-1)f&=n'h-n'^2e,\\
\Frob\cdot e&=qe, & \Frob\cdot h&=h-2\varepsilon c'e, & \Frob\cdot f&=q^{-1}\bigl(f+\varepsilon c'h-c'^2e\bigr).
\end{align*}
Also $(t-1)^3=0$ on $\ad T$, and $(t-1)e_1=0$, $(t-1)e_2=n'e_1$ on $T$.
\end{lemma}

\begin{proof}
$\rho(t)=1+n'e$, since $\chi$, $\varepsilon$ are trivial on $I$, and $\rho(t)^{-1}=1-n'e$ as $e^2=0$. The matrix units give $eh=-e$, $he=e$, $ehe=0$, $ef-fe=h$, $efe=e$.
So $t\cdot x=(1+n'e)x(1-n'e)=x+n'(ex-xe)-n'^2exe$. For $x=e$: $e$. For $x=h$: $h+n'(-2e)$. For $x=f$: $f+n'h-n'^2e$.

(Here $eh=E_{12}(E_{11}-E_{22})=-E_{12}$, $he=E_{12}$, $ef=E_{11}$, $fe=E_{22}$, using $E_{ij}E_{kl}=\delta_{jk}E_{il}$.)

For $A=\sm\alpha\beta0\delta$, $A^{-1}=\sm{\alpha^{-1}}{-\beta\alpha^{-1}\delta^{-1}}0{\delta^{-1}}$, and a direct multiplication gives
\begin{align*}
AeA^{-1}&=\Mtwo0\alpha00A^{-1}=\Mtwo0{\alpha\delta^{-1}}00=\alpha\delta^{-1}e,\\
AhA^{-1}&=\Mtwo\alpha{-\beta}0{-\delta}A^{-1}=\Mtwo1{-2\beta\delta^{-1}}0{-1}=h-2\beta\delta^{-1}e,\\
AfA^{-1}&=\Mtwo\beta0\delta0A^{-1}=\Mtwo{\beta\alpha^{-1}}{-\beta^2\alpha^{-1}\delta^{-1}}{\delta\alpha^{-1}}{-\beta\alpha^{-1}}=\delta\alpha^{-1}f+\beta\alpha^{-1}h-\beta^2\alpha^{-1}\delta^{-1}e.
\end{align*} With $\alpha=\varepsilon q$, $\beta=c'$, $\delta=\varepsilon$ and $\varepsilon^{-1}=\varepsilon$:
$\alpha\delta^{-1}=q$, $\beta\delta^{-1}=\varepsilon c'$, $\delta\alpha^{-1}=q^{-1}$, $\beta\alpha^{-1}=\varepsilon c'q^{-1}$, $\beta^2\alpha^{-1}\delta^{-1}=c'^2q^{-1}$. Finally
$(t-1)^2f=n'(t-1)h=-2n'^2e$ and $(t-1)^3f=0$. The symbolic computation was also checked with PARI/GP (\Cref{ex:cohomology:identities}).
\end{proof}

\begin{proposition}[the model computation]\label{prop:cohomology:model}
Let $\ell$ be odd, $M=\ad T$ as in \Cref{lem:cohomology:tate-shape}, and $w=v_\ell(n')$. (We write $w$, not $v$, because in the statement of the paper's Lemma
lem:tamfactor $v$ denotes $v_\ell(q-1)$.)
\begin{enumerate}[label=(\alph*)]
\item $(t-1)M=n'(\Zl e\oplus\Zl h)$, so
\begin{gather*}
H^1(I,M)=M/(t-1)M=(\Zl/n')\bar e\oplus(\Zl/n')\bar h\oplus\Zl\bar f,\\
H^1(I,M)_{\mathrm{tor}}=(\Z/\ell^w)\bar e\oplus(\Z/\ell^w)\bar h\simeq(\Z/\ell^w)^2 .
\end{gather*}
\item $\overline\Frob$ acts on the torsion by $\bar e\mapsto\bar e$ and $\bar h\mapsto q^{-1}(\bar h-2\varepsilon c'\bar e)$.
\item With $m=\min\bigl(w,\,v_\ell(q-1),\,v_\ell(c')\bigr)$,
\[
\bigl(H^1(I,M)_{\mathrm{tor}}\bigr)^{G/I}=\{x\bar e+y\bar h:\ (q-1)y\equiv0,\ c'y\equiv0\pmod{\ell^w}\}\ \simeq\ \Z/\ell^w\oplus\Z/\ell^m,
\]
of order $\ell^{w+m}$, length $w+m$ and Fitting ideal $\ell^{w+m}\Zl$.
\item $m$ depends on $c'$ only modulo $\ell^{\min(w,v_\ell(q-1))}\Zl=(n',q-1)\Zl$. This is exactly the indeterminacy of $c'=b(\Frob)$: replacing $\Frob$ by $\Frob\,t^a$ ($a\in\Zl$)
changes $c'$ by $\varepsilon qn'a$, and replacing $e_2$ by $e_2+re_1$ (which changes the cocycle $b/\varepsilon$ by a coboundary, \Cref{prop:cohomology:corner}(ii)) changes $c'$
by $\varepsilon r(q-1)$; rescaling $e_1$, $e_2$ multiplies $c'$ by a unit.
\item (Check with $M=T$.) $H^1(I,T)=(\Zl/n')\bar e_1\oplus\Zl\bar e_2$, with torsion $(\Z/\ell^w)\bar e_1$, on which $\overline\Frob$ acts by $q^{-1}\cdot\varepsilon q=\varepsilon$. So
$(H^1(I,T)_{\mathrm{tor}})^{G/I}$ is $(\Z/\ell^w)\bar e_1$ if $\varepsilon=1$ and $0$ if $\varepsilon=-1$, with Fitting ideal $\ell^w\Zl$, resp. $\Zl$.
\end{enumerate}
\end{proposition}

\begin{proof}
(a) By \Cref{lem:cohomology:tate-shape}, $(t-1)M$ is spanned by $-2n'e$ and $n'h-n'^2e$, that is $(t-1)M=n'(\Zl\cdot2e+\Zl(h-n'e))$. As $\ell$ is odd, $2\in\Zl^\times$, so this
is $n'(\Zl e+\Zl(h-n'e))=n'(\Zl e\oplus\Zl h)$. Since $e,h,f$ is a basis, $M/(t-1)M=(\Zl/n')\bar e\oplus(\Zl/n')\bar h\oplus\Zl\bar f$, and $\Zl/n'\Zl=\Z/\ell^w$ (if $w=0$ it is $0$).
\Cref{prop:cohomology:frob-h1}(a) identifies this with $H^1(I,M)$.

(b) By \Cref{prop:cohomology:frob-h1}(b), $\overline\Frob$ acts by $q^{-1}\Ad\rho(\Frob)$, that is $x\mapsto q^{-1}\,\Frob\cdot x$; by \Cref{lem:cohomology:tate-shape}, $q^{-1}\Frob\cdot e=e$ and $q^{-1}\Frob\cdot h=q^{-1}(h-2\varepsilon c'e)$.

(c) $\overline\Frob(x\bar e+y\bar h)-(x\bar e+y\bar h)=-2\varepsilon c'q^{-1}y\,\bar e+(q^{-1}-1)y\,\bar h$. This vanishes if and only if $c'y\equiv0$ and $(q-1)y\equiv0\pmod{\ell^w}$, because
$2$, $\varepsilon$, $q$ are units and $\bar e,\bar h$ are independent generators of $(\Z/\ell^w)^2$. By \Cref{prop:cohomology:frob-h1}(c) these are the $G/I$-invariants. The
condition on $y\in\Z/\ell^w$ is $\ell^{\min(w,v_\ell(q-1),v_\ell(c'))}y=0$, which cuts out the subgroup $\ell^{w-m}\Z/\ell^w\simeq\Z/\ell^m$; $x$ is free. In detail: for
$a\in\Zl$ with $\alpha=v_\ell(a)$ (and $\alpha=\infty$ if $a=0$), the condition $ay\equiv0\pmod{\ell^w}$ on $y\in\Z/\ell^w$ is $\ell^{\min(\alpha,w)}y=0$, because $a=\ell^\alpha\cdot$unit. Two
such conditions, for $a=q-1$ and $a=c'$, hold together if and only if $\ell^{\mu}y=0$ with $\mu$ the smaller exponent, $\mu=\min(w,v_\ell(q-1),v_\ell(c'))=m$. The
elements of $\Z/\ell^w$ killed by $\ell^m$ are the multiples of $\ell^{w-m}$, a cyclic group of order $\ell^m$. The order and length follow,
and $\Fitt_{\Zl}(A)=\ell^{\len A}\Zl$ for finite $A$ (\Cref{ch:modules}).

(d) If $c''=c'+\ell^\mu z$ with $\mu=\min(w,v_\ell(q-1))$ and $z\in\Zl$, then $\min(\mu,v_\ell(c''))=\min(\mu,v_\ell(c'))$, so $m$ is unchanged. The first change:
$\rho(\Frob\,t^a)=\rho(\Frob)\rho(t)^a=\sm{\varepsilon q}{c'}{0}{\varepsilon}\sm1{n'a}01=\sm{\varepsilon q}{c'+\varepsilon qn'a}{0}{\varepsilon}$. The second: by the formula in the proof of
\Cref{prop:cohomology:corner}(ii), the corner becomes $b+r(\chi\varepsilon-\varepsilon)$, whose value at $\Frob$ is $c'+\varepsilon r(q-1)$; on $I$ it is unchanged, so $n'$ is unchanged.
In both cases the change is a multiple of $n'$ or of $q-1$, hence lies in $(n',q-1)\Zl=\ell^\mu\Zl$.

(e) By \Cref{lem:cohomology:tate-shape}, $(t-1)T=n'\Zl e_1$, which gives $H^1(I,T)$ by \Cref{prop:cohomology:frob-h1}(a). Here $\rho_M=\rho$, and the action is $q^{-1}\rho(\Frob)\bar e_1=q^{-1}\varepsilon q\bar e_1$.
If $\varepsilon=-1$ the invariants are the $x$ with $2x\equiv0\pmod{\ell^w}$, which is $x\equiv0$ for $\ell$ odd.
\end{proof}

A numerical instance, with every step visible. Take $\ell=3$, $q=7$, $n'=9$ and $c'=3$, so $w=2$, $v_3(q-1)=v_3(6)=1$, $v_3(c')=1$ and $m=1$. In the basis $e,h,f$ the
matrix of $t-1$ has columns $0$, $(-18,0,0)$, $(-81,9,0)$; its elementary divisors over $\Z$ are $9,18,0$ (PARI/GP, \texttt{matsnf}), that is $9,9,0$ up to units of $\Z_3$, so
$H^1(I,\ad T)_{\mathrm{tor}}\simeq(\Z/9)^2$. The invariance conditions on $x\bar e+y\bar h$ are $6y\equiv0$ and $3y\equiv0\pmod9$, satisfied by $y\in\{0,3,6\}$, and $x$
is arbitrary: the invariants are $\Z/9\oplus3\Z/9\simeq\Z/9\oplus\Z/3$, of order $27=3^{w+m}$. Replacing $c'$ by a unit, say $c'=1$, gives $m=0$ and invariants
$\Z/9$; replacing $q=7$ by $q=19$ ($v_3(18)=2$) and keeping $c'=3$ still gives $m=1$, now limited by $c'$.

\begin{corollary}[the Fitting ideal of the invariants]\label{cor:cohomology:fitting}
In the model setting with $\ell$ odd, the ideal $\Fitt_{\Zl}\bigl((H^1(I,M)_{\mathrm{tor}})^{G/I}\bigr)$ is defined (the module is finite), and equals $\ell^{w+m}\Zl$ for $M=\ad T$ and
$\ell^w\Zl$ or $\Zl$ (as $\varepsilon=1$ or $-1$) for $M=T$. For $G=G_{\Q_q}$ this ideal is the Tamagawa factor $\Tam_q(M)$ of the paper's Section 1.2, defined in
\Cref{ch:blochkato}; so
\[
v_\ell\bigl(\Tam_q(\ad T)\bigr)=v_\ell(n')+\min\bigl(v_\ell(n'),v_\ell(q-1),v_\ell(c')\bigr),
\]
which is the formula of Lemma lem:tamfactor once $v_\ell(n')=v_\ell(n)$ (\Cref{ch:local,ch:galrep}) and $\min(v_\ell(n),v_\ell(q-1),v_\ell(c'))=\min(v_\ell(n),v_\ell(q-1),v_\ell(c))$ for the
paper's class $c$ of the unit part of the Tate parameter (\Cref{ch:galcoh}). Lemma lem:tamfactor itself is assembled in \Cref{ch:blochkato}; the Kummer
identification of $c'$ and its indeterminacy modulo $(n,q-1)$ are proved in \Cref{ch:galcoh}.
\end{corollary}

\begin{proof}
\Cref{prop:cohomology:model}(c), (e) and $\Fitt_{\Zl}(\bigoplus\Z/\ell^{a_i})=\ell^{\sum a_i}\Zl$ (\Cref{ch:modules}). The identifications of $v_\ell(n')$ and of $c'$ are proved in the
chapters cited.
\end{proof}

\begin{remark}\label{rem:cohomology:model-remarks}
\begin{enumerate}[label=(\roman*)]
\item If $w=0$ every torsion module above is $0$ and every Fitting ideal is $\Zl$, consistently with the formula ($m\le w=0$).
\item The free part $\Zl\bar f$ of $H^1(I,\ad T)$ is not $\Frob$-stable: $q^{-1}\Frob\cdot\bar f=q^{-2}(\bar f+\varepsilon c'\bar h-c'^2\bar e)$ by \Cref{lem:cohomology:tate-shape}. Only the
torsion is needed, and it is stable (\Cref{prop:cohomology:frob-h1}(c)).
\item For $M=T$ the answer $\ell^{v_\ell(n)}$ (split, $\varepsilon=1$) or $1$ (nonsplit, $\varepsilon=-1$) is the $\ell$-part of the Tamagawa number $c_q(E)$, which is $n$ for
split and $1$ or $2$ for nonsplit multiplicative reduction (\Cref{ch:reduction}). This is the paper's ``check''.
\item The hypothesis $\ell$ odd enters only through $2\in\Zl^\times$ in (a), (c) and (e). For $\ell=2$ the answer changes (\Cref{ex:cohomology:ell2}).
\end{enumerate}
\end{remark}

\begin{paperlink}\label{pl:cohomology:tamfactor}
The proof of Lemma lem:tamfactor in the revised paper, sentence by sentence (quotations verbatim, from the proof):
\begin{itemize}
\item ``the kernel of $I\to\langle t\rangle$ has trivial pro-$\ell$ quotient, so by inflation-restriction $H^1(I,M')=H^1(\langle t\rangle,M')=M'/(t-1)M'$ for every
$\Z_\ell$-module $M'$ of finite type with this action'': \Cref{cor:cohomology:pro-l-quotient}(a) (which is \Cref{thm:cohomology:inf-res}(d) with
\Cref{thm:cohomology:prime-to-l}(c)) and \Cref{thm:cohomology:procyclic-compact}; the identification is $[c]\mapsto c(t)$ (\Cref{prop:cohomology:frob-h1}(a)).
\item ``$(t-1)e=0$, $(t-1)h=-2n'e$, $(t-1)f=n'h-n'^2e$, and as $2$ is a unit, $(t-1)M=n'(\Z_\ell e\oplus\Z_\ell h)$; so
$H^1(I,M)=(\Z_\ell/n')\bar e\oplus(\Z_\ell/n')\bar h\oplus\Z_\ell\bar f$ and its torsion is $(\Z/\ell^{v_\ell(n)})^2$'': \Cref{lem:cohomology:tate-shape} and
\Cref{prop:cohomology:model}(a), with $v_\ell(n')=v_\ell(n)$.
\item ``The group $G/I$ acts on $H^1(I,M)$ by $(\sigma\cdot c)(i)=\sigma c(\sigma^{-1}i\sigma)$'': \Cref{prop:cohomology:conjugation}.
\item ``for the arithmetic Frobenius $\Frob$, $\Frob^{-1}t\Frob=t^{q^{-1}}$ and $c(t^k)\equiv k\,c(t)$ modulo $(t-1)M$, so $\Frob$ acts on $M/(t-1)M$ as $q^{-1}\rho(\Frob)$'': the
relation is \Cref{constr:cohomology:model} (from \Cref{ch:local}), the congruence is \Cref{lem:cohomology:powers}(ii) with $k=q^{-1}\in\Zl$, the conclusion is
\Cref{prop:cohomology:frob-h1}(b). Here the paper's $\rho$ is the representation on $T$, so on $M=\ad T$ its $q^{-1}\rho(\Frob)$ is a slip for the conjugation
$q^{-1}\Ad\rho(\Frob)$, which is what its next sentence (``the adjoint action'') computes.
\item ``the adjoint action sends $e$ to $qe$ and $h$ to $h-2c'\varepsilon e$, so on the torsion $\bar e\mapsto\bar e$ and $\bar h\mapsto q^{-1}(\bar h-2c'\varepsilon\bar e)$'':
\Cref{lem:cohomology:tate-shape} and \Cref{prop:cohomology:model}(b).
\item ``The fixed points are the $x\bar e+y\bar h$ with $(q-1)y\equiv0$ and $c'y\equiv0\pmod{\ell^{v_\ell(n)}}$'': \Cref{prop:cohomology:model}(c), where
\Cref{prop:cohomology:frob-h1}(c) supplies the fact, implicit in the paper, that invariance under $G/I$ is invariance under $\Frob$.
\item ``in general it is well defined only modulo $\ell^{\min(v_\ell(n),v_\ell(q-1))}$ (changing the lift changes it by a multiple of $n'$, changing the cocycle by a
multiple of $q-1$)'': \Cref{prop:cohomology:model}(d).
\item ``As a check, the same computation for $M=T$ gives $(\Z/\ell^{v_\ell(n)})\bar e_1$ with $\Frob$ acting by $\varepsilon$'': \Cref{prop:cohomology:model}(e).
\end{itemize}
The remaining steps (that $\kappa|_I$ is $n$ times the tame character, the splitting $\kappa=n\kappa_q+\kappa_u$, and the value of $c'$) are Kummer theory and are proved in
\Cref{ch:galcoh}; the assembly is in \Cref{ch:blochkato}. In the paper's sentence the invariants are taken in $(\Z/\ell^{v_\ell(n)})^2$, the torsion, as the definition of
$\Tam$ requires; with this reading, and with $\Ad\rho(\Frob)$ for $\rho(\Frob)$ in the fourth quoted sentence, every step above is correct.
\end{paperlink}

\section{Cup products, Pontryagin duality and lengths}\label{sec:cohomology:cup}

The paper uses duality in one place: the exact sequence of Lemma lem:dfg21 has Pontryagin duals as its last two terms, and the proof of Theorem thm:main
adds lengths along it. The duality that produces those dual terms is Tate's local duality, a cup product pairing proved in \Cref{ch:galcoh}. Here we define
cup products, recall what Pontryagin duality does to lengths, and prove the length bookkeeping.

\begin{definition}[cup products]\label{def:cohomology:cup}
Let $G$ be a group (or a profinite group, with continuous cochains and modules of type (D), or of type (C) and (V) with the tensor product over $\Zl$) and $A$, $B$
two $G$-modules, $A\otimes B$ with the diagonal action. For $a\in C^p(G,A)$, $b\in C^q(G,B)$ put
\[
(a\cup b)(g_1,\dots,g_{p+q})=a(g_1,\dots,g_p)\otimes g_1\cdots g_p\,b(g_{p+1},\dots,g_{p+q})\in A\otimes B .
\]
For $p=q=1$: $(a\cup b)(g,h)=a(g)\otimes g\,b(h)$. A $G$-equivariant pairing $A\times B\to C$ composed with $\cup$ gives $C^p(G,A)\times C^q(G,B)\to C^{p+q}(G,C)$.
\end{definition}

\begin{proposition}[properties of the cup product]\label{prop:cohomology:cup}
\begin{enumerate}[label=(\roman*)]
\item $d(a\cup b)=da\cup b+(-1)^pa\cup db$. Hence $\cup$ induces bilinear maps $H^p(G,A)\times H^q(G,B)\to H^{p+q}(G,A\otimes B)$.
\item For $p=q=1$ directly: if $a$, $b$ are $1$-cocycles then $a\cup b$ is a $2$-cocycle; if $a=d^0m$ then $a\cup b=d^1(m\otimes b)$ with $(m\otimes b)(g)=m\otimes b(g)$; if
$b=d^0m'$ then $a\cup b=-d^1\chi$ with $\chi(g)=a(g)\otimes gm'$.
\item (Stated.) The cup product is associative and graded commutative, $a\cup b=(-1)^{pq}\tau_*(b\cup a)$ with $\tau\colon B\otimes A\to A\otimes B$ the flip; it commutes with
restriction and inflation; it is compatible with connecting homomorphisms (up to the sign $(-1)^p$ when the connecting map acts on the second factor); and
$\Cor(\Res a\cup b)=a\cup\Cor b$.
\end{enumerate}
\end{proposition}

\begin{proof}
(i) in general is a direct computation; see \cite[Chapter I, Section 4]{NSW} and \cite[Chapter V, Section 3]{Brown}. We check the case used most, (ii).
Let $\psi(g,h)=a(g)\otimes gb(h)$ with $a$, $b$ cocycles. In $(d^2\psi)(g,h,k)=g\psi(h,k)-\psi(gh,k)+\psi(g,hk)-\psi(g,h)$ substitute $a(gh)=a(g)+ga(h)$ and $b(hk)=b(h)+hb(k)$:
\[
ga(h)\otimes ghb(k)-\bigl(a(g)+ga(h)\bigr)\otimes ghb(k)+a(g)\otimes\bigl(gb(h)+ghb(k)\bigr)-a(g)\otimes gb(h)=0 .
\]
If $a(g)=gm-m$: $d^1(m\otimes b)(g,h)=gm\otimes gb(h)-m\otimes b(gh)+m\otimes b(g)=gm\otimes gb(h)-m\otimes gb(h)=a(g)\otimes gb(h)$, so $a\cup b=d^1(m\otimes b)$. If $b(h)=hm'-m'$: $d^1\chi(g,h)=ga(h)\otimes ghm'-a(gh)\otimes ghm'+a(g)\otimes gm'=-a(g)\otimes ghm'+a(g)\otimes gm'=-(a\cup b)(g,h)$. So the cup product of
classes is well defined in degrees $(1,1)$. (iii) is stated with the same references and \cite[Chapter VIII, Section 3]{SerreLocalFields}.
\end{proof}

\begin{remark}\label{rem:cohomology:cup-uses}
In \Cref{ch:galcoh} the cup product $H^i(\Q_p,M)\times H^{2-i}(\Q_p,M^*(1))\to H^2(\Q_p,\mu)\simeq\Q/\Z$ is Tate's local duality pairing, from which the dual terms
$H^1_f(\Q_p,T^D)^*$ and $H^0(\Q,W^D)^*$ of Lemma lem:dfg21 come. In \Cref{ch:curves} the cup product on $H^1(\pi_1(X),\Z)$ for a compact Riemann surface $X$ agrees
with the topological one.
\end{remark}

\begin{proposition}[Pontryagin duality, recalled and stated]\label{prop:cohomology:pontryagin}
\begin{enumerate}[label=(\roman*)]
\item (Recalled from \Cref{ch:modules}.) $A^*=\Hom(A,\Q/\Z)$ for finite $A$, which is $\Hom(A,\Qlmod)$ when $A$ is $\ell$-primary; $T^*=\Hom_{\mathrm{cont}}(T,\Qlmod)$ for $T$ of
type (C) is discrete $\ell$-primary ($(\Zl^r)^*=(\Qlmod)^r$), and $W^*$ is of type (C) for $W$ discrete $\ell$-primary with $W[\ell]$ finite. The functor $X\mapsto X^*$ is exact,
$X^{**}=X$, and $\len X^*=\len X$ for finite $X$.
\item For a $G$-module $X$ of one of these kinds, $X^*$ is a $G$-module by $(g\varphi)(x)=\varphi(g^{-1}x)$, of the dual kind, and $(X^*)^G=(X_G)^*$.
\item (Stated.) For $G$ profinite and $A$ a discrete torsion $G$-module, $H^n(G,A)^*\simeq H_n(G,A^*)$, the homology of the profinite group $G$ with coefficients in the compact
module $A^*$; for $G$ finite and $A$ finite, $\hat H^n(G,A)^*\simeq\hat H^{-n-1}(G,A^*)$.
\end{enumerate}
\end{proposition}

\begin{proof}
(i) is \Cref{ch:modules}. (ii) $g\varphi$ is again a (continuous) homomorphism and the action is continuous because $G$ acts through finite quotients on each
finite piece. A homomorphism $\varphi$ is invariant if and only if $\varphi(g^{-1}x)=\varphi(x)$ for all $g$, $x$, that is, if and only if it kills $I_GX$. (iii) We do not prove
this; see \cite[Chapter I, Section 3 and Chapter II, Section 7]{NSW} and \cite[Chapter VI, Section 7]{Brown}.
\end{proof}

\begin{lemma}[lengths in exact sequences]\label{lem:cohomology:lengths}
\begin{enumerate}[label=(\roman*)]
\item \sloppy For a finite $\Zl$-module $A$, $\len(A^*)=\len A$; dualising an exact sequence of finite $\Zl$-modules gives an exact sequence with the arrows reversed.
\item In an exact sequence $0\to A\to B\to C\to D\to0$ of $\Zl$-modules of finite length, $\len B=\len A+\len C-\len D$. More generally the alternating sum of the lengths in
a finite exact sequence of modules of finite length is $0$.
\item Consequently, in the exact sequence of Lemma lem:dfg21,
\[
0\to H^1_f(\Q,W)\to H^1_\Sigma(\Q,W)\to\bigoplus_{p\in\Sigma\cup\{\infty\}}H^1_f(\Q_p,T^D)^*\to H^0(\Q,W^D)^*\to0,
\]
if all terms are finite,
\[
\len H^1_\Sigma(\Q,W)=\len H^1_f(\Q,W)+\sum_{p\in\Sigma\cup\{\infty\}}\len H^1_f(\Q_p,T^D)-\len H^0(\Q,W^D),
\]
and the term at $\infty$ is $0$ for $\ell$ odd (\Cref{pl:cohomology:real-place}).
\end{enumerate}
\end{lemma}

\begin{proof}
(i) is \Cref{ch:modules}. (ii) Length is additive on short exact sequences (Jordan-Holder, \Cref{ch:modules}). Split $0\to A\to B\to C\to D\to0$ at the image $X$ of $B\to C$:
$0\to A\to B\to X\to0$ and $0\to X\to C\to D\to0$ give $\len B=\len A+\len X$ and $\len C=\len X+\len D$; eliminating $\len X$ gives the formula. For a general exact
sequence $0\to A_1\to\dots\to A_r\to0$, let $X_i$ be the image of $A_i\to A_{i+1}$; then $0\to X_{i-1}\to A_i\to X_i\to0$ is exact (with $X_0=X_r=0$), so
$\len A_i=\len X_{i-1}+\len X_i$, and in $\sum_i(-1)^i\len A_i$ the terms $\len X_i$ cancel in pairs. (iii) Apply (ii), then (i) to the
last two terms, and use that the length of a finite direct sum is the sum of the lengths.
\end{proof}

\begin{paperlink}\label{pl:cohomology:lengths}
The proof of Theorem thm:main in the paper uses (iii) with $H^0(\Q,W^D)=0$ (Lemma lem:h0), and Lemma lem:fpr to evaluate each local term. Remark rem:ns uses it with
$\len H^0(\Q,W^D)=h_\ell(E)$, which is the displayed identity there: ``$\len H^1_\Sigma(\Q,W)=\len H^1_f(\Q,W)+\sum_{p\in\Sigma}\bigl(v_\ell(L_p(A_f,1)^{-1})+v_\ell(\Tam_p(A))\bigr)-h_\ell(E)$''. The paper's
``$^*$ is the Pontryagin dual'' in Lemma lem:dfg21 is the dual of finite groups, and the only properties used are those of (i). The finiteness of the terms is
part of the input of \Cref{ch:blochkato} (DFG04's theorems), not a consequence of this lemma.
\end{paperlink}

\section{Galois cohomology: notation and the road ahead}\label{sec:cohomology:galois}

This short section fixes the notation of Galois cohomology and proves the one formula that is purely group-theoretic: the unramified part of $H^1$ of a
local field is $M^{I}/(\Frob-1)M^I$. Everything arithmetic (Hilbert 90, Kummer theory, local duality, Selmer groups) is in \Cref{ch:galcoh}.

\begin{convention}[Galois cohomology]\label{conv:cohomology:galois}
For a field $K$ with separable closure $\overline K$, $G_K=\Gal(\overline K/K)$ is profinite (\Cref{ch:galois}); a \emph{Galois module} is a $G_K$-module of type (D), (C) or
(V), and $H^n(K,M):=H^n(G_K,M)$, continuous cochain cohomology. For $K=\Q_p$: $I_p\trianglelefteq G_{\Q_p}$, and $G_{\Q_p}/I_p=\Gal(\overline\F_p/\F_p)\simeq\widehat\Z$ is generated by the
arithmetic Frobenius $\Frob_p$ ($x\mapsto x^p$); $H^n(I_p,M)$ carries the conjugation action of $G_{\Q_p}/I_p$; $H^n(\R,M):=H^n(\Gal(\C/\R),M)$; $\Zl(1)=\varprojlim\mu_{\ell^k}$
with $\chi_\ell(\Frob_p)=p$ for $p\ne\ell$; $\Frob_p$ never denotes the geometric Frobenius $\Frob_p^{-1}$. The \emph{unramified part} of $H^1$ is
\[
\Hur(\Q_p,M):=H^1(G_{\Q_p}/I_p,M^{I_p})\ \xrightarrow[\ \sim\ ]{\ \Inf\ }\ \ker\bigl(\Res\colon H^1(\Q_p,M)\to H^1(I_p,M)\bigr),
\]
the two descriptions agreeing by \Cref{thm:cohomology:inf-res}(a) and \Cref{rem:cohomology:inf-res-continuous}.
\end{convention}

\begin{proposition}[unramified cohomology, a preview]\label{prop:cohomology:unramified}
Let $p$ be a prime and $\Gamma=G_{\Q_p}/I_p\simeq\widehat\Z$.
\begin{enumerate}[label=(\roman*)]
\item For $M$ of type (C): $\Hur(\Q_p,M)=M^{I_p}/(\Frob_p-1)M^{I_p}$.
\item For $V$ of type (V): $\Hur(\Q_p,V)=V^{I_p}/(\Frob_p-1)V^{I_p}$, and it is $0$ if and only if $\det(1-\Frob_p\mid V^{I_p})\ne0$, if and only if
$\det(1-\Frob_p^{-1}\mid V^{I_p})\ne0$.
\end{enumerate}
\end{proposition}

\begin{proof}
$M^{I_p}$ is a closed $\Zl$-submodule of $M$, hence of type (C) (resp. a subspace, of type (V)), with a continuous action of $\Gamma$, and $\widehat\Z$ surjects onto $\Zl$.
Apply \Cref{thm:cohomology:procyclic-compact} with $F=\Frob_p$. The last equivalence holds because $1$ is an eigenvalue of an invertible operator if and only if it is
an eigenvalue of its inverse.
\end{proof}

\begin{paperlink}\label{pl:cohomology:unramified}
After Lemma lem:dfg21: ``For $p\in\Sigma$ the group $H^1_f(\Q_p,V^D)=H^1_{\mathrm{ur}}(\Q_p,V^D)=(V^D)^{I_p}/(\Frob_p-1)$ vanishes, because
$\det(1-\Frob_p^{-1}\mid(V^D)^{I_p})=L_p(A_f,1)^{-1}\ne0$''. The last equality in the chain and the vanishing are \Cref{prop:cohomology:unramified}(ii). (The image of $\Frob_p-1$ equals that of
$\Frob_p^{-1}-1=-\Frob_p^{-1}(\Frob_p-1)$, so the quotient does not depend on which Frobenius is meant.) The equality
$H^1_f=\Hur$ at $p\ne\ell$ is \Cref{ch:blochkato}, and the value of the determinant is \Cref{ch:lfunctions}. The next sentences, ``So $H^1_f(\Q_p,T^D)$ is the torsion of
$H^1(\Q_p,T^D)$, a finite group'', use \Cref{prop:cohomology:les-cont}(iii): the preimage of $0$ under $H^1(\Q_p,T^D)\to H^1(\Q_p,V^D)$ is the torsion; finiteness is
\Cref{prop:cohomology:h1-fg}(ii) with the finiteness theorem of \Cref{ch:galcoh}.
\end{paperlink}

\begin{remark}[the road ahead]\label{rem:cohomology:road}
\Cref{ch:galcoh} proves: Hilbert's theorem 90 and Kummer theory, $H^1(K,\mu_n)=K^\times/(K^\times)^n$ and $H^1(K,\Zl(1))=\varprojlim K^\times/(K^\times)^{\ell^k}$ (with
\Cref{thm:cohomology:tate-limits}); the Kummer class of the Tate parameter and the value of $c'$ in \Cref{prop:cohomology:model}; finiteness and Tate local duality for
$G_{\Q_p}$; the unramified classes as exact annihilators of each other; Selmer groups and the comparison sequence behind Lemma lem:dfg21. \Cref{ch:blochkato}
defines the Bloch-Kato conditions $H^1_f$ and the Tamagawa factors $\Tam_p$, proves Lemma lem:fpr, and assembles Lemma lem:tamfactor from \Cref{prop:cohomology:model}.
\end{remark}

\section{Exercises}\label{sec:cohomology:exercises}

\begin{exercise}\label{ex:cohomology:dd3}
Verify $d^3d^2=0$ on $C^2(G,M)$ by expanding $(d^3d^2\psi)(g_1,g_2,g_3,g_4)$; there are $16$ terms, which cancel in pairs.
\emph{Hint.} Group the terms by the pattern of products of the $g_i$ in the arguments of $\psi$.
\end{exercise}

\begin{exercise}\label{ex:cohomology:cyclic-ex}
(a) For $G$ cyclic of order $m$ acting trivially on $\Z/n$, show $H^1(G,\Z/n)\simeq H^2(G,\Z/n)\simeq\Z/\gcd(m,n)$. (b) Let $t$ generate $G=\Z/4$ and act on $\Z[i]$ by multiplication by $i$.
Compute $H^0$, $H^1$, $H^2$.
\emph{Answer.} (a) $N=m$ and $t-1=0$, so $H^1=\ker(m\mid\Z/n)=(n/\gcd(m,n))\Z/n\Z$ and $H^2=(\Z/n)/m(\Z/n)=\Z/\gcd(m,n)$, both cyclic of order $\gcd(m,n)$. (b) $H^0=0$; $N=1+i+i^2+i^3=0$, so $H^1=\Z[i]/(i-1)\simeq\Z/2$ and $H^2=0/0=0$.
\end{exercise}

\begin{exercise}\label{ex:cohomology:real-place-ex}
Let $\sigma$ be complex conjugation, acting on $\Zl(1)=\varprojlim\mu_{\ell^k}$ by $\zeta\mapsto\zeta^{-1}$. Show $H^0(\R,\Zl(1))=0$ and $H^1(\R,\Zl(1))\simeq\Zl/2\Zl$, which is $\Z/2$
for $\ell=2$ and $0$ for $\ell$ odd; and $H^1(\R,\Zl)=0$ for every $\ell$ (trivial action). Deduce that for $\ell$ odd $H^1(\R,M)=0$ for every subquotient $M$ of
$\ad T(1)$, in accordance with the paragraph after Lemma lem:dfg21.
\emph{Answer.} \Cref{thm:cohomology:cyclic}(i): on $\Zl(1)$, $N=1+\sigma=0$ and $(\sigma-1)\Zl(1)=2\Zl(1)$; on $\Zl$, $N=2$ is injective. The last part is
\Cref{cor:cohomology:z2}.
\end{exercise}

\begin{exercise}\label{ex:cohomology:z4}
(a) Show that $H^2(\Z/2,\Z/2)\simeq\Z/2$ (trivial action). (b) For the extension $0\to2\Z/4\Z\to\Z/4\Z\to\Z/2\Z\to0$ and the section $s(0)=0$, $s(1)=1$, compute the cocycle
$\psi$ of \Cref{thm:cohomology:h2-ext} and show it is not a coboundary. (c) Which extension corresponds to $0$?
\emph{Answer.} (a) \Cref{thm:cohomology:cyclic}(ii): $M^G/NM=(\Z/2)/2(\Z/2)$. (b) $\psi(1,1)=s(1)+s(1)-s(0)=2$, the nonzero element, and $\psi$ vanishes elsewhere. For a coboundary,
$(d^1\varphi)(0,0)=\varphi(0)$ and $(d^1\varphi)(1,1)=2\varphi(1)-\varphi(0)=\varphi(0)$ in $\Z/2$, so a coboundary takes equal values at $(0,0)$ and $(1,1)$; as
$\psi(0,0)=0\ne\psi(1,1)$, $\psi$ is not one. (c) $\Z/2\times\Z/2$.
\end{exercise}

\begin{exercise}\label{ex:cohomology:tvw}
Let $\widehat\Z$ act trivially on $T=\Zl$, $V=\Ql$, $W=\Qlmod$. Compute $H^0$ and $H^1$ of all three, write down the long exact sequence of
\Cref{prop:cohomology:les-cont}(c), and check that the connecting map $H^0(\widehat\Z,W)\to H^1(\widehat\Z,T)$ is zero, in accordance with $H^1(\widehat\Z,T)_{\mathrm{tor}}=0$.
Then do the same with $F$ acting on $T=\Zl(1)$ by $q$, $q\ne\ell$, and check that the connecting map is onto $H^1(\widehat\Z,\Zl(1))=\Z/\ell^{v_\ell(q-1)}$.
\emph{Answer.} Trivial action: $\Zl,\Ql,\Qlmod$ in degree $0$ and $\Zl,\Ql,\Qlmod$ in degree $1$ (\Cref{thm:cohomology:procyclic-compact,thm:cohomology:procyclic-discrete}), and $H^0(V)\to H^0(W)$ is onto.
Second case: $H^0(V)=0$, $H^0(W)=\mu_{\ell^\infty}^{\Frob}\simeq\Z/\ell^{v_\ell(q-1)}$, $H^1(V)=0$. So $\delta$ is injective (its kernel is the image of $H^0(V)=0$) between two
groups of the same finite order $\ell^{v_\ell(q-1)}$, hence bijective.
\end{exercise}

\begin{exercise}\label{ex:cohomology:twists}
Let $\widehat\Z$, with generator $F$, act on $\Zl(j)=\Zl(1)^{\otimes j}$ ($j\in\Z$) by $q^j$, where $q\ne\ell$ is a prime. Show that $H^1(\widehat\Z,\Zl(j))\simeq\Z/\ell^{v_\ell(q^j-1)}$ for $j\ne0$, and
compute $v_\ell(q^j-1)$ for $\ell=3$, $q=7$, $j=1,2,3$.
\emph{Answer.} \Cref{thm:cohomology:procyclic-compact}; $v_3(6)=1$, $v_3(48)=1$, $v_3(342)=2$.
\end{exercise}

\begin{exercise}\label{ex:cohomology:identities}
Verify the six identities of \Cref{lem:cohomology:tate-shape} and $(t-1)^3=0$ on $\ad T$ by machine. In PARI/GP (this code was run):
\begin{lstlisting}[language=GP]
e=[0,1;0,0]; h=[1,0;0,-1]; f=[0,0;1,0];
Ad(A,x)=A*x*A^(-1);
R=[1,n;0,1];
[Ad(R,e)-e, Ad(R,h)-h+2*n*e, Ad(R,f)-f-(n*h-n^2*e)]
D(x)=Ad(R,x)-x;
D(D(D(f)))
for(s=-1,1, if(s, F=[s*q,c;0,s]; print([Ad(F,e)-q*e, Ad(F,h)-(h-2*s*c*e), Ad(F,f)-(f+s*c*h-c^2*e)/q])))
\end{lstlisting}
Every output is a zero matrix. Then check the paper's claim that $H^1(I_q,\ad T)_{\mathrm{tor}}\simeq(\Z/\ell^{v_\ell(n)})^2$ for $\ell$ odd, by computing the elementary divisors of the
$3\times3$ matrix of $t-1$ in the basis $e,h,f$.
\emph{Answer.} The columns are $0$, $(-2n',0,0)$, $(-n'^2,n',0)$; for $\ell$ odd the elementary divisors are $n',n',0$ up to units.
\end{exercise}

\begin{exercise}\label{ex:cohomology:ell2}
Redo \Cref{prop:cohomology:model}(a) for $\ell=2$. Show that $(t-1)M=n'(2\Z_2e\oplus\Z_2h')$ with $h'=h-n'e$, so that $H^1(I,\ad T)_{\mathrm{tor}}\simeq\Z/2^{w+1}\oplus\Z/2^w$, not $(\Z/2^w)^2$; and that
for $M=T$ and $\varepsilon=-1$ the invariants are $\{x:2x\equiv0\}\simeq\Z/2$ when $w\ge1$. Where else does the proof of \Cref{prop:cohomology:model} use $\ell$ odd?
\emph{Answer.} $e,h',f$ is a basis, and $(t-1)h=-2n'e$, $(t-1)f=n'h'$. In (c) the factor $2$ in $2\varepsilon c'q^{-1}y$ is no longer a unit.
\end{exercise}

\begin{exercise}\label{ex:cohomology:end}
Let $\ell$ be odd and $\End T=\ad T\oplus\Zl\cdot1$ in the model setting. Show that $H^1(I,\End T)_{\mathrm{tor}}=H^1(I,\ad T)_{\mathrm{tor}}$, compatibly with the action of $G/I$, so
that the Fitting ideal of the invariants is the same for $\End T$ and for $\ad T$.
\emph{Hint.} The decomposition is $G$-stable for $\ell$ odd ($x=(x-\frac12\tr x)+\frac12\tr x$), and $H^1(I,\Zl)=\Zl$ is torsion-free (\Cref{ex:cohomology:procyclic-examples}(i)).
\end{exercise}

\begin{exercise}[the curve \curve{11a1} at $q=11$, $\ell=5$]\label{ex:cohomology:11a1}
The curve $E=\curve{11a1}$ has split multiplicative reduction at $11$ with $n=v_{11}(\Delta_E)=5$ ($\Delta_E=-11^5$), Tamagawa number $c_{11}=5$, $j_E=-2^{12}\cdot31^3/11^5$, and its Tate parameter
is $q_E=10\cdot11^5+O(11^6)$, so the unit part $u=q_E\,11^{-5}$ satisfies $u\equiv-1\pmod{11}$ (computed with PARI/GP: \texttt{ellinit(E,O(11\^{}8)).tate}). (a) Check
$u\equiv-1$ by hand from $q_E=j_E^{-1}(1+O(11^5))$, using $2^{12}\equiv4$ and $31^3\equiv3\pmod{11}$. (b) The fifth powers in $\F_{11}^\times$ are $\{\pm1\}$, so $\bar u$ is a fifth power; granting
from \Cref{ch:galcoh} that $\min(v_5(n),v_5(q-1),v_5(c'))=\min(v_5(n),v_5(q-1),v_5(c))$ with $v_5(c)=v_5(q-1)=1$ here, deduce from \Cref{prop:cohomology:model} that
$(H^1(I_{11},\ad T_5E)_{\mathrm{tor}})^{G/I}$ has order $25$. (c) Check \Cref{prop:cohomology:model}(e) against $c_{11}=5$. (d) The modular degree of \curve{11a1} is $1$. Explain why this does not
contradict Theorem thm:main.
\emph{Answer.} (b) $w=1$, $v_5(10)=1$, $m=1$, order $5^{1+1}$. (c) $\varepsilon=1$, invariants $\Z/5$. (d) Theorem thm:main requires $E[5]$ irreducible, and $E(\Q)\simeq\Z/5$, so $E[5]$
is reducible; the computation illustrates Lemma lem:tamfactor, which needs no irreducibility.
\begin{lstlisting}[language=GP]
E = ellinit([0,-1,1,-10,-20], O(11^8));
Q = E.tate[3];  \\ Tate parameter q_E
[valuation(Q,11), (Q/11^5) % 11]   \\ returns [5, 10]
\end{lstlisting}
\end{exercise}

\begin{exercise}\label{ex:cohomology:sylow}
Let $G$ be finite, $P$ a Sylow $\ell$-subgroup and $M$ a $G$-module. Show that $\Res\colon H^1(G,M)\to H^1(P,M)$ is injective on the $\ell$-primary part of $H^1(G,M)$. This is the
step ``by restriction and corestriction'' in the proof of Lemma lem:tamadd.
\emph{Hint.} $\Cor\circ\Res=[G:P]$ is prime to $\ell$ (\Cref{prop:cohomology:res-inf-cor}(d), proved in degree $1$).
\end{exercise}

\begin{exercise}\label{ex:cohomology:c3}
Let $C_3=\langle\sigma\rangle$. (a) Show $H^1(C_3,\Z_3[C_3])=0$, using \Cref{thm:cohomology:cyclic}. (b) Let $T=\Z_3^2$ with $\sigma$ acting with characteristic polynomial $X^2+X+1$. Show
$H^1(C_3,T)=T/(\sigma-1)T\simeq\F_3$. These are the two facts behind Lemma lem:tamadd at $\ell=3$ and the sentence after it, ``then
$H^1(\langle\sigma\rangle,T)=T/(1-\zeta_3)T\simeq\F_3$''.
\emph{Answer.} (a) On $\Z_3[C_3]$, $\ker N$ is the augmentation ideal, which equals $(\sigma-1)\Z_3[C_3]$. (b) $N=1+\sigma+\sigma^2=0$ on $T$, and $T\simeq\Z_3[\zeta_3]$ with $\sigma=\zeta_3$
(\Cref{ch:ant}), so $T/(\sigma-1)T=\Z_3[\zeta_3]/(\zeta_3-1)\simeq\F_3$.
\end{exercise}

\begin{exercise}\label{ex:cohomology:h0-length}
Let $\chi\colon G\to\Zl^\times$ be a character with $\chi(G)\subseteq1+\ell^h\Zl$ but $\chi(G)\not\subseteq1+\ell^{h+1}\Zl$, $T=\Zl(\chi)$ and $W=T\otimes\Qlmod$. Show that $H^0(G,T/\ell^kT)$ is all
of $T/\ell^kT$ for $k\le h$ and has no element of order $\ell^k$ for $k>h$, so that $H^0(G,W)\simeq\Z/\ell^h$ by \Cref{lem:cohomology:h0-torsion}(d).
\emph{Answer.} For $k\le h$, $g$ acts on $T/\ell^kT=\Z/\ell^k$ by $\chi(g)\equiv1\pmod{\ell^k}$, trivially. For $k>h$ pick $g$ with $v_\ell(\chi(g)-1)=h$; for $x\notin\ell T$,
$(\chi(g)-1)x$ has valuation $h<k$, so $gx\ne x$ in $T/\ell^kT$. By \Cref{lem:cohomology:h0-torsion}(b), $H^0(G,W)[\ell^{h+1}]\simeq H^0(G,T/\ell^{h+1}T)$ has no element of
order $\ell^{h+1}$, so $H^0(G,W)$ is killed by $\ell^h$ and lies in the finite group $W[\ell^h]$; and $\#H^0(G,T/\ell T)\le\#(T/\ell T)=\ell$. So (d) applies, with largest $k$
equal to $h$.
\end{exercise}

\begin{exercise}\label{ex:cohomology:zp-not-zl}
Let $p\ne\ell$ and let $\Gamma=\Zp$ act on $M$ of type (C) through an element $u$ of order $p^a$. Show that $H^1(\Gamma,M)=0$, that $\ker N/(u-1)M=0$ for $N=\sum_{i<p^a}u^i$, and that
$M/(u-1)M\ne0$ in general (take $u=1$). So the formula of \Cref{thm:cohomology:procyclic-compact} needs the surjection onto $\Zl$.
\emph{Hint.} \Cref{thm:cohomology:prime-to-l}(b); and \Cref{thm:cohomology:cyclic} for the finite quotient, whose order is a unit of $\Zl$.
\end{exercise}

\begin{exercise}[harder]\label{ex:cohomology:transgression}
Prove \Cref{thm:cohomology:inf-res}(b). Let $c\in Z^1(H,M)$ have $G/H$-invariant class. (a) Show that $c$ can be changed by a coboundary and extended to a cochain $\tilde c$
on $G$ with $\tilde c(gh)=\tilde c(g)+gc(h)$ and $\tilde c(hg)=c(h)+h\tilde c(g)$ for $g\in G$, $h\in H$. (b) Show that $d^1\tilde c$ takes values in $M^H$ and factors through
$G/H\times G/H$, defining a $2$-cocycle whose class $\operatorname{tg}[c]\in H^2(G/H,M^H)$ is well defined. (c) Prove exactness at $H^1(H,M)^{G/H}$ and at $H^2(G/H,M^H)$.
\emph{Hint.} \cite[Chapter I, Section 6]{NSW}; for (a) fix coset representatives and use the invariance $\sigma\cdot c-c=d^0m_\sigma$.
\end{exercise}

\section{Notes and references}\label{sec:cohomology:notes}

The homological algebra of \Cref{sec:cohomology:complexes,sec:cohomology:derived} follows Weibel \cite{Weibel}: the long exact sequence, the snake lemma, $\Ext^1$ and
extensions, and the Lyndon-Hochschild-Serre spectral sequence are in his Sections 1.3, 3.4 and 6.8; we recall the theorem numbers 1.3.1, 1.3.2, 3.4.3, 3.5.8 and 6.8.2
from memory and have not checked them against a printed copy. Group cohomology by cochains, restriction, corestriction, inflation and the conjugation action
follow Brown \cite[Chapters III to VI]{Brown}, Serre's \emph{Local Fields} \cite[Chapters VII and VIII]{SerreLocalFields}, and Atiyah and Wall's chapter in
Cassels-Frohlich \cite[Chapter IV]{CasselsFrohlich}; the chapter and section pointers to these books are likewise unverified against the printed editions. The
explicit cochain $\psi$ in the proof of \Cref{thm:cohomology:order-kills} avoids the corestriction in higher degrees.

Continuous cohomology of profinite groups with discrete coefficients is in Serre's \emph{Galois Cohomology} \cite[Chapter I]{SerreGaloisCoh} and in
Neukirch, Schmidt and Wingberg \cite[Chapters I and II]{NSW}. For compact coefficients the standard reference is Tate \cite[Section 2]{Tate95}, where the
$\varprojlim^1$ sequence of \Cref{thm:cohomology:tate-limits} is proved; Rubin \cite[Appendix B]{Rubin} gives a summary adapted to Galois representations. Our proof
of \Cref{thm:cohomology:procyclic-compact} is direct: it constructs the cocycle from the geometric sums of \Cref{lem:cohomology:geometric-sums} instead of passing to
the limit over $M/\ell^kM$, so that the identity the paper uses does not depend on the stated part of Tate's theorem. The procyclic computation with discrete
coefficients is Serre's \cite[Chapter XIII, Section 1]{SerreLocalFields} (pointer unverified). The supernatural order and the maximal pro-$\ell$ quotient are in
\cite[Chapter I, Section 1]{SerreGaloisCoh}; Milne \cite[Chapter II]{MilneCFT} is a readable alternative for inflation-restriction and Tate cohomology.

The spectral sequence of \Cref{thm:cohomology:hochschild-serre} was found by Lyndon \cite{Lyndon48} and in general form by Hochschild and Serre \cite{HochschildSerre53};
only its five-term exact sequence is used in the book, and only through \Cref{cor:cohomology:pro-l-quotient}(b) and \Cref{prop:cohomology:tame-inertia}(iii), neither of
which the paper needs.

The definition of the Tamagawa factor as the Fitting ideal of the $G_{\Q_p}$-invariants of the torsion of $H^1(I_p,M)$, which \Cref{sec:cohomology:model} computes, is
Fontaine and Perrin-Riou's \cite[I.4.2.2]{FPR94}, as recalled by Diamond, Flach and Guo \cite[p.~708]{DFG04}. The shape of the Tate representation is
\cite[Chapter V]{SilvermanATAEC} and \cite[Theorem 3.1(e)]{DDT}; the corner $\varepsilon\kappa$ is the revised paper's. The conjugation action is the paper's
$(\sigma\cdot c)(i)=\sigma c(\sigma^{-1}i\sigma)$; with the other common convention the Frobenius acts on $H^1(I,M)$ by $q\rho_M(\Frob)^{-1}$ (\Cref{rem:cohomology:other-convention}).
